\documentclass[11pt,a4paper]{article}
\usepackage[margin=21mm]{geometry}
\usepackage{amsmath,amssymb,amsthm,mathtools,mathrsfs}
\usepackage{fontspec,unicode-math}
\setmainfont{Latin Modern Roman}
\setmathfont{Latin Modern Math}
\usepackage{longtable,booktabs,array,enumitem,microtype,xurl}
\usepackage[unicode,hidelinks]{hyperref}
\usepackage{bookmark}
\newtheorem{theorem}{Theorem}[section]
\newtheorem{proposition}[theorem]{Proposition}
\newtheorem{lemma}[theorem]{Lemma}
\newtheorem{corollary}[theorem]{Corollary}
\theoremstyle{definition}
\newtheorem{definition}[theorem]{Definition}
\newtheorem{remark}[theorem]{Remark}
\providecommand{\Tr}{\operatorname{Tr}}
\providecommand{\Spec}{\operatorname{Spec}}
\providecommand{\coker}{\operatorname{coker}}
\providecommand{\diag}{\operatorname{diag}}
\providecommand{\rem}{\operatorname{rem}}
\providecommand{\im}{\operatorname{im}}
\providecommand{\C}{\mathbb C}
\providecommand{\R}{\mathbb R}
\providecommand{\Q}{\mathbb Q}
\allowdisplaybreaks[2]
\setlength{\emergencystretch}{4em}
\setlength{\parskip}{4pt}
\setlength{\parindent}{0pt}
\hypersetup{pdftitle={Reviewed mixed-boundary continuation},pdfauthor={}}
\title{Reviewed mixed-boundary continuation\\\large Original coefficients, derived fibres, inertia and arithmetic source maps}
\author{}
\date{14 September 2026}
\begin{document}
\maketitle
\begin{abstract}
This manuscript reconstructs the missing proof bodies in the supplied
mixed-boundary intake and corrects their coefficient, specialization and
source domains. It joins the derived primary fibre and finite-field
boundary to the original arithmetic source through explicit maps. The
new whole-amplitude period comparison retains its derivative and boundary
terms. The full current programme and its inherited proofs accompany
this continuation. No Riemann-hypothesis endpoint is asserted.
\end{abstract}
\tableofcontents
\clearpage
\section{How this intake changes the current calculation}
The received archive has an intact checksum but an incomplete proof body:
MB20--44 and MB46--47 are absent, and its recorded PDF validation failed.
The following complete reconstructions replace the missing and mistyped
claims. The unmodified archive is retained as provenance, not as a
second competing mathematical authority.

The original total object and all its earlier arrows remain. IAL supplies
the exact formal and ramified maps with their corrected coefficient rings;
FF computes the actual finite-field descent and parameter-restriction
map; IPM and ECB construct the whole-amplitude period observation, its
relation image and explicit error estimates in the original source metric.
IMS attaches them to both original legs and each independent coefficient
face. The complete imported HCD, PGB and PGMB proofs then give the original
holonomy comparison, the completed graph quotient and its finite tails.
The inherited proof dependencies are supplied in the accompanying complete
cumulative source, including the original analytic input, SPG, PGS, PGM,
CMI, AEX and the exact AAM exclusion scale.

In particular the primary relative socle is over the full parameter ring.
It is not the ordinary module-theoretic socle. The formal normal-coordinate
raising matrix is not identified with the arithmetic nilpotent for a
general full packet; IAL32 computes their entire marked relationship.
The finite period observation remains a valid coefficient observation.
IPM9--15d computes its separate relationship to the actual entire amplitude,
including the full unit and every derivative remainder. Its relation image
is retained before any receiving quotient.

The prior SPG Stokes construction applies on its proved marked slice.
The intake's earlier statement that no such Stokes calculation was
available does not supersede that existing construction. Neither its
coordinate isometry nor the common ramification determinant cancellation
is used as an arithmetic exponent bound.


% BEGIN COMPLETE CORRECTED SOURCE retained_base_source.tex
\section{The retained absolute base and theta source}
For a commutative ring $R$ keep the original semiring
\begin{equation}
 G(R)=\{\tau\}\sqcup R^\bullet,\qquad e_R=0_R^\bullet,\qquad
 r^\bullet+s^\bullet=(r+s)^\bullet,\quad
 r^\bullet s^\bullet=(rs)^\bullet.                    \tag{MB1}
\end{equation}
External $\tau$ is the additive identity and is absorbing for multiplication.
The two observations are
\[
 p_R(\tau)=0,\quad p_R(r^\bullet)=r,\qquad
 b_R(\tau)=0,\quad b_R(r^\bullet)=1\in\mathbb B.
\]
The map $(p_R,b_R)$ is injective: its images are $(0,0)$ and $(r,1)$.
In particular it retains both $(0,0)$ and $(0,1)$.
The pointed/preaddition base is
\begin{equation}
 \mathbf F_{1,\tau}=\{\tau_0,1\},\quad
 \mathfrak b_\tau=\operatorname{Spec}_{\rm mon}\mathbf F_{1,\tau},
 \quad j_R(\tau_0)=\tau,\quad j_R(1)=1_R^\bullet .       \tag{MB2}
\end{equation}
Its preaddition deletes $\tau_0$ terms and preserves the number of unit
terms.  Thus $j_R$ preserves each base relation, and $1+1=1$ is not imposed.
The coefficient realization is the original square
\begin{equation}
\begin{array}{ccc}
G(\mathbb Z)&\xrightarrow{G(\iota)}&G(\mathbb C)\\
p_{\mathbb Z}\downarrow&&\downarrow p_{\mathbb C}\\
\mathbb Z&\xrightarrow{\iota}&\mathbb C.
\end{array}                                                    \tag{MB3}
\end{equation}
The target of $p_{\mathbb Z}$ is infinite.  The coefficient stalk of the
original direct image to the one-point base retains $G(R)$; its support
observation is the further map $b_R:G(R)\to\mathbb B$.
These are the inherited structural maps, not a new general six-functor
formalism for semirings.

For a coefficient vector space $V$, the lift of a linear map $f:V\to W$ is
$f^\tau(v^\bullet)=f(v)^\bullet$, $f^\tau(\tau)=\tau$.
For the original support diagram it is applied in each fibre with its
original transport. At a present nonbottom outer label, a vector mapped
to zero reaches that fibre's supported zero. If an outer support map
itself reaches bottom, its original bottom fibre must also be retained;
no bottom-reflection property is assumed.

\subsection{The original theta source and finite observation}
Retain, with their original topologies,
\begin{align}
 V&=\{\phi\in\mathcal S(\mathbb R)_{\rm ev}:\phi(0)=0,
                         \ \int_{\mathbb R}\phi=0\},\nonumber\\
 \mathscr B&=\{F\in C^\infty(\mathbb R_{>0}):
 \sup_{x>0}x^b|D^jF(x)|<\infty\quad(b\in\mathbb Z,j\geq0)\},
\nonumber\\
 D&=-x\partial_x,\nonumber\\
 \widehat\phi(\nu)&=\int_{\mathbb R}\phi(x)e^{-2\pi ix\nu}\,dx,
 \qquad\Theta\phi(x)=2\sum_{n\geq1}\phi(nx),
 \qquad Q=\mathscr B/\Theta V.                         \tag{MB4}
\end{align}
The Mellin map is $\mathcal MF(s)=\int_0^\infty F(x)x^s dx/x$.
The unchanged source identities are
\begin{equation}
 \phi_*=(4\pi^2x^4-6\pi x^2)e^{-\pi x^2},\qquad
 g=2\xi,\qquad \mathcal M\Theta(P(D)\phi_*)=gP.          \tag{MB5}
\end{equation}
Their analytic range and Euler inverse theorems are inherited written
results in AG/CAU and \path{arithmetic_input.tex}. They are not newly
certified by the finite checks of this note.

Let $h=\prod_\rho(s-\rho)^{m_\rho}$ be a nonempty packet of actual
nontrivial zeros at their complete orders.  Set
\begin{equation}
 E_h=\mathbb C[s]/h,\quad A_h=M_s,\quad v_h=g/h,\quad
 \upsilon_h=j_hv_h\in E_h^\times,\quad
 \mathcal MF_h=v_h,\quad \mathcal T_h(P)=P(D)F_h.         \tag{MB6}
\end{equation}
The source inverse defines $F_h$ at precisely this domain.  With
$\alpha(P)=P(D)\phi_*$ the actual chain map is
\begin{equation}
\begin{array}{ccc}
\mathbb C[s]&\xrightarrow{\times h}&\mathbb C[s]ds\\
\alpha\downarrow&&\downarrow\mathcal T_h\\
V&\xrightarrow{\Theta}&\mathscr B,
\end{array}
\qquad
J_h\mathcal T_h(P)=\upsilon_h[P]_h.                    \tag{MB7}
\end{equation}
Its top map is $E_h\xrightarrow{M_{\upsilon_h}}E_h
\xrightarrow{\sigma_h}Q$.  Thus observed arithmetic coordinates enter
polynomial coordinates through $M_{\upsilon_h^{-1}}$, not by replacing
$\upsilon_h$ by one. The exact original observation sequence is
\begin{equation}
0\to Q[h(D)]\to E_h\xrightarrow{\times j_hg}E_h
                  \to Q/h(D)Q\to0.                    \tag{MB8}
\end{equation}
A finite observation's full kernel on $\mathscr B$ is not silently
identified with $\Theta V$: (MB8) and the specified restriction of (MB7)
are the maps used here.

On the critical strip use the enlargement from the MCF off-critical
object to $\Omega=\{0<\Re s<1\}$. Its included critical and off-critical
primary modules stay separately present. No nonempty off-critical packet
is assumed in any theorem below.


% END COMPLETE CORRECTED SOURCE retained_base_source.tex


% BEGIN COMPLETE CORRECTED SOURCE algebra.tex
\section{Reconstructed formal, ramified, and mixed coefficient calculation}
\label{sec:intake-algebra}

This proof body reconstructs the algebra missing from the received
\texttt{NEW\_BODY.tex}: the received proof stops after MB19 and resumes
inside an unrelated display MB45 and the proof of MB47.  The same deletion
occurs in \texttt{CONTINUATION.tex} and its integration patch.  Consequently
the following proofs are new reconstructions, not quotations of an intact
received proof.  They use the original marked primitive, polynomial
coordinates, one-form orientation, and coefficient rings.  In particular,
the parameter $a$ is never specialized without displaying that map.

\subsection{The original marked polynomial complex and its connection}
Let $h(s)=s^n+\sum_{i<n}h_i s^i$, $n\geq1$, be a nonempty monic actual
packet polynomial with its complete root multiplicities.  Set
$E_h=\mathbb C[s]/(h)$ and write $A_h=A_h(0)$ for multiplication by
$s$ on this marked algebra.  Put
\begin{equation}
 \Phi_h(s)=\frac{s^{n+1}}{n+1}+\sum_{i<n}\frac{h_i s^{i+1}}{i+1},
 \qquad f_t(s)=\Phi_h(s)-ts,\qquad
 L_h=u\partial_s+h(s)-t.                         \tag{IAL1}
\end{equation}
Over $R=\mathbb C[u,t]$ consider the complex in degrees zero and one
$C_h=[R[s]\xrightarrow{L_h}R[s]ds]$.  If a nonzero polynomial $P$ has
$s$-degree $r$, $L_hP$ has degree $r+n$ and the same leading coefficient.
Thus $L_h$ is injective.  Repeatedly subtracting $L_h(c s^{r-n})$ from a
polynomial of degree $r\geq n$ terminates and constructs unique $R$-linear
maps $\mathfrak H_h,r_h$ with
\begin{equation}
 P=L_h\mathfrak H_h(P)+r_h(P),\quad \deg_s r_h(P)<n,
 \quad \mathfrak H_h L_h=I,\quad r_h L_h=0.        \tag{IAL2}
\end{equation}
Uniqueness follows from the same leading-degree calculation.  Hence
$H^1(C_h)$ is the free $R$-module with the ordered representatives
$1,s,\ldots,s^{n-1}$ times $ds$.

After inverting $u$, on degree one the two connection operators are
$\partial_u-f_t/u^2$ and $\partial_t-s/u$.  On degree zero their
counterparts are $\partial_u-f_t/u^2+1/u$ and $\partial_t-s/u$.
Direct expansion gives
\begin{equation}
 (\partial_u-f_t/u^2)L_h=L_h(\partial_u-f_t/u^2+1/u),
 \qquad (\partial_t-s/u)L_h=L_h(\partial_t-s/u).
                                                        \tag{IAL3}
\end{equation}
For the first equality the commutator contributions are
$\partial_s+(h-t)/u=L_h/u$; for the second they are $-1+1=0$.
Thus these are actual operators on the complex and on its cohomology.
For the original ordinary monic division
\begin{equation}
 f_t s^b=(h-t)q_b+r_b,\quad
 C_h(t)e_b=\operatorname{coeff}(r_b),\quad
 B_h(t)e_b=\operatorname{coeff}(q_b'),                 \tag{IAL4}
\end{equation}
one has $[(h-t)q_b\,ds]=[-u q_b'\,ds]$.  The connection in the original
ordered top lattice is therefore
\begin{equation}
 \nabla_u^h=\partial_u-C_h(t)/u^2+B_h(t)/u,
 \qquad \nabla_t^h=\partial_t-A_h(t)/u,                \tag{IAL5}
\end{equation}
where $A_h(t)$ is ordinary multiplication by $s$ modulo $h-t$.
For $s$ times a representative of degree below $n$, its only possible
quotient in division by $h-t$ is constant, so its derivative term is zero;
this proves the stated formula for $A_h(t)$.
At $u\ne0$, this is a connection coefficient matrix, not a descended
multiplication action on the differential quotient: indeed
$sL_hQ-L_h(sQ)=-uQ$.  The original multiplication action is recovered
on the marked algebraic quotient $u=t=0$.
For an ordered $k$-fold external complex, the differential is the sum
of $L_{h,i}ds_i\wedge$, and the degree-$j$ $u$-operator is
$\partial_u-\sum_i f_{t_i}/u^2+(k-j)/u$.  Tensoring (IAL3), with the
original exterior order, proves the chain identity including that degree
term.  Tensoring the coefficient decompositions (IAL2) proves that only
top cohomology survives and that its basis consists of all ordered
$s_1^{b_1}\cdots s_k^{b_k}ds_1\wedge\cdots\wedge ds_k$, $b_i<n$.

\subsection{Formal primary coordinates, all coefficients, and their inverse}
Fix an actual root $\rho$ of multiplicity $m$, write $d=m+1$ and
$z=s-\rho$, and retain
\begin{equation}
 h(\rho+z)=z^m a_\rho(z),\quad a_\rho(0)\ne0,
 \quad \kappa_\rho=\Phi_h(\rho),\quad
 R_\rho(z)=d\sum_{j\geq0}\frac{a_\rho^{(j)}(0)}{j!}
                         \frac{z^j}{d+j}.             \tag{IAL6}
\end{equation}
The sum is finite.  Integrating $h$ from $\rho$ to $\rho+z$ proves
$\Phi_h(\rho+z)-\kappa_\rho=z^dR_\rho(z)/d$.
Choose and retain one $d$-th root $B_\rho(0)$ of $R_\rho(0)$.
On a sufficiently small disc, write $R_\rho(z)=R_\rho(0)(1+w(z))$
with $|w|<1$ and define
$B_\rho(z)=B_\rho(0)\sum_{r\geq0}{1/d\choose r}w(z)^r$.
This converges, has the prescribed root and constant term, and is unique
among germs with that constant term (divide two such germs and use that
a continuous $d$-th root of unity is locally constant).  The coordinate
$y=zB_\rho(z)$ has derivative $B_\rho(0)\ne0$, hence an analytic inverse
$z=\zeta_\rho(y)$.  With this precise branch,
\begin{equation}
 \Phi_h(\rho+\zeta_\rho(y))-t(\rho+\zeta_\rho(y))
 =\kappa_\rho+\frac{y^d}{d}-t(\rho+\zeta_\rho(y)),
 \qquad ds=\zeta_\rho'(y)dy.                         \tag{IAL7}
\end{equation}
For $r\geq1$, the residue substitution $y=zB_\rho(z)$, followed by
integration of a formal derivative, gives
\begin{equation}
 [y^r]\zeta_\rho(y)=\frac1r[z^{r-1}]B_\rho(z)^{-r}.
                                                        \tag{IAL8}
\end{equation}
More generally the particularly useful exact coefficient identity is
\begin{equation}
 [y^r]\bigl(Q(\rho+\zeta_\rho(y))\zeta_\rho'(y)\bigr)
   =[z^r]Q(\rho+z)B_\rho(z)^{-r-1}.                  \tag{IAL9}
\end{equation}
Indeed the left side is the residue of
$Q(\rho+\zeta(y))\zeta'(y)y^{-r-1}dy$; substituting $z=\zeta(y)$
turns it into $Q(\rho+z)z^{-r-1}B_\rho(z)^{-r-1}dz$.
This proof keeps every power of the chosen nonzero constant
$B_\rho(0)$, which need not equal one.

Use the iterated ring $\mathbb C[[z]][[u]]$ and then its coordinate
isomorphism to $\mathbb C[[y]][[u]]$.  At the marked value $t=0$, the
cochain maps are substitution on degree zero and substitution times
$\zeta_\rho'$ on degree one.  Since
$h(\rho+\zeta_\rho)\zeta_\rho'=y^m$, these maps turn $L_h$ exactly into
$u\partial_y+y^m$.  An arbitrary formal series in the latter target
has a unique reduction of degree below $m$ over $\mathbb C[[u]]$:
divide its $u^0$ coefficient by $y^m$, subtract the differential image,
and repeat at each successive $u$ coefficient.  The process converges
in the stated $u$-adic topology.  If an image had degree below $m$,
the first nonzero $u$ coefficient of its preimage would be multiplied
by $y^m$ at that order; this contradicts the degree restriction.  The
same argument proves zero kernel in degree zero.

For $0\leq b<m$ and $\ell\geq0$ the full reduction is
\begin{equation}
 [y^{b+\ell d}dy]=(-u)^\ell
       \prod_{r=0}^{\ell-1}(b+1+rd)[y^bdy],\qquad
 [y^{m+\ell d}dy]=0.                                \tag{IAL10}
\end{equation}
Here the empty product is one.  Applying the differential to $y^n$
gives $[y^{m+n}dy]=-un[y^{n-1}dy]$; iteration proves both formulas,
including their signs and integer coefficients.  For an input
$\sum_i u^i\sum_r c_{i,r}y^r$, any specified $u^N$ coefficient of the
reduction uses only the finitely many $i+\ell=N$.  Thus (IAL10) is
also a well-defined continuous formula on the entire iterated ring.

Put
$a_{\rho,j,r}=[y^r](\rho+\zeta_\rho(y))^j\zeta_\rho'(y)$.
In the original ordered global polynomial basis the comparison matrix
has entries
\begin{equation}
 K_{\rho,bj}(u)=\sum_{\ell\geq0}
 a_{\rho,j,b+\ell d}(-u)^\ell
          \prod_{r=0}^{\ell-1}(b+1+rd),\qquad 0\leq b<m.
                                                        \tag{IAL11}
\end{equation}
Use (IAL9) to compute each coefficient without a formal inversion
approximation.  At $u=0$ this is the original full primary-jet map,
followed by the invertible coordinate substitution and multiplication
by the invertible jet of $\zeta_\rho'$.  The Chinese remainder map
$\mathbb C[s]/h\to\bigoplus_\rho\mathbb C[z]/z^{m_\rho}$ is invertible
because its distinct factors are relatively prime: Bezout inverses
construct the inverse idempotents.  Consequently the square matrix
$K_h(0)$, formed by all these rows, is invertible.  Write
$K_h(u)=\sum_{r\geq0}K_ru^r$.  Its unique inverse is constructed by
\begin{equation}
 L_0=K_0^{-1},\qquad
 L_r=-K_0^{-1}\sum_{i=1}^r K_iL_{r-i}\quad(r\geq1).
                                                        \tag{IAL12}
\end{equation}
The coefficient identities prove $KL=I$; over the commutative formal
coefficient ring $\det K$ is a unit, so the adjugate identity also gives
$LK=I$.  This proves the full formal primary isomorphism, including
different roots with equal critical values, which remain separate rows.

Coordinate substitution commutes with the chain connections (IAL3).
In the local normal basis the connection on $[y^bdy]$ is obtained by
reducing $-(\kappa_\rho+y^d/d)y^b/u^2$.  Formula (IAL10) gives exactly
\begin{equation}
 \nabla_{u,\rho}=\partial_u-\frac{\kappa_\rho}{u^2}I
 +\frac1u\operatorname{diag}\left(\frac1d,\ldots,\frac md\right),
 \qquad K_h'+\Omega_{\rm loc}K_h=K_h\Omega_h.          \tag{IAL13}
\end{equation}
This proves the connection intertwining, rather than only equality of
the special fibres.  The coordinate construction here is at $t=0$;
the full $t$ dependence still has the displayed phase (IAL7).

For an analytic amplitude $A(y)$ bounded by $M_R$ on $|y|\leq R$, Cauchy's
coefficient formula gives $|[y^n]A|\leq M_R R^{-n}$.  Since
$b+1+rd\leq d(r+1)$, (IAL10) proves the explicit estimate
\begin{equation}
 |[u^\ell]\operatorname{red}_b(A)|
 \leq M_R R^{-b}(d/R^d)^\ell\ell!.                  \tag{IAL14}
\end{equation}
For the arithmetic theta amplitude one must use exactly
\begin{equation}
 A(y)=(g/h)(\rho+\zeta_\rho(y))
       P(\rho+\zeta_\rho(y))\zeta_\rho'(y),\qquad g=2\xi,
                                                        \tag{IAL15}
\end{equation}
and its own $M_R$.  The factorial estimate is a formal coefficient
bound, not a proof of convergence of $K_h(u)$ or of global sectorial
gluing of the theta integral.

\subsection{The exact ramified complete-primary family}
For the actual complete primary subpacket $h_\rho(s)=(s-\rho)^m$, its
own original primitive is
\begin{equation}
 \Phi_{h_\rho}(s)=\frac{(s-\rho)^d-(-\rho)^d}{d},
 \quad \kappa=-\frac{(-\rho)^d}{d},\quad
 u=v^d,\quad t=av^m,\quad s=\rho+vx,
 \quad F=\kappa v^{-d}-\rho a v^{-1}.                \tag{IAL16}
\end{equation}
This $\kappa$ is not silently identified with $\Phi_h(\rho)$ for a
larger packet.  Substitution proves the exact phase identity
$f_t/u=F+x^d/d-ax$ and the one-form identity $ds=v\,dx$.
Write $e_b=z^b$, $0\leq b<m$, for the original translated coefficient
basis.  Define $Ne_b=e_{b+1}$ for $b<m-1$, $Ne_{m-1}=0$, and
$Re_{m-1}=e_0$, $Re_b=0$ otherwise.  These definitions include $m=1$,
when $N=0$ and $R=I$.  Direct ordinary division gives
\begin{equation}
 \begin{aligned}
 A(t)&=\rho I+N+tR,\\
 B&=\operatorname{diag}((b+1)/d)_{b=0}^{m-1},\\
 C(t)&=(\kappa-\rho t)I-\frac{mt}{d}(N+tR).
 \end{aligned}                                      \tag{IAL17}
\end{equation}
Indeed $f_t=(\kappa-\rho t)+z^d/d-tz$.  Subtracting
$(z^m-t)z^{b+1}/d$ leaves
$(\kappa-\rho t)z^b-mt z^{b+1}/d$; in the last column the additional
constant quotient $-mt/d$ reduces $z^m$ to $t$ and has derivative zero.
The derivative of the first quotient is $(b+1)z^b/d$, proving every
column of (IAL17).

Put $D_v=\operatorname{diag}(v^{b+1})_{b=0}^{m-1}$ and $J(a)=N+aR$.
Using $\mathrm d u=d\,v^{d-1}\mathrm d v$ where $d=m+1$ is the integer,
and $\mathrm d t=v^m\mathrm d a+ma v^{m-1}\mathrm d v$,
substitute (IAL17) in (IAL5).
The nonscalar $dv$ terms cancel.  The scalar $dv$ term is
$-d\kappa v^{-d-1}+\rho a v^{-2}$ and the scalar $da$ term is
$-\rho v^{-1}$.  The remaining coefficients are
$\operatorname{diag}(b+1)dv/v$ and $-v^{-1}(N+av^mR)da$.
Since $D_v^{-1}ND_v=v^{-1}N$ and
$D_v^{-1}RD_v=v^{m-1}R$, this proves
\begin{equation}
 \nabla_{v,a}=d+dF\,I+D_v^{-1}dD_v
                  -D_v^{-1}J(a)D_v\,da.             \tag{IAL18}
\end{equation}
Here $d$ outside fractions denotes the exterior differential.  The
map $D_v$ sends this meromorphic connection to
$d+dF\,I-J(a)da$.  After tensoring with the inverse of the explicitly
retained exponential line $(d+dF)$, the factor is $d-J(a)da$.

There is also an actual chain explanation on $v\ne0$.
Let $S(P)=P(\rho+vx)$.  Then
$S(L_hP)=v^m(\partial_x+x^m-a)S(P)$.
The maps in degree zero and one
\begin{equation}
 S^0(P)=v^{m+1}S(P),\qquad
 S^1(Q\,ds)=vS(Q)\,dx                              \tag{IAL19}
\end{equation}
therefore intertwine $L_h$ with $\partial_x+x^m-a$ exactly.  On the
ordered top cohomology bases the resulting matrix is $D_v$.
This explains why dropping the one-form factor and replacing its
entries by $v^b$ would change the original chain map.

\subsection{The torsion extension, derived fibre, and dual action}
Put $S=\mathbb C[a]$, $\mathcal O=S[[v]]$, $S_0=\mathcal O/(v)=S$.
All statements in this paragraph are over this ring, with $a$ variable.
Multiplication by each $v^{b+1}$ is injective; hence
\begin{equation}
 0\longrightarrow\mathcal O^m\xrightarrow{D_v}\mathcal O^m
 \longrightarrow T_\rho\longrightarrow0,\qquad
 T_\rho=\bigoplus_{b=0}^{m-1}\mathcal O/(v^{b+1})\varepsilon_b.
                                                        \tag{IAL20}
\end{equation}
This module is free over $S$ of rank $m(m+1)/2$, with basis
$v^j\varepsilon_b$, $0\leq j\leq b<m$.  It is not a module of finite
length over $\mathcal O$: already its summand $S=\mathcal O/(v)$ has
the infinite strict descending chain $S\supset aS\supset a^2S\supset\cdots$.
The number $m(m+1)/2$ is its length after passing to $\mathbb C(a)[[v]]$
or after the explicitly specified fibre $a=a_0$ over $\mathbb C[[v]]$.

Let $E_\rho=\mathbb C[z]/z^m$.  The exact corrected socle isomorphism is
\begin{equation}
 \sigma_\rho:S\otimes_{\mathbb C}E_\rho
       \xrightarrow{\ \sim\ }\ker(v:T_\rho\to T_\rho),
 \qquad 1\otimes z^b\longmapsto v^b\varepsilon_b.    \tag{IAL21}
\end{equation}
Indeed in $\mathcal O/(v^{b+1})$ the annihilator of $v$ consists precisely
of $S v^b$; direct summation proves injectivity and surjectivity.
The source cannot be $E_\rho$ alone while $a$ remains variable.
After tensoring over $S$ with $S/(a-a_0)$, (IAL21) recovers the
specified $\mathbb C$-linear map from $E_\rho$.
The word socle here means this \emph{relative $v$-socle}.  The ordinary
module-theoretic socle over $\mathcal O$ is zero: $v$ is in its
Jacobson radical since every $1-vf$ is invertible, so its maximal
ideals are $(v,a-a_0)$; a vector in a simple submodule would be killed
by some $a-a_0$, which is impossible in this $S$-free module.

The resolution $[\mathcal O\xrightarrow{v}\mathcal O]$ of $S_0$ gives
the actual derived specialization
\begin{equation}
 H^{-1}(T_\rho\otimes^{\mathbf L}_{\mathcal O}S_0)=\ker v,
 \qquad H^0(T_\rho\otimes^{\mathbf L}_{\mathcal O}S_0)=T_\rho/vT_\rho.
                                                        \tag{IAL22}
\end{equation}
Equivalently reduce the two free terms of (IAL20) modulo $v$; the
differential is zero.  The connecting identification of its left
$S^m$ with $\ker v$ sends $e_b$ to $v^b\varepsilon_b$: lifting that
class and multiplying by $v$ gives $D_v e_b$, which proves the map
without choosing a noncanonical dimension comparison.

The source arithmetic action is $A(av^m)=\rho I+N+av^mR$, and
\begin{equation}
 L=D_vA(av^m)D_v^{-1}=\rho I+v(N+aR),\qquad
 LD_v=D_vA(av^m).                                    \tag{IAL23}
\end{equation}
This polynomial identity shows that $L$ preserves the image lattice
and descends to $T_\rho$.  For $b<m-1$ it sends
$v^b\varepsilon_b$ to
$\rho v^b\varepsilon_b+v^{b+1}\varepsilon_{b+1}$.
For the last basis vector its wrap is $av^m\varepsilon_0=0$ in
$\mathcal O/(v)\varepsilon_0$.  Thus (IAL21) intertwines the action
with $\rho I+N$.  On $T_\rho/vT_\rho$ the action is $\rho I$.
The two actions are the reductions of the two different free terms
of the resolution and are both retained.

After removing the specified exponential line, consider stability
under the regular \emph{$a$-direction} operator $\partial_a-J(a)$.
The smallest $\mathcal O$-submodule containing $D_v\mathcal O^m$ and
stable under this operator is $v\mathcal O^m$.  This module is stable
because $v$ is independent of $a$ and $J(a)$ has entries in $S$.
Conversely it must contain $v\varepsilon_0$ and then
$(-1)^b(\partial_a-J)^b(v\varepsilon_0)=v\varepsilon_b$ for $b<m$;
before the last column no $R$ term occurs.  Hence it contains all of
$v\mathcal O^m$.  The added quotient is
\begin{equation}
 v\mathcal O^m/D_v\mathcal O^m
       =\bigoplus_{b=0}^{m-1}\mathcal O/(v^b),
 \qquad \operatorname{rank}_S=\frac{m(m-1)}2.         \tag{IAL24}
\end{equation}
This is not stability under the full ordinary $v$ derivative:
$\partial_v(v\varepsilon_0)=\varepsilon_0\notin v\mathcal O^m$.
The $a$-direction qualification is essential.
For the full ordinary regular connection the minimum is instead
$\mathcal O^m$: applying $\partial_v$ first produces
$\varepsilon_0$, and successive regular $a$ derivatives then produce
all $\varepsilon_b$.  This lattice is stable under both regular
derivatives, and its added quotient over $D_v\mathcal O^m$ has
$S$-rank $m(m+1)/2$.  Neither assertion treats the retained irregular
exponential line as though its polar derivative preserved these lattices.

With $K=\mathcal O[v^{-1}]$, the torsion linking pairing is
\begin{equation}
 \langle[x],[y]\rangle=x^TD_v^{-1}y\pmod{\mathcal O}
     \quad\text{in }K/\mathcal O.                    \tag{IAL25}
\end{equation}
Changing either lift by $D_v$ times an integral vector changes the
value by an integral scalar.  On each summand, a homomorphism to
$K/\mathcal O$ is determined by the image of one, and its annihilation
by $v^{b+1}$ is exactly the class of $v^{-b-1}y$ modulo $\mathcal O$.
This constructs the inverse and proves perfectness over $\mathcal O$.
Its adjoint arithmetic matrix is
\begin{equation}
 L^\dagger=D_vL^TD_v^{-1}
         =\rho I+N^T+av^mR^T,
 \qquad \langle Lx,y\rangle=\langle x,L^\dagger y\rangle.
                                                        \tag{IAL26}
\end{equation}
It descends because $L^\dagger D_v=D_vL^T$.  On its ordinary fibre
the action is $\rho I+N^T$; on its socle the $N^T$ terms vanish since
$v^b\varepsilon_{b-1}=0$, and the wrap terms vanish as well.  Thus
that socle action is $\rho I$.
The coefficient $[v^{-1}]$ of (IAL25), equivalently the residue after
multiplication by the oriented differential $\mathrm d v$, gives a
perfect $S$-pairing between the original socle and the dual ordinary
fibre: on the displayed bases its matrix is the identity.  This proves
the exact exchange of the ordinary and
derived fibres under this linking duality.
For a positive real dilation parameter $\lambda$, separately from
the polynomial coordinate $a$, the original spectral exponential and
its contragredient are
\begin{equation}
 \lambda^\rho\sum_{j=0}^{m-1}\frac{(\log\lambda)^j}{j!}N^j,
 \qquad
 \lambda^{-\rho}\sum_{j=0}^{m-1}\frac{(-\log\lambda)^j}{j!}(N^T)^j.
                                                        \tag{IAL27}
\end{equation}
They follow by exponentiating $\rho I+N$ and its negative transpose.
No logarithm of the arbitrary affine coordinate $a$ is required.

\subsection{Strict spectral filtrations and the character-twisted maps}
On $E_\rho$ put
\begin{equation}
 M_jE_\rho=\operatorname{span}\{z^b:m-1-2b\leq j\},
 \quad Ne_b=e_{b+1},\quad H e_b=(m-1-2b)e_b.
                                                        \tag{IAL28}
\end{equation}
Then $N M_j\subset M_{j-2}$.  The only nonzero graded pieces have
indices $m-1,m-3,\ldots,1-m$, each of dimension one.  For $r\geq0$
the map $N^r:\operatorname{gr}^M_r\to\operatorname{gr}^M_{-r}$
is either the zero map between two zero spaces or sends its displayed
basis to its displayed basis with coefficient one.  It is an
isomorphism in both cases.  This proves the centered nilpotent
filtration identities directly on the original Jordan block.

For $0\leq r\leq m$ put $K_r=z^rE_\rho$ and use the induced
filtration on this subspace and on its actual quotient.  The basis
$z^{r+c}$ has ambient index $m-1-2r-2c$, while its centered index in
the length $m-r$ block is $m-r-1-2c$.  In the quotient, $z^b$, $b<r$,
has centered index $r-1-2b$.  Consequently
\begin{equation}
 \begin{gathered}
 M_jK_r=M^{\rm centered}_{j+r}K_r,\qquad
 M_j(E_\rho/K_r)=M^{\rm centered}_{j-(m-r)}(E_\rho/K_r),
 \\
 0\to M_jK_r\to M_jE_\rho\to M_j(E_\rho/K_r)\to0.
 \end{gathered}                                        \tag{IAL29}
\end{equation}
Every equality and exactness statement follows by those actual basis
indices, including the zero spaces when $r=0$ or $r=m$.
The dual filtration $M_j(E_\rho^*)=\{\ell:\ell(M_{-j-1}E_\rho)=0\}$
has dual basis weight $-(m-1-2b)$.  The negative transpose $-N^T$
lowers that weight by two, and dualizing the strict sequence in
(IAL29) gives its strict dual sequence with the reversed arrows.

For the normal meromorphic block (IAL13), use $e_b=y^bdy$ and write
$B e_b=(b+1)e_b/d$.  Define $N_y e_b=e_{b+1}$ with last image zero.
For each $j\geq0$, direct evaluation on $e_b$ gives
\begin{equation}
 [B,N_y^j]=(j/d)N_y^j,\qquad
 N_y^j:\mathcal V_\rho\longrightarrow
 \mathcal V_\rho\otimes\mathcal L_{-j/d},\qquad
 \mathcal L_{-j/d}=(\mathbb C((u)),d-(j/d)du/u).
                                                        \tag{IAL30}
\end{equation}
For $b+j<m$ this is the subtraction of the two actual diagonal
coefficients; for $b+j\geq m$ both sides vanish.  The source-to-target
connection defect is
$((B-j/d)N_y^j-N_y^jB)du/u=0$; the scalar irregular term is the same
on both sides.  This proves horizontality with the specified target
line.  For $0\leq j<m$, the kernel is the span of
$e_{m-j},\ldots,e_{m-1}$ and the image the span of
$e_j,\ldots,e_{m-1}$, so the rank is $m-j$.  For $j\geq m$ the map
is zero.  Its dual is the transpose map with the source and target
dual lines reversed, by transposing that same zero defect equation.

For $\zeta_d=\exp(2\pi i/d)$ the horizontal monodromy of the normal
block is $T e_b=\zeta_d^{-(b+1)}e_b$; its horizontal exponential
$\exp(-\kappa_\rho/u)$ is single-valued.  Therefore
\begin{equation}
 TN_y^jT^{-1}=\zeta_d^{-j}N_y^j,\qquad
 (\zeta_d^jT)N_y^j=N_y^jT.                           \tag{IAL31}
\end{equation}
This is a statement about the displayed normal connection; the
formal comparison of a general full packet has not been declared
convergent and does not eliminate its analytic Stokes maps.
For the next formula let $e_\rho\in E_h$ be the CRT idempotent that
is one on the complete $\rho$ primary factor and zero on the other
complete primary factors, and let $E_h(\rho)=e_\rho E_h$.  The square
restriction $k_\rho:E_h(\rho)\to\mathbb C[y]/y^m$ of the local
comparison is an isomorphism.  The rectangular row block in (IAL11)
is exactly $K_\rho(0)=k_\rho e_\rho$, with $e_\rho$ here denoting
the projection $E_h\to E_h(\rho)$; only $k_\rho$ is inverted.
At that marked jet fibre, the original arithmetic operator is
\begin{equation}
 k_\rho\bigl(A_h(0)|_{E_h(\rho)}\bigr)k_\rho^{-1}
   =\rho I+\sum_{j=1}^{m-1}([y^j]\zeta_\rho(y))N_y^j,
 \qquad [y]\zeta_\rho(y)=B_\rho(0)^{-1}.             \tag{IAL32}
\end{equation}
Substitution conjugates multiplication by $s$ to multiplication by
$\rho+\zeta_\rho(y)$ modulo $y^m$, and multiplication by the Jacobian
commutes with that multiplication.  This proves the formula.  Each
term in (IAL32) has its individual horizontal target (IAL30); the
sum is an operator on the marked coefficient fibre and is not a
horizontal endomorphism after discarding those target lines.

\subsection{The entire ordered tensor cone and its cyclic map}
For possibly unequal primary orders $m_i\geq1$ and roots $\rho_i$,
let $D=\operatorname{lcm}(m_i+1)$ and $r_i=D/(m_i+1)$.  Put
$S=\mathbb C[a_1,\ldots,a_k]$, $\mathcal O=S[[v]]$ and retain
$u=v^D$, $t_i=a_i v^{r_i m_i}$ and $z_i=v^{r_i}x_i$.
The original tensor basis is $z^{\mathbf b}=\prod_i z_i^{b_i}$,
$0\leq b_i<m_i$, with the ordered one-form orientation.  Set
\begin{equation}
 \ell_{\mathbf b}=\sum_i r_i(b_i+1),\quad
 D_{\otimes}=\operatorname{diag}(v^{\ell_{\mathbf b}}),\quad
 T_{\boldsymbol\rho}=\operatorname{coker}D_{\otimes}
   =\bigoplus_{\mathbf b}\mathcal O/(v^{\ell_{\mathbf b}})
                          \varepsilon_{\mathbf b}.   \tag{IAL33}
\end{equation}
This is the cokernel of the tensor lattice map; it is not a tensor
product of its one-factor cokernels.  Its exact socle isomorphism is
\begin{equation}
 S\otimes_{\mathbb C}\bigotimes_i\mathbb C[z_i]/z_i^{m_i}
 \longrightarrow\ker v,\qquad
 z^{\mathbf b}\longmapsto
       v^{\ell_{\mathbf b}-1}\varepsilon_{\mathbf b}.
                                                        \tag{IAL34}
\end{equation}
Each summand has annihilator $S v^{\ell_{\mathbf b}-1}$; this proves
bijectivity and keeps every ordered tensor direction.  The source
sum action and its target are
\begin{equation}
 A_{\Sigma}=\sum_i\bigl(\rho_i I+N_i+a_i v^{r_i m_i}R_i\bigr),
 \quad L_{\Sigma}=\left(\sum_i\rho_i\right)I
                     +\sum_i v^{r_i}(N_i+a_iR_i),
 \quad L_{\Sigma}D_{\otimes}=D_{\otimes}A_{\Sigma}.
                                                        \tag{IAL35}
\end{equation}
Raising $b_i$ increases $\ell$ by $r_i$, so a raising term in (IAL34)
has exactly the next socle exponent.  A wrap from $b_i=m_i-1$ to zero
lands with exponent $\ell_{\mathbf b}-1+r_i$, whereas the receiving
summand length is $\ell_{\mathbf b}-r_i(m_i-1)$; their difference is
$r_i m_i-1\geq0$.  Thus every wrap vanishes in its specified quotient.
It follows that (IAL34) intertwines $L_{\Sigma}$ with the entire
original operator $\sum_i\rho_i I+\sum_iN_i$.

Writing $M=\prod_i m_i$, summing the exponents gives
\begin{equation}
 \operatorname{rank}_S T_{\boldsymbol\rho}
 =\sum_{\mathbf b}\ell_{\mathbf b}
 =\frac M2\sum_i r_i(m_i+1)=\frac{kDM}{2}.             \tag{IAL36}
\end{equation}
For equal orders this is $k(m+1)m^k/2$.  These are again $S$-ranks,
generic $v$-lengths or fixed-parameter fibre lengths, not total
$\mathcal O$-lengths.  Placing the free tensor lattices in top degree
$k$ puts their mapping cone in degrees $k-1,k$; after derived
specialization its socle and ordinary fibre occur in those respective
degrees, with the actions just computed.

For equal order $m$, put $q=k(m-1)+1$ and $Z=\sum_i z_i$ in the actual
tensor algebra.  The cyclic coefficient inclusion is
\begin{equation}
 I_{\rm cyc}:\mathbb C[Z]/Z^q\longrightarrow
          \bigotimes_i\mathbb C[z_i]/z_i^m,
 \qquad Z^j\longmapsto
   \sum_{|\mathbf b|=j}\frac{j!}{\prod_i b_i!}z^{\mathbf b}.
                                                        \tag{IAL37}
\end{equation}
For every $0\leq j\leq k(m-1)$ a tuple with this degree and
$0\leq b_i<m$ exists (allocate $j$ successively to the $k$ slots).
Its displayed coefficient is nonzero in characteristic zero.  Since
different degrees have disjoint basis support, all columns are
independent.  The next power $Z^q$ is zero because no admissible tuple
has that degree.  This proves the domain, nilpotency order, and
injectivity, including $m=1$.  In the equal-order chart
\begin{equation}
 D_{\otimes}I_{\rm cyc}
 =I_{\rm cyc}\operatorname{diag}(v^k,v^{k+1},\ldots,v^{k+q-1}).
                                                        \tag{IAL38}
\end{equation}
Each summand in a column has total degree $j$, proving the identity
without deleting its multinomial coefficient.  In particular for
$m=k=2$, $(N_1+N_2)^2(1)=2z_1z_2$, while $z_1-z_2$ is a nonzero
four-dimensional tensor direction outside the three-dimensional
cyclic image.  The cyclic restriction therefore does not replace
the full tensor object.  The coefficient identities for individual
$N_i$ and their twists extend on the ordered tensor product; for
equal $d$ their sum has target the common line $\mathcal L_{-1/d}$,
and its powers have the multinomial coefficients in (IAL37).  For
unequal $d_i$ the individual target lines $\mathcal L_{-1/d_i}$
remain separately specified.

\subsection{Strict packet maps with their original parameter defects}
For monic $h\mid H$ put $b=H/h$.  On the two polynomial complexes
over $\mathbb C[u,t]$ define
\begin{equation}
 f_{hH}^0=I,\qquad
 f_{hH}^1=L_H\mathfrak H_h+M_b r_h.                   \tag{IAL39}
\end{equation}
Then $f_{hH}^1L_h=L_H$ by (IAL2), so this is a strict cochain map.
For $h\mid H\mid K$ put $c=K/H$.  The $H$ decomposition of
$f_{hH}^1P$ is exactly $L_H\mathfrak H_hP+b r_hP$, since the second
term has degree below $\deg H$.  Uniqueness of (IAL2) therefore gives
\begin{equation}
 f_{HK}^1f_{hH}^1P=L_K\mathfrak H_hP+cb r_hP=f_{hK}^1P.
                                                        \tag{IAL40}
\end{equation}
The degree-zero maps compose as identities.  Moreover direct
subtraction, with $M_b L_h$ meaning multiplication \emph{after}
$L_h$ rather than application of $L_h$ to $bP$, proves
\begin{equation}
 f_{hH}^1-M_b=(1-b)(u\partial_s-t)\mathfrak H_h,
 \qquad f_{hH}^1\bmod(u,t)=M_b.                       \tag{IAL41}
\end{equation}
At $u=t=0$, the differential division is ordinary monic division,
which also proves the specialization directly.

On the top free lattices the induced map is $I_b:P\mapsto bP$ on
degree-below-$\deg h$ representatives.  Its image is free and has
a free complement: the ordered set
$1,s,\ldots,s^{\deg b-1},b,bs,\ldots,bs^{\deg h-1}$ has distinct
degrees and leading coefficients one and is thus a triangular basis
of all polynomials of degree below $\deg H$.  If $\ell_h$ extracts
the coefficient of $s^{\deg h-1}$, ordinary reduction of $sP$ and
$sbP$ gives
\begin{equation}
 A_H(t)I_b-I_bA_h(t)=t(1-b)\ell_h,\qquad
 \nabla_t^H I_b-I_b\nabla_t^h=-\frac{t}{u}(1-b)\ell_h.
                                                        \tag{IAL42}
\end{equation}
Indeed the respective subtractions are $(h-t)\ell_h(P)$ and
$(H-t)\ell_h(P)$, whose difference after multiplication by $b$ is
$t(1-b)\ell_h(P)$.  Since $I_b$ is independent of $u$ and $t$,
the other exact defect is
\begin{equation}
 \nabla_u^H I_b-I_b\nabla_u^h
 =\frac{I_bC_h(t)-C_H(t)I_b}{u^2}
       +\frac{B_H(t)I_b-I_bB_h(t)}{u}.                \tag{IAL43}
\end{equation}
These formulas keep the original matrices and constants.  No
horizontal inclusion is inferred from the chain-map identity alone.
The displayed formula also repairs the literal form-feed character
in the received MB57, where its second \texttt{\string\frac} command
was corrupted.

For the actual theta source let $v_h=g/h$, $\mathcal T_h(P)$ be the
original inverse-Mellin source with Mellin transform $v_hP$, and
$g=2\xi$.  Since $v_H b=v_h$ as actual entire functions, injectivity
of the original Mellin realization proves
\begin{equation}
 \mathcal T_H(bP)=\mathcal T_h(P).                    \tag{IAL44}
\end{equation}
If $h$ and $b$ are complementary complete primary factors of the
actual packet, their CRT injection $\iota_{hH}$ includes the $h$
jets and assigns zero to the other complete primary factors.
At those other roots $v_hP$ vanishes to the full indicated orders;
at the $h$ roots its jet is $(j_hv_h)[P]_h$.  Thus
\begin{equation}
 J_H\mathcal T_h(P)=\iota_{hH}(j_hv_h)[P]_h.          \tag{IAL45}
\end{equation}
This last CRT interpretation has the stated complete-primary
complement hypothesis; the polynomial chain maps (IAL39)--(IAL43)
hold for every monic divisibility $h\mid H$.  Ordered tensoring of
those cochain maps retains their source degrees and exterior signs.
Extending the theta source maps coefficientwise in $u,t$, (IAL41)
also supplies their actual off-marked discrepancy:
\begin{equation}
 \mathcal T_Hf_{hH}^1(P)-\mathcal T_h(P)
 =\mathcal T_H\bigl((1-b)(u\partial_s-t)\mathfrak H_hP\bigr).
                                                        \tag{IAL45a}
\end{equation}
For ordered coefficient maps $I_1,\ldots,I_k$ the connection on their
external product is the sum of its slot connections, hence its defect
is exactly
\begin{equation}
 \nabla^{H}_{\otimes}\bigotimes_i I_i
 -\bigotimes_i I_i\nabla^{h}_{\otimes}
 =\sum_i I_1\otimes\cdots\otimes
 (\nabla^{H_i}I_i-I_i\nabla^{h_i})\otimes\cdots\otimes I_k.
                                                        \tag{IAL45b}
\end{equation}
This follows by the product rule; all $I_i$ have cochain degree zero,
so no extra orientation sign is introduced.

\subsection{Independent coefficient faces and exact receiving labels}
All preceding scalar extension, lattice, socle, packet, and ordered
tensor maps are coefficient maps over the original tau absolute
base.  For a finite coefficient mask $A$, use the direct sum of its
labelled copies, without identifying it with the terminal mask.  If
$A\subset B$, let $z_{AB}$ be identity on each admitted old slot and
zero in each newly admitted slot.  For any of the displayed linear
maps $F$, define $F_A=\bigoplus_{a\in A}F$.  Componentwise evaluation
proves
\begin{equation}
 F_Bz_{AB}=z_{AB}F_A.                                \tag{IAL46}
\end{equation}
At a present outer support this zero is its supported zero; lifting
a linear map by $v^\bullet\mapsto F(v)^\bullet$ and $\tau\mapsto\tau$
preserves that distinction.  An empty coefficient mask has the zero
coefficient module at its unchanged present outer label.  The
independent original leg mask remains a separate index throughout.

For the retained original source operator $D$ and any proper relation
subspace $W\subset V$, put $W_D=W+DW$.  Its exact quotient map is
\begin{equation}
 D_{W,W_D}:V/W\longrightarrow V/W_D,\qquad
 [\phi]_W\longmapsto[D\phi]_{W_D}.                   \tag{IAL47}
\end{equation}
If a representative changes by $w\in W$, its image changes by
$Dw\in W_D$, proving well-definedness.  Repeated action has receiving
space $W+DW+\cdots+D^rW$, by induction on $r$.  For the original
invertible dilation $U_\lambda$, the corresponding quotient map
$V/W\to V/U_\lambda W$ is $[\phi]\mapsto[U_\lambda\phi]$, with
inverse induced by $U_{\lambda^{-1}}$; changes of representative and
composition are checked by these same exact formulas.

The original theta operator is injective on its source $V$.  For
$x>0$, Schwarz decay makes
$\sum_{n,r\geq1}|\phi(nrx)|<\infty$: bounding the divisor count of
$j=nr$ by $j$ and using $|\phi(jx)|\leq C_xj^{-3}$ suffices.
Consequently the exact inversion is
\begin{equation}
 \frac12\sum_{n\geq1}\mu(n)\Theta\phi(nx)
 =\sum_{j\geq1}\phi(jx)\sum_{n\mid j}\mu(n)=\phi(x).
                                                        \tag{IAL47a}
\end{equation}
Here $\mu(n)$ is zero if a squared prime divides $n$ and is $(-1)^r$
if $n$ is a product of $r$ distinct primes.  The inner sum is one
for $j=1$ and is $(1-1)^r=0$ otherwise.  Thus $\Theta\phi=0$ implies
$\phi(x)=0$ for $x>0$; evenness and $\phi(0)=0$ finish the proof.

Finally let $r_h^0:E_h\to\mathbb C[s]$ be the original degree-below-
$\deg h$ representative map at $u=t=0$, and use the actual source
section $s_h^{\rm pol}(v)=\mathcal T_h(r_h^0v)$.  Direct ordinary
division gives $s r_h^0(v)-r_h^0(A_hv)=h\ell_h(v)$; since
$\mathcal T_h(hP)=\Theta(P(D)\phi_*)$, this proves
$Ds_h^{\rm pol}-s_h^{\rm pol}A_h=\Theta(\phi_*\ell_h)$.
The observed arithmetic coordinates are obtained with the retained
unit $\upsilon_h=j_h(g/h)\in E_h^\times$:
\begin{equation}
 s_h^{\rm obs}(x)=s_h^{\rm pol}(\upsilon_h^{-1}x),\qquad
 J_hs_h^{\rm obs}=I,\qquad
 Ds_h^{\rm obs}-s_h^{\rm obs}A_h
 =\Theta\bigl(\phi_*\ell_h M_{\upsilon_h^{-1}}\bigr).
                                                        \tag{IAL48a}
\end{equation}
Multiplication by $\upsilon_h^{-1}$ commutes with $A_h$, which proves
the last equality and keeps all Taylor coefficients of the unit.
For polynomial coordinates, passing to the receiving quotient in
(IAL47) gives the actual block map
\begin{equation}
 \begin{pmatrix}D_{W,W_D}&[\phi_*\ell_h]_{W_D}\\0&A_h\end{pmatrix}.
                                                        \tag{IAL48}
\end{equation}
In observed coordinates its upper-right entry is instead
$[\phi_*\ell_h M_{\upsilon_h^{-1}}]_{W_D}$.
For a specified complex-linear marked coefficient isomorphism
$K:E_h\to E'$ (in particular the square $K_h(0)$), choose the
transported section
$s'(w)=s_h^{\rm pol}(K^{-1}w)$.  The exact change of extension
coordinates is then $\operatorname{diag}(I,K)$.  Multiplying block
matrices shows that its upper-right term is
$[\phi_*\ell_h]_{W_D}K^{-1}$ and its bottom block is $KA_hK^{-1}$.
This diagonal assertion depends on the displayed transported section.
If the actual new section instead satisfies
$\widetilde s'(w)=s'(w)+\Theta\beta(w)$ for a specified
$\beta:E'\to V$, the exact change of coordinates is
\begin{equation}
 ([\phi]_W,v)\longmapsto
       ([\phi-\beta Kv]_W,Kv),\qquad
 \begin{pmatrix}I&-[\beta K]_W\\0&K\end{pmatrix}.
                                                        \tag{IAL49}
\end{equation}
Equality of represented source functions proves this formula directly.
Applying $D$ to $\widetilde s'$ gives its new upper-right entry
\begin{equation}
 [\phi_*\ell_hK^{-1}+D\beta-\beta(KA_hK^{-1})]_{W_D}.
                                                        \tag{IAL50}
\end{equation}
Thus a change of source section carries its actual additional
primitive; it cannot be replaced by a diagonal coefficient assertion.
The same proof applies after (IAL48a), with its full unit factors.
For the actual dilation primitive
$U_\lambda s_h^{\rm pol}-s_h^{\rm pol}\lambda^{A_h}
=\Theta\Psi_{\lambda,h}$, the receiving source label is
$U_\lambda W$ and the represented block map is
$\left(\begin{smallmatrix}U_\lambda&[\Psi_{\lambda,h}]_{U_\lambda W}
\\0&\lambda^{A_h}\end{smallmatrix}\right)$.
Its composition law follows by expanding two dilations:
\begin{equation}
 \Psi_{\lambda\mu,h}
 =U_\lambda\Psi_{\mu,h}
            +\Psi_{\lambda,h}\mu^{A_h}.              \tag{IAL51}
\end{equation}
The expansion first proves equality after applying $\Theta$; the
injectivity just proved in (IAL47a) proves equality of the actual
source primitives.
For a transported section it becomes $\Psi_{\lambda,h}K^{-1}$;
for the further section $\widetilde s'$ its additional term is
$U_\lambda\beta-\beta(K\lambda^{A_h}K^{-1})$.
These equations prove the exact primitive and receiving-label
transport at every coefficient face by (IAL46).
For the entire formal $K_h(u)$ of (IAL11), the precise corresponding
source objects are $E_h[[u]]$, $V[[u]]$, $W[[u]]$, and
$\mathscr B[[u]]$.  Extend $\Theta,D,U_\lambda$ and the original
section coefficientwise; for example
$\Theta[[u]](\sum_i u^i\phi_i)=\sum_i u^i\Theta\phi_i$.
Every coefficient of a product by $K_h(u)$ or its inverse uses a
finite sum, so the formal section
\begin{equation}
 s'_{\rm formal}=s_h^{\rm pol}[[u]]K_h(u)^{-1},\qquad
 A'_{\rm formal}=K_h(u)A_hK_h(u)^{-1}                 \tag{IAL52}
\end{equation}
and all the block identities (IAL48)--(IAL51) are well-defined
formal-series identities.  Reduction modulo $u$ recovers exactly
their marked source maps.  Inclusion of constant source functions
into these formal series is injective, as is $\Theta[[u]]$ by
(IAL47a) coefficientwise.  No analytic convergence or source-norm
identity is assigned to this formal completion.  Thus it constructs
the exact formal and marked bridge while preserving the original
actual source and its receiving relation label.
On the original absolute scalar objects
$G(R)=\{\tau\}\sqcup R^\bullet$, the coefficient specialization
$G(\mathbb C[a][[v]])\to G(\mathbb C[a])$ sends $v^\bullet$ to
$0^\bullet$ and $\tau$ to $\tau$.  This preserves the original
amplitude/support pair and the infinite integer realization target;
it does not equate the absolute base point with $v=0$.

The reconstructed calculation therefore supplies the exact formal
comparison, the ramified torsion object with its two different
specialization actions, every labelled tensor direction, and the
strict arrows back to the original packet and mixed source.  Its
formal coefficient estimate (IAL14) alone does not bound the
analytic Stokes gluing or the original signed four-endpoint theta
quantity uniformly in tensor order.

% END COMPLETE CORRECTED SOURCE algebra.tex


% BEGIN COMPLETE CORRECTED SOURCE finite_field.tex
% New intake finite-field repair and audit, 2026-09-14.
% Original NEW_BODY.tex is preserved without alteration.
\section{The finite-field primary family, its descent, and the actual boundary restriction}

This calculation supplies the missing definitions and proof behind the
intake's references to MB46--MB47, verifies MB48--MB52 and MB77--MB79
on their specified domains, and calculates the restriction map that is
needed when a boundary parameter varies.  The original polynomial
primitive, ordered factors, finite Taylor unit, coefficient specialization,
and the two different support observations are retained throughout.

\subsection{The coefficient point and the additive and Kummer conventions}

Let \(h_\rho(s)=(s-\rho)^m\) be one complete primary subpacket of the
original actual packet, with \(m\geq1\), \(d=m+1\).  The primitive is the
one with constant coefficient zero in the original \(s\) coordinate:
\[
 \Phi_{h_\rho}(s)=\frac{(s-\rho)^d-(-\rho)^d}{d}
 =\kappa_\rho+\frac{(s-\rho)^d}{d},
 \qquad \kappa_\rho=-\frac{(-\rho)^d}{d}.                 \tag{FF1}
\]
In particular, \(\kappa_\rho\) is not an independent coefficient.  The
symbol \(\Phi_h(\rho)\) for a larger packet is a different, separately
specified critical value; it is not substituted into (FF1).

Use the original unsplit finitely generated coefficient ring inside
\(\mathbb C\), with the full matrices of \(\upsilon_h\) and
\(\upsilon_h^{-1}\).  Enlarge it by the finitely many labelled primary
roots and their primitive coefficients that occur in this calculation.
Inverting the factorials through the largest phase degree makes every
displayed denominator a unit.  Whenever distinct labels are to remain
distinct under specialization, include the inverse of their original
nonzero differences.  These operations specify a finite type coefficient
model and do not replace the original complex coefficients.
This ring is taken before adjoining the idempotent coefficients for
all distinct tensor-sum centres.  The different full CRT localization
is retained explicitly below; its characteristic-\(p\) point is not
silently assumed to exist.
Fix an actual residue-field map
\[
 R\longrightarrow K_0=\mathbb F_{Q_0},\qquad
 \operatorname{char}K_0=p>\max_i d_i.                    \tag{FF2}
\]
Bars below mean this very map.  It specializes the matrix identity
\(M_{\upsilon_h}M_{\upsilon_h^{-1}}=1\) to the same identity, and makes
no assertion that the entire analytic function \(g=2\xi\) is a
finite-field function.

Fix a prime \(\ell\ne p\), an algebraically closed \(\ell\)-adic
coefficient field \(E\), and an embedding used for absolute values.
Fix a nontrivial additive character \(\psi_0:\mathbb F_p\to E^\times\).
For a finite extension \(K/\mathbb F_p\), write
\(\psi_K(x)=\psi_0(\operatorname{Tr}_{K/\mathbb F_p}x)\), and define
\(\mathcal L_\psi(f)\) by the character \(\psi_0\) of the torsor
\(Y^p-Y=f\).  Its geometric Frobenius trace is \(\psi_K(f)\).
A separately specified nontrivial additive character of \(K\) is
\(\psi_K(cf)\) for its unique \(c\in K^\times\); that coefficient must
then multiply every phase and every Gauss sum below.  For simplicity of
notation the displayed formulas use the trace character just fixed.

For an ordered collection put \(D=\operatorname{lcm}(d_i)\).
Make the explicit residue extension \(K/K_0\) of degree
\(e=\operatorname{ord}_D(Q_0)\), taking \(e=1\) when \(D=1\), so
\(D\mid Q-1\), \(Q=Q_0^e\).  For a multiplicative character
\(\chi:K^\times\to E^\times\) with \(\chi^d=1\), let
\(\mathcal K_\chi\) be the Kummer sheaf whose geometric Frobenius
trace at \(x\in K^\times\) is \(\chi(x)\).  Its trace over an extension
is \(\chi\circ\operatorname{Nm}(x)\).  Its multiplication isomorphism
is the one induced by multiplication of roots in the power torsor.
It retains the fibre \(\mathcal K_\chi(d)\); that fibre has Frobenius
\(\chi(d)\), and is not replaced by the unit line.

All characterwise decompositions below are made over this specified
splitting field.  Before this extension the sheaf itself is defined by
its original direct image over \(K_0\).  Frobenius can permute the
character summands.  More explicitly, if \(\delta_\zeta\) is a deck
transformation and \(F_0\) is geometric Frobenius, then
\[
 F_0\delta_\zeta F_0^{-1}=\delta_{\zeta^{Q_0^{-1}}},
 \qquad
 e_\theta=\frac1D\sum_{\zeta\in\mu_D}
                 \theta(\zeta)^{-1}\delta_\zeta,
 \qquad F_0e_\theta F_0^{-1}=e_{\theta^{Q_0}}.           \tag{FF3}
\]
Here \(Q_0^{-1}\) is taken modulo \(D\).  The second equality follows
from the first by substituting
\(\eta=\zeta^{Q_0^{-1}}\) in the finite sum.  Thus the original
Frobenius action, not a list of artificially descended one-dimensional
lines, is retained.  A Gauss eigenvalue calculated over \(K\) is an
eigenvalue of \(F_0^e\).  If its absolute value is \(Q^{w/2}\), an
eigenvalue \(\alpha\) of \(F_0\) above it has
\(|\alpha|^e=Q_0^{ew/2}\), hence \(|\alpha|=Q_0^{w/2}\).

\subsection{The missing marked Kummer--Gauss calculation}

On \(\mathbb G_{m,u}\times\mathbb A^1_t\) let
\[
 \pi:\mathbb A^1_s\times
          (\mathbb G_{m,u}\times\mathbb A^1_t)
             \longrightarrow\mathbb G_{m,u}\times\mathbb A^1_t,
 \qquad
 \mathcal F_{\rho,m}
   =R^1\pi_!\mathcal L_\psi
        \left(\frac{\Phi_{h_\rho}(s)-ts}{u}\right).
                                                               \tag{FF4}
\]
All coefficients in (FF4) are their specified reductions.
For \(t=0\), the original translation \(z=s-\bar\rho\) and the
projection formula give
\[
 \mathcal F_{\rho,m}|_{t=0}
 =\mathcal L_\psi(\bar\kappa_\rho/u)\otimes
       R^1\pi_{z!}\mathcal L_\psi(z^d/(du)).             \tag{FF5}
\]
The following is an isomorphism of sheaves, including its constant
Frobenius lines:
\[
 \begin{split}
 \mathcal F_{\rho,m}|_{t=0}
 &\simeq\mathcal L_\psi(\bar\kappa_\rho/u)\otimes
   \bigoplus_{\substack{\chi^d=1\\\chi\ne1}}
       \mathcal K_\chi(u)\otimes\mathcal K_\chi(d)
                           \otimes\mathcal G_{\chi,\psi},\\
 \mathcal G_{\chi,\psi}
 &=H_c^1(\mathbb G_{m,\overline K},
                 \mathcal K_\chi(x)\otimes\mathcal L_\psi(x)),
 \qquad
 F_Q|_{\mathcal G_{\chi,\psi}}=-G_K(\chi,\psi),\\
 G_K(\chi,\psi)&=\sum_{x\in K^\times}\chi(x)\psi_K(x).
 \end{split}                                                  \tag{FF6}
\]
Each \(\mathcal G_{\chi,\psi}\) has dimension one.

\begin{proof}
Let \(f:\mathbb A^1_z\to\mathbb A^1_y\), \(y=z^d\), and
\(j_y:\mathbb G_m\hookrightarrow\mathbb A^1_y\).
The degree-\(d\) finite power map carries its actual deck action.
On \(\mathbb G_m\) its direct image is the regular representation,
split by the finite character idempotents.  At \(y=0\) its fibre
has a single point and a trivial deck action.  It follows, on all
geometric stalks and with the actual projectors, that
\[
 f_*E_{\mathbb A^1_z}
 \simeq E_{\mathbb A^1_y}\oplus
       \bigoplus_{\substack{\chi^d=1\\\chi\ne1}}
                         j_{y!}\mathcal K_\chi(y).     \tag{FF7}
\]
For a nontrivial character its stalk and costalk boundary invariants
vanish, so \(j_{y!}\mathcal K_\chi=j_{y*}\mathcal K_\chi\) in degree
zero.  The finite map has no higher direct image.

Tensor (FF7) with \(\mathcal L_\psi(y/(du))\) and push with compact
supports over \(y\).  The trivial summand has no compact cohomology:
it is a nonconstant linear Artin--Schreier sheaf on \(\mathbb A^1\).
For example its \(H_c^0\) is zero by nonproper connected support,
its \(H_c^2\) is the dual of invariant \(H^0\) with twist and is
zero by nontrivial monodromy, and its Euler characteristic is
\(1-\operatorname{Swan}_\infty=0\).
On a nontrivial Kummer summand substitute \(y=du\,x\).
This is an isomorphism of the relative punctured lines on \(u\ne0\),
and multiplication in the Kummer torsor gives exactly
\(\mathcal K_\chi(d)\otimes\mathcal K_\chi(u)
 \otimes\mathcal K_\chi(x)\).
Proper-support base change and the projection formula give (FF6)
at the level of complexes.

For each such summand \(H_c^0=0\); Poincare duality gives \(H_c^2=0\)
because its dual has nontrivial tame monodromy at zero.  Its degree-one
Frobenius trace is therefore \(-G_K(\chi,\psi)\).
The following calculation proves that this Gauss sum is nonzero.
Under the fixed complex embedding,
\[
 \begin{split}
 |G_K(\chi,\psi)|^2
 &=\sum_{r\in K^\times}\chi(r)
       \sum_{y\in K^\times}\psi_K((r-1)y)\\
 &=(Q-1)-\sum_{r\ne1}\chi(r)=Q.
 \end{split}\tag{FF8}
\]
The last equality uses \(\sum_{r\in K^\times}\chi(r)=0\).
Thus every nontrivial character summand has positive dimension.
Deligne's degree-\(d\), one-variable theorem, Weil II 3.7.2.3,
gives total dimension \(d-1\); there are precisely \(d-1\)
nontrivial characters.  Since all their cohomology is in degree one,
each dimension is one.  This finishes the sheaf-level construction,
including its Frobenius sign.
\end{proof}

As a check of the constants rather than a substitute for that
construction, the compact-support trace identity reads
\[
 -\operatorname{Tr}(F_Q,\mathcal F_{\rho,m,u,0})
 =\sum_{z\in K}\psi_K((\bar\kappa_\rho+z^d/d)/u)
 =\psi_K(\bar\kappa_\rho/u)
       \sum_{\substack{\chi^d=1\\\chi\ne1}}
                \chi(du)G_K(\chi,\psi).               \tag{FF9}
\]
The constant \(1\) from \(z=0\) cancels the trivial-character
sum \(\sum_{y\ne0}\psi_K(y/(du))=-1\).  Reversing the
degree-one trace sign or dropping \(\chi(d)\) violates this identity.

\subsection{The actual inertia, every ordered invariant, and the Tate term}

At \(u=0\) the Kummer summands of (FF6) have nontrivial finite tame
characters.  The line \(\mathcal L_\psi(\bar\kappa_\rho/u)\) has
nontrivial wild character of break one when \(\bar\kappa_\rho\ne0\),
and is trivial when \(\bar\kappa_\rho=0\).
Wild and tame characters cannot cancel one another, because the tame
character restricts trivially to wild inertia.  Therefore
\[
 V_{\rho,m}^{I_u}=0,\qquad
 N_{I_u}=0\ \hbox{on a finite-index inertia subgroup}. \tag{FF10}
\]
The first assertion uses the nontrivial tame character even when the
wild line is trivial.  The second assertion follows because every
inertia summand factors through a finite group, killed by a finite
extension.  It does not identify \(N_{I_u}\) with multiplication by
the original spectral jet \(s-\rho\).

For \(k\) labelled primary factors, write
\(\boldsymbol\chi=(\chi_1,\ldots,\chi_k)\), with
\(\chi_i^{d_i}=1,\chi_i\ne1\), and set
\[
 \kappa_\Sigma=\sum_i\bar\kappa_i,\quad
 \chi_\Sigma=\prod_i\chi_i,\quad
 C_{\boldsymbol\chi}
   =\bigotimes_{i=1}^k
        \bigl(\mathcal K_{\chi_i}(d_i)
                   \otimes\mathcal G_{\chi_i,\psi}\bigr).
                                                               \tag{FF11}
\]
The factor labelled \(\boldsymbol\chi\) is
\[
 \mathcal L_\psi(\kappa_\Sigma/u)\otimes
        \mathcal K_{\chi_\Sigma}(u)\otimes C_{\boldsymbol\chi},
 \qquad
 F_Q|_{C_{\boldsymbol\chi}}
       =(-1)^k\prod_i\chi_i(d_i)G_K(\chi_i,\psi).        \tag{FF12}
\]
Consequently its invariant projector is the identity precisely when
\(\kappa_\Sigma=0\) in the specified residue field and
\(\chi_\Sigma=1\), and is zero otherwise.  In particular
\[
 V_0=V^{I_u}
   =\begin{cases}
      \displaystyle\bigoplus_{\chi_\Sigma=1}
                         C_{\boldsymbol\chi},&\kappa_\Sigma=0,\\
      0,&\kappa_\Sigma\ne0 .
     \end{cases}                                             \tag{FF13}
\]
Every ordered tuple is an individual summand even when its character
or Frobenius scalar coincides with another tuple's scalar.

Here is the rank formula also for unequal orders.  Embed all characters
in the cyclic group of \(D\)-th characters, using their actual orders.
For a generator \(\zeta_D\), a nontrivial \(d_i\)-th character has
exponent \(l_iD/d_i\), \(1\le l_i\le d_i-1\).
The number of tuples with product one is
\[
 r_{(d_i)}
  =\frac1D\sum_{r=0}^{D-1}
       \prod_{i=1}^k
       \begin{cases}d_i-1,&d_i\mid r,\\-1,&d_i\nmid r.\end{cases}
                                                               \tag{FF14}
\]
Indeed the indicator of a zero total exponent is
\(D^{-1}\sum_{r=0}^{D-1}
       \zeta_D^{r\sum_i l_iD/d_i}\).
Sum it over all independently labelled \(l_i\); the \(i\)-th inner
geometric sum is \(d_i-1\) when \(d_i\mid r\), and \(-1\) otherwise.
For identical \(m=d-1\), this becomes exactly
\[
 r_{m,k}=\frac{m^k+m(-1)^k}{m+1}.                       \tag{FF15}
\]
This is the dimension of (FF13) only when the original wild sum
\(\kappa_\Sigma\) is zero.

Let \(j_u\) denote the inclusion of the punctured strict local
\(u\)-trait, and \(i_u\) its closed point.  The exact local cohomology is
\[
 H^0(i_u^*Rj_{u*}\mathcal F)=V_0,\qquad
 H^1(i_u^*Rj_{u*}\mathcal F)=V_0(-1),\qquad
 H^r=0\quad(r\ne0,1).                                  \tag{FF16}
\]
To prove this, take wild inertia invariants first.  A pro-\(p\)
group has exact invariants on these \(\ell\)-adic representations:
at finite levels its averaging operator is available because
\(p\ne\ell\), and the continuous cohomology in positive degree vanishes.
On tame inertia, use the two-term continuous cochain complex after
removing the prime-to-\(\ell\) part.  A nontrivial finite character
has no degree-zero or degree-one cohomology: when it survives that
prime-to-\(\ell\) projector its topological generator has invertible
\(T-1\) over \(E\).  On a trivial summand, continuous homomorphisms
from the tame \(\ell\)-quotient give \(E(-1)\).  The natural
Frobenius action on this homomorphism space is multiplication by \(Q\).
This proves (FF16), including the Tate twist.  The usual sheaf sequence
is, on the whole \(u\)-curve,
\[
 0\longrightarrow j_{u!}\mathcal F
 \longrightarrow j_{u*}\mathcal F
 \longrightarrow i_{u*}V_0\longrightarrow0.            \tag{FF17}
\]
It is exact on the open by the identity and at the missing point by
the computation \(j_{u!}\mathcal F=0\), \(j_{u*}\mathcal F=V_0\).

Every Gauss line in (FF12) has weight one by (FF8).
The Kummer constant has finite order.  Hence \(V_0\) has weight \(k\)
and \(V_0(-1)\) has weight \(k+2\).  This conclusion also descends to
the original finite field by (FF3); no splitting over that field is
assumed.  The local logarithm is zero and its centred monodromy
filtration is concentrated at index zero.  The weight statement agrees
with Weil II 1.8.4 and 1.8.8 on this actual lisse finite-field family.

\subsection{Opposite phase, the dual line, and the ordered cohomological sign}

Retain \(\psi^{-1}_K(x)=\psi_K(-x)\).
For \(\chi\ne1\), direct summation gives
\[
 G_K(\chi,\psi)G_K(\chi^{-1},\psi^{-1})=Q.              \tag{FF18}
\]
To verify it without changing a constant, write \(x=ry\):
the product is
\(\sum_{r\ne0}\chi(r)\sum_{y\ne0}\psi_K((r-1)y)\),
which is exactly the right side of (FF8).
Thus the two negative degree-one eigenvalues multiply to \(Q\).
The two Kummer factors, including their fibres at \(d\), are inverse;
the two scalar Artin--Schreier phases are opposite.  It follows that
\[
 \mathcal F_{\rho,m,\psi}^{\vee}
       \simeq\mathcal F_{\rho,m,\psi^{-1}}(1)
                       \quad\hbox{on the marked \(u\)-curve}.   \tag{FF19}
\]
This is the actual duality pairing.  Equivalently, Poincare duality
pairs \(H_c^1\) for the positive phase with \(H^1\) for the negative
phase into \(E(-1)\).  For the rank-one phase on the \(s\)-line the
only omitted point is infinity, where its pole degree is \(d\),
prime to \(p\); it has no inertia cohomology.  Thus the map
\(H_c^1\to H^1\) is an isomorphism and gives the stated
compact-support pairing.  The same argument applies with the
parameter \(t\) retained.

Tensor in the original order.  The target becomes \(E(-k)\),
and the dual of the ordered \(k\)-fold sheaf is its opposite-phase
realization twisted by \((k)\).  If the tensor pairing is written
as the product of the \(k\) degree-one pairings, moving each second
factor past the subsequent first factors makes
\(\sum_{i=1}^k(k-i)=k(k-1)/2\) interchanges.  The sign is therefore
\((-1)^{k(k-1)/2}\), in addition to the \((-1)^k\) already present
in the compact-support trace (FF12).  These two signs have different
origins and neither is removed.

\subsection{The ramified family with its actual parameter retained}

On \(v\ne0\) put \(u=v^d\), \(t=av^{d-1}\), and use the actual
coordinate isomorphism \(s=\bar\rho+vx\).  Its inverse is
\(x=(s-\bar\rho)/v\).  From (FF1),
\[
 \frac{\Phi_{h_\rho}(s)-ts}{u}
     =\frac{\bar\kappa_\rho}{v^d}
         -\frac{\bar\rho a}{v}+\frac{x^d}{d}-ax.        \tag{FF20}
\]
Define the regular parameter sheaf by its full direct image
\[
 \mathcal A_m=R^1\pi_{x!}\mathcal L_\psi(x^d/d-ax),
        \qquad \pi_x:\mathbb A^1_x\times\mathbb A^1_a
                                  \longrightarrow\mathbb A^1_a .
                                                               \tag{FF21}
\]
The coordinate isomorphism, the projection formula, and
proper-support base change give
\[
 \mathcal F_{\rho,m}|_{u=v^d,t=av^{d-1}}
   \simeq\mathcal L_\psi
       (\bar\kappa_\rho/v^d-\bar\rho a/v)
                            \otimes\operatorname{pr}_a^*\mathcal A_m.
                                                               \tag{FF22}
\]
The multiplication law of the Artin--Schreier torsors proves the
separation of the two phases in (FF22).  The change of coordinate is
a bijection of the relative affine lines; it introduces no finite-field
counting factor.  In the separate de Rham realization its differential
is \(ds=v\,dx\), retained by the original diagonal period map.

The sheaf \(\mathcal A_m\) is lisse of rank \(m\), pointwise pure of
weight one, and all its other \(R^j\pi_{x!}\) vanish.
Here is the precise application of Deligne.  The parameter map sends
\(a\) to the actual polynomial \(x^d/d-ax\) in the universal affine
space of polynomials of degree at most \(d\).  Its leading form is
\(x^d/d\), with \(p\nmid d\).  Its projective zero hypersurface in
\(\mathbb P^0\) is empty and hence smooth.  Thus the entire parameter
map lands in Deligne's good open.  Weil II 3.7.3 makes every universal
proper-support cohomology sheaf lisse there.  Base change gives lissity
on \(\mathbb A^1_a\), and 3.7.2.3, in affine dimension one, gives
the rank, concentration and purity at every closed point.
Consequently the regular factor in (FF22) extends across the whole
divisor \(v=0\) as the lisse sheaf
\(\operatorname{pr}_a^*\mathcal A_m\).

For a \emph{fixed} finite-field parameter \(a=a_0\), the punctured
\(v\)-trait has the scalar inertia character in (FF22).
Its break is \(d\) when \(\bar\kappa_\rho\ne0\), is one when
\(\bar\kappa_\rho=0,\bar\rho a_0\ne0\), and is zero with trivial
inertia when both coefficients are zero.  Its Swan conductor is
the break times \(m\), and the invariant space is respectively zero
or the complete fibre \((\mathcal A_m)_{a_0}\).
For this actual singleton, however, the middle case is empty:
(FF1) and \(p>d\) give
\[
       \bar\kappa_\rho=0\quad\Longleftrightarrow\quad\bar\rho=0.
                                                               \tag{FF23}
\]
The break-one case does occur after the tensor cancellations below.
This distinction prevents an inconsistent choice of primary constants.

For completeness, the integer break calculation is as follows.
Over the geometric residue field, a scalar Artin--Schreier equation
\(Y^p-Y=cv^{-b}+\text{lower poles}\), \(c\ne0\), \(p\nmid b\),
has degree \(p\): a Laurent series with a pole of order \(r>0\)
has highest pole order \(pr\) after applying \(Z\mapsto Z^p-Z\),
whereas a series with no pole has no pole after that operation.
The extension is totally ramified.  Normalize its valuation by
\(w(v)=p\); then \(w(Y)=-b\).  Choose
\(1\le s<p\), \(-bs\equiv1\pmod p\), and
\(r=(1+bs)/p\).  The element \(v^rY^s\) has valuation one.
A nonidentity translation \(Y\mapsto Y+c_0\) changes it by a
leading term \(sc_0v^rY^{s-1}\) of valuation \(1+b\).
Every further binomial term has larger valuation.
The lower ramification groups are therefore the degree-\(p\)
group through index \(b\) and trivial afterwards.
The Herbrand function is the identity up to its sole break, so
the upper break is also \(b\).  This proves every stated scalar
break and its rank multiple.

For equal degrees and independently labelled parameters, the
ordered product is exactly
\[
 \mathcal L_\psi
     \left(\frac{\sum_i\bar\kappa_i}{v^d}
               -\frac{\sum_i\bar\rho_i a_i}{v}\right)
    \otimes\bigotimes_{i=1}^k\operatorname{pr}_{a_i}^*\mathcal A_m .
                                                               \tag{FF24}
\]
At a fixed parameter tuple its regular rank is \(m^k\), its weight
is \(k\), and its full covered invariant fibre survives precisely
when the two displayed polar coefficients vanish.
If they do not both vanish, the largest nonzero pole has order
\(d\) or one and the Swan conductor is its order times \(m^k\).

For unequal primary orders, keep \(D=\operatorname{lcm}(d_i)\),
put \(e_i=D/d_i\), and use
\[
 u=v^D,\quad t_i=a_i v^{D-e_i},\quad
           s_i=\bar\rho_i+v^{e_i}x_i .
                                                               \tag{FF25}
\]
The same substitution and projection formula give the complete phase
and regular factor
\[
 \mathcal L_\psi
     \left(\frac{\sum_i\bar\kappa_i}{v^D}
             -\sum_i\frac{\bar\rho_i a_i}{v^{e_i}}\right)
       \otimes\bigotimes_i\operatorname{pr}_{a_i}^*\mathcal A_{m_i}.
                                                               \tag{FF26}
\]
Every \(e_i<D\) and \(p\nmid D\).  Group equal exponents without
identifying their ordered slots:
\(c_D=\sum_i\bar\kappa_i\) and
\(c_e=-\sum_{i:e_i=e}\bar\rho_i a_i\).
On a fixed parameter trait, invariants are the entire regular fibre
of dimension \(\prod_i m_i\) exactly when every \(c_e\) is zero.
Otherwise the largest exponent \(b\) with \(c_b\ne0\) is prime
to \(p\), the break is \(b\), and the Swan conductor is
\(b\prod_i m_i\), by the valuation calculation just given.

On the marked slice \(\boldsymbol a=0\), the cover descends to
the original \(u\)-curve.  Over the splitting field it is Galois,
with tame quotient of order \(D\), and
\[
 V^{I_u}=(V^{I_v})^{I_u/I_v}.                           \tag{FF27}
\]
A vector fixed by the right side is fixed by every coset and
the normal subgroup, hence by all of \(I_u\); the reverse implication
is immediate.  Thus (FF27) is an equality of actual fixed spaces.
Its finite quotient projector is exactly the product-character
projector of (FF13)--(FF14).  The covered regular fibre is not
silently substituted for the original unramified invariant fibre.

\subsection{Boundary restriction does not commute with varying parameters}

The preceding inertia calculations concern fixed parameter traits.
The canonical map relating them to the boundary on the full parameter
space is generally not an isomorphism.  The following calculation
uses an ordered tensor of an original actual primary.

Since a nontrivial zeta zero \(\rho\) is nonzero, include its actual
inverse in the finite coefficient model and use a good point (FF2)
of that model.  Thus \(\bar\rho\ne0\).  Take \(k=p\) ordered
copies of \(h_\rho\); no original slot is identified with another.
In the residue field \(p\bar\kappa_\rho=0\).  Pull (FF24) back
along the actual parameter slice
\[
       a_1=a,\qquad a_2=\cdots=a_p=0.                  \tag{FF28}
\]
Writing \(X=\mathbb A^1_v\times\mathbb A^1_a\),
\(U=\mathbb G_{m,v}\times\mathbb A^1_a\), and
\(j:U\hookrightarrow X\), the resulting sheaf is exactly
\[
 \mathcal F
   =\mathcal L_\psi(-\bar\rho a/v)\otimes\mathcal B(a),
 \qquad
 \mathcal B(a)=\mathcal A_m(a)\otimes
                         (\mathcal A_m(0))^{\otimes(p-1)}.
                                                               \tag{FF29}
\]
The regular factor extends lisse to \(X\) and has rank \(m^p\).
At the origin \(o=(0,0)\) one has
\[
                 (j_*\mathcal F)_{\bar o}=0.           \tag{FF30}
\]

\begin{proof}
Let \(R^{\mathrm{sh}}\) be the strict henselization of
\(\mathcal O_{X,o}\).  It is a regular local domain with parameters
\(v,a\).  Every lisse finite-level coefficient sheaf pulled back from
\(X\) is constant on this strictly henselian local scheme.  Passing
through its compatible finite levels trivializes \(\mathcal B\)
there as a local system.

The height-one prime \(v=0\) defines a discrete valuation of
\(\operatorname{Frac}(R^{\mathrm{sh}})\).  Its residue contains the
nonzero image of \(a\), because \(a\notin(v)\).  The function
\(-\bar\rho a/v\) has pole order exactly one in that valuation.
It cannot be \(Z^p-Z\) in this fraction field: a pole of \(Z\)
would make the highest pole of \(Z^p-Z\) a positive multiple of
\(p\), while an integral \(Z\) makes \(Z^p-Z\) integral.
Thus its Artin--Schreier torsor has degree \(p\), and its
nontrivial character sheaf is nontrivial on the connected scheme
\(\operatorname{Spec}R^{\mathrm{sh}}[1/v]\).
The invariant subspace of this rank-one character is zero.
Tensoring with the constant fibre of \(\mathcal B\) preserves that
zero invariant space.  By the definition of the geometric stalk of
\(j_*\), this is (FF30).
\end{proof}

Let \(c:\mathbb A^1_v\to X\) be \(a=0\), let
\(j_0:\mathbb G_{m,v}\to\mathbb A^1_v\), and let \(c_U\)
be the induced map of opens.  The usual unit of adjunction and the
cartesian square construct the canonical map
\[
 c^*Rj_*\mathcal F\longrightarrow Rj_{0*}c_U^*\mathcal F.
                                                               \tag{FF31}
\]
Explicitly, its unit and counit construction is
\[
 \begin{split}
 c^*Rj_*\mathcal F
 &\longrightarrow Rj_{0*}j_0^*c^*Rj_*\mathcal F\\
 &=Rj_{0*}c_U^*j^*Rj_*\mathcal F
       \longrightarrow Rj_{0*}c_U^*\mathcal F .
 \end{split}                                                  \tag{FF31a}
\]
The first arrow is the unit for \(j_0^*,Rj_{0*}\), and the last
is induced by the counit \(j^*Rj_*\mathcal F\to\mathcal F\).
On the open this counit is the identity isomorphism.
On this fixed trait the phase of (FF29) is zero, so
\(c_U^*\mathcal F=(\mathcal A_m(0))^{\otimes p}\) is constant.
The degree-zero stalk map of (FF31) at zero is therefore exactly
\[
       0\longrightarrow(\mathcal A_m(0))^{\otimes p}.   \tag{FF32}
\]
Its kernel is zero and its cokernel has dimension \(m^p\).
The right-hand derived stalk also has degree one equal to
\((\mathcal A_m(0))^{\otimes p}(-1)\), by (FF16).
No assertion about the higher-degree left-hand stalk is required
for this calculated failure of degree-zero base change.

This example preserves (FF1): it uses the actual cancellation
\(p\bar\kappa_\rho=0\) and the independently varying actual
parameters.  A singleton with
\(\bar\kappa_\rho=0,\bar\rho\ne0\) would contradict (FF1) and
is not an example in the program.  Equations (FF29)--(FF32)
show why fixed-trait invariant fibres cannot simply be assembled
into the ordinary boundary stalks on the full parameter space.

The coefficient arrows used by this example are exactly
\[
 R[u,\boldsymbol t,\boldsymbol s]\longrightarrow
 K_0[u,\boldsymbol t,\boldsymbol s]\longrightarrow
 K[v,v^{-1},\boldsymbol a,\boldsymbol x],
 \quad
 s_i\longmapsto\bar\rho+vx_i,\quad
 u\longmapsto v^d,\quad t_i\longmapsto a_i v^{d-1}.
                                                               \tag{FF32a}
\]
The first arrow is coefficient reduction, the second includes the
specified splitting field extension and the explicit ramified
substitution.  It is the unsplit ordered tensor model.  In particular,
it does not use an inverse of \(p\), a factorial \(r!\) divisible
by \(p\), or a projector separating all tensor-sum centres.

The relation to the full offcritical-quartet CRT vertex is also exact.
At a stage \(k\ge p\), two of its retained centres are
\[
 \lambda_0=k\rho,\qquad
 \lambda_1=(k-p)\rho+p(1-\overline\rho),\qquad
 \lambda_0-\lambda_1=p(\rho+\overline\rho-1).
                                                               \tag{FF32b}
\]
Here the bar on \(\overline\rho\) means complex conjugation,
as explicitly written, and the offcritical hypothesis makes the
last factor nonzero.  These are complex coefficients at this
separate vertex, before any residue reduction.
Let \(A_{\rm CRT}\subset\mathbb C\) contain the actual coefficients
of its idempotent \(e(Z)=\sum_j e_jZ^j\), with
\(e(\lambda_0)=1\) and \(e(\lambda_1)=0\).
Retain its divided difference
\[
 \mathscr D=\sum_{j\ge1}e_j
             \sum_{r=0}^{j-1}\lambda_0^{j-1-r}\lambda_1^r,
 \qquad
       1=p(\rho+\overline\rho-1)\mathscr D.             \tag{FF32c}
\]
Multiplying the inner sum by \(\lambda_0-\lambda_1\)
telescopes to \(\lambda_0^j-\lambda_1^j\).  Summing gives
\(e(\lambda_0)-e(\lambda_1)=1\), proving (FF32c).
Consequently the exact common characteristic-zero coefficient
extension \(R\hookrightarrow R'=R[A_{\rm CRT}]\subset\mathbb C\)
contains
\(p^{-1}=(\rho+\overline\rho-1)\mathscr D\).
There is no unital extension of the map (FF2) from \(R'\) to
characteristic \(p\): it would send the equality (FF32c) to
\(1=0\).  Thus the finite-field example is attached to the common
unsplit source by (FF32a); it is not asserted to exist at that
incompatible full CRT corner.  This reproduces the exact
JPP.45--46 and TO.J18 coefficient obstruction without deleting
either vertex from the total object.

\subsection{The exact relation to the larger packet and to coefficient faces}

For a larger original monic packet \(H\), the marked-primary
Artin--Schreier sheaf is related to its full polynomial phase by
an explicit tensor kernel before proper-support pushforward.
On the same \((s,u,t)\) space put
\[
 \mathcal D_{h_\rho,H}
   =\mathcal L_\psi((\Phi_H(s)-\Phi_{h_\rho}(s))/u).
 \qquad
 \mathcal L_H
    \simeq\mathcal L_{h_\rho}\otimes\mathcal D_{h_\rho,H}.
                                                               \tag{FF33}
\]
The isomorphism is induced by addition of the two
Artin--Schreier coordinates; their right sides add to the exact
full phase.  For \(h_\rho\mid H\mid K\) the kernels satisfy
\(\mathcal D_{h_\rho,K}\simeq
  \mathcal D_{h_\rho,H}\otimes\mathcal D_{H,K}\)
by the same addition, with its associative torsor law.
Thus the corresponding direct image is exactly
\[
       R\pi_!\mathcal L_H
       \simeq R\pi_!
              (\mathcal L_{h_\rho}\otimes\mathcal D_{h_\rho,H}).
                                                               \tag{FF34}
\]
The kernel depends on \(s\); it is not moved outside this
direct image by a false projection formula.
This gives the precise sheaf-level relationship to the larger
actual packet.  The original polynomial complex also has its
separate strict packet maps and their nonzero connection defects.
A formal characteristic-zero primary splitting does not, by
itself, identify (FF34) with a direct sum of the marked-primary
finite-field sheaves (FF4).

For an original outer support label and a coefficient face \(A\),
take the direct sum of the corresponding labelled sheaves or
coefficient realizations.  If \(A\subset B\), the face map sends
a vector to the same labelled coordinate in \(B\) and inserts
zero in every other coordinate.  Each character projector in
(FF13) is a scalar identity or zero on that same labelled
summand.  Consequently, on every vector,
\[
        P_B\,\iota_{AB}=\iota_{AB}\,P_A .               \tag{FF35}
\]
The kernel of \(P_A\) is the direct sum of precisely the
noninvariant labelled summands; its image is the invariant
sum.  This is an exact direct-sum splitting for the marked
trait representations with finite inertia.  It does not
replace the varying-parameter comparison (FF31) by a projector
having the same fibre formula.
All these maps lift under the original
\(v\mapsto v^\bullet\), \(\tau\mapsto\tau\) rule.
A represented zero maps to the receiving supported zero;
external absence remains external absence.
The amplitude/support pair therefore remains injective, and the
original infinite arithmetic square has not been altered.

\subsection{Primary reference and mathematical scope}

Deligne, \emph{La conjecture de Weil II},
Publications mathematiques de l'IHES \textbf{52} (1980), 137--252,
3.7.2.3 and 3.7.3, supplies the explicitly applied concentration,
rank, purity and lissity theorem.  His 1.8.4 and 1.8.8 concern
the local monodromy filtration and invariant boundary stalk
of the actual lisse finite-field representation.
The original French pages 175--176 and 215--216 were inspected
for this calculation; the source is
\url{https://www.numdam.org/item/PMIHES_1980__52__137_0.pdf}.
The Kummer decomposition, Gauss signs, fixed-trait cohomology,
ramification breaks, unequal-order invariant count, and the
restriction map (FF31)--(FF32) are proved above.
These are operations on the specified residue-field sheaves and
their original coefficient model.  They supply no new identification
of finite-field Frobenius with the original complex dilation
operator or with its spectral nilpotent.

% END COMPLETE CORRECTED SOURCE finite_field.tex


% BEGIN COMPLETE CORRECTED SOURCE period_metric.tex
% Reconstructed and corrected period/source intake, 2026-09-14.
% The delivered NEW_BODY.tex and INTEGRATION.patch both omit MB20--44.
\section{The actual entire theta amplitude and the finite period observation}
\label{sec:ipm-entire-source}

This calculation uses the actual packet
\(h(s)=\prod_\rho(s-\rho)^{m_\rho}\), with all complete orders,
\(d_h=\deg h\ge1\), \(\Phi_h'=h\), \(\Phi_h(0)=0\),
\(g=2\xi\), \(v_h=g/h\), and
\(\upsilon_h=j_hv_h\in E_h^\times\), where
\(E_h=\mathbb C[s]/(h)\). No root outside this actual packet is
introduced. The analytic statements have these data as their domain;
they do not assert that an off-critical packet has been found.
The finite-source map is the original injective map
\(\mathcal T_hP=P(D)F_h\), with Mellin transform \(v_hP\).
Thus all maps below on \(\mathcal T_h\mathcal P_N\) are defined
unambiguously through that injective source parametrization.

\subsection{An explicit growth bound proves the global integration domain}

Write
\[
 \psi(x)=\sum_{n\ge1}e^{-\pi n^2x},\qquad
 C_\theta=(1-e^{-3\pi})^{-1},\qquad
 \beta_R=(R+1)/2.
 \tag{IPM.1}
\]
The original theta representation, also recorded in DLMF 25.5.13--14,
gives the entire identity
\[
 g(s)=1+s(s-1)\int_1^\infty
       \bigl(x^{s/2}+x^{(1-s)/2}\bigr)\psi(x)\,\frac{dx}{x}.
 \tag{IPM.2}
\]
For completeness, the identity follows initially for \(\Re s>1\)
by integrating each Gaussian:
\(\int_0^\infty\psi(x)x^{s/2}dx/x
=\pi^{-s/2}\Gamma(s/2)\zeta(s)\).
Split the integral at one and use the original Poisson identity
\(\psi(x)=x^{-1/2}\psi(1/x)+(x^{-1/2}-1)/2\).
The elementary integral on \((0,1)\) is
\(1/(s-1)-1/s=1/(s(s-1))\). Multiplication by \(s(s-1)\)
proves (IPM.2) in that half-plane. Its right side is entire, since its
integral and every coefficient derivative converge uniformly on compact
sets by Gaussian decay. Analytic continuation proves the identity
everywhere, including \(s=0,1\).

Since \(n^2\ge1+3(n-1)\), for \(x\ge1\)
one has \(\psi(x)\le C_\theta e^{-\pi x}\). Consequently
\[
 \begin{gathered}
 |g(s)|\le\mathcal X(R):=
  1+2C_\theta R(R+1)\pi^{-\beta_R}\Gamma(\beta_R,\pi)\\
 \le 1+2C_\theta R(R+1)\pi^{-\beta_R}\Gamma(\beta_R),
 \qquad |s|\le R,\ R\ge1.
 \end{gathered}
 \tag{IPM.3}
\]
Indeed both powers of \(x\) in the integral have absolute value at most
\(x^{(R+1)/2}\), and \(|s(s-1)|\le R(R+1)\).
The upper incomplete gamma is the literal integral after \(z=\pi x\).
An elementary bound sufficient here is
\(\Gamma(\beta)\le2(2\beta)^\beta\) for \(\beta\ge1\).
To prove it, write its integrand as
\((t^{\beta-1}e^{-t/2})e^{-t/2}\), bound the first factor by
its maximum \([2(\beta-1)/e]^{\beta-1}\), and integrate the second.
At \(\beta=1\) use the value one for the zeroth power. This maximum
is at most \((2\beta)^\beta\). In particular
\(\log\mathcal X(R)=O(R\log(R+2))\); the explicit function
\(\mathcal X\), rather than a substituted constant, is retained in
all tail integrals.

Let \(R_h=\max_\rho|\rho|\). For \(|s|=r\ge\max(1,2R_h)\),
\[
 |h(s)|\ge(r/2)^{d_h},\qquad
 |v_h(s)|\le 2^{d_h}r^{-d_h}\mathcal X(r).
 \tag{IPM.4}
\]
The complete-divisor hypothesis makes \(v_h\) entire, so it is bounded
on the remaining compact disc. Cauchy's integral formula on radius-one
circles gives the same \(\exp(O(r\log(r+2)))\) growth for each fixed
derivative. The formula retains both the degree and every original root
in its compact-disc and exterior bounds.

Fix \(u\ne0\), \(t\in\mathbb C\), a branch of \(\arg u\), and
the original BC rays and oriented contours
\[
 \theta_j=\frac{\arg u+\pi+2\pi j}{d_h+1},\quad
 \ell_j(r)=re^{i\theta_j},\quad
 \Gamma_j=-\ell_0+\ell_j\quad(1\le j\le d_h),\quad
 f_t=\Phi_h-ts.
 \tag{IPM.5}
\]
Their leading real exponent is
\(-r^{d_h+1}/((d_h+1)|u|)\). If
\(f_t(s)=s^{d_h+1}/(d_h+1)+\sum_{\nu=0}^{d_h}c_\nu s^\nu\),
then, for \(r\ge1\), its full real exponent is at most
\[
 -\frac{r^{d_h+1}}{(d_h+1)|u|}
       +\frac{\sum_{\nu=0}^{d_h}|c_\nu|r^\nu}{|u|}.
 \tag{IPM.6}
\]
Combining (IPM.3)--(IPM.6) proves absolute convergence of
\[
 \mathcal I_{u,t}(a)=
  \left(\int_{\Gamma_j}e^{f_t(s)/u}a(s)\,ds\right)_{j=1}^{d_h}
 \tag{IPM.7}
\]
for \(a=v_hP\), every polynomial \(P\), every derivative used below,
and every entire quotient by \(h\) constructed below. On a compact
parameter set with \(u\ne0\), the same domination proves coefficient
differentiation and the needed contour deformations inside the fixed
decaying sectors. This is a proof of the domain for the actual theta
amplitude. It makes no assertion for arbitrary undeclared rays.

\subsection{The exact finite-unit comparison and its nonzero receiving term}

For any entire amplitude \(a\) in this domain, let
\(r_h(a)\) be its unique polynomial of degree below \(d_h\) having
all its original primary jets. Existence and uniqueness are ordinary
Hermite interpolation at the complete packet. Thus
\[
 a=r_h(a)+h\,\mathcal Q_ha,\qquad
 \mathcal Q_ha:=\frac{a-r_h(a)}h
 \quad\hbox{is entire}.
 \tag{IPM.8}
\]
The same bounds just proved apply to \(\mathcal Q_h(v_hP)\) and
its derivatives: outside the packet disc division is bounded by (IPM.4),
inside it every apparent pole is removable. Define
\(\Pi_h\) as the matrix of (IPM.7) on \(1,s,\ldots,s^{d_h-1}\),
in that order. The following three maps have the explicitly common
original source \(\mathcal T_h\mathcal P_N\):
\[
 \begin{aligned}
 \mathcal O^{\rm coeff}(\mathcal T_hP)
    &=\Pi_h M_{\upsilon_h^{-1}}J_h\mathcal T_hP
      =\Pi_h[P]_h,\\
 \mathcal O^{\rm jet}(\mathcal T_hP)
    &=\Pi_hJ_h\mathcal T_hP
      =\Pi_h M_{\upsilon_h}[P]_h
      =\mathcal I_{u,t}(r_h(v_hP)),\\
 \mathcal O^{\rm entire}(\mathcal T_hP)
    &=\mathcal I_{u,t}(v_hP).
 \end{aligned}
 \tag{IPM.9}
\]
Here \(\Pi_h[P]_h\) means evaluation on the declared remainder frame;
it does not identify \(\mathcal H_{h,u,t}\) with \(E_h\) as a
connection. The equalities in the first two lines follow exactly from
\(J_h\mathcal T_hP=\upsilon_h[P]_h\). They retain the full finite
unit and its inverse.

For an oriented finite chain \(\Gamma\), write
\(\partial_\Gamma b=\sum_p n_p b(p)\), where
\(\partial\Gamma=\sum_p n_p[p]\). Integration of the exact
derivative \(u(e^{f_t/u}A)'=e^{f_t/u}(uA'+(h-t)A)\), with
\(A=\mathcal Q_h(v_hP)\), gives
\[
 \begin{aligned}
 \mathcal I_\Gamma(v_hP)-\mathcal I_\Gamma(r_h(v_hP))
   &=t\mathcal I_\Gamma(A)-u\mathcal I_\Gamma(A')
          +u\partial_\Gamma(e^{f_t/u}A).
 \end{aligned}
 \tag{IPM.10}
\]
All orientations and the positive sign of the endpoint term follow
from this derivative. On the global contours (IPM.5), its two outer
endpoint limits are zero by (IPM.3)--(IPM.6), and the two copies of
the common origin cancel. Therefore the global original-source map is
\[
 \begin{aligned}
 \mathcal O^{\rm entire}-\mathcal O^{\rm jet}
   &=(t\mathcal I_{u,t}\mathcal Q_h
            -u\mathcal I_{u,t}\partial_s\mathcal Q_h)M_{v_h},\\
 \mathcal O^{\rm entire}(\mathcal T_hP)
       -\mathcal O^{\rm coeff}(\mathcal T_hP)
   &=\Pi_h(M_{\upsilon_h}-I)[P]_h
       +t\mathcal I_{u,t}(A)-u\mathcal I_{u,t}(A').
 \end{aligned}
 \tag{IPM.11}
\]
In the first line, the right side acts on the polynomial parametrization
\(P\) of the source. Thus this is an actual map with a convergent
target, not a comparison asserted from a formal coefficient matrix.
In particular finite inverse-unit cancellation in MB67 proves the first
line of (IPM.9), and does not delete either the other finite-unit term
or the explicitly calculated receiving term in (IPM.11).

The boundary action is equally explicit. On an original relation
\(P=hQ\), the first two observations vanish, whereas
\[
 \mathcal O^{\rm entire}(\mathcal T_h(hQ))
     =t\mathcal I_{u,t}(v_hQ)-u\mathcal I_{u,t}((v_hQ)').
 \tag{IPM.12}
\]
At a truncated local chain the additional endpoint is the one in
(IPM.10) with \(A=v_hQ\). The relation source
\(\mathcal T_h(hQ)=\Theta(Q(D)\phi_*)\) is the original one.
Equation (IPM.12) specifies its receiving observation; it does not
discard its source norm or its proper-source cohomology class.

For a finite cutoff \(N\ge d_h-1\), set
\[
 \mathscr R_N^{\rm per}
    =\mathcal O^{\rm entire}\mathcal T_h
                    (h\mathcal P_{N-d_h})\subset\mathbb C^{d_h}.
 \tag{IPM.12a}
\]
Negative polynomial degree denotes the zero space. The strongest
canonical receiving quotient at this cutoff is the literal quotient
\(\mathbb C^{d_h}/\mathscr R_N^{\rm per}\). The maps
\[
 \begin{array}{ccc}
 \mathcal P_{N-d_h}&\xrightarrow{\ \times h\ }&\mathcal P_N\\
 {\scriptstyle Q\mapsto\mathcal I_{u,t}(v_hhQ)}\downarrow&&
       \downarrow{\scriptstyle P\mapsto\mathcal I_{u,t}(v_hP)}\\
 \mathscr R_N^{\rm per}&\lhook\joinrel\longrightarrow&\mathbb C^{d_h}
 \end{array}
 \tag{IPM.12b}
\]
form an exact cochain square: both composites on \(Q\) are the same
integral, evaluated by (IPM.12). The target differential is inclusion;
the source differential is multiplication by the original \(h\).
Its induced degree-one map is
\[
 \overline{\mathcal O}_N:E_h\longrightarrow
             \mathbb C^{d_h}/\mathscr R_N^{\rm per},\qquad
 [P]_h\longmapsto[\mathcal I_{u,t}(v_hP)].
 \tag{IPM.12c}
\]
Two lifts differ by an original \(h\)-relation and hence by precisely
the receiving subspace in (IPM.12a). Conversely a linear target quotient
through which the entire map factors via \(E_h\) must kill every such
image; it therefore factors uniquely through (IPM.12c). This proves
its universal quotient property and not merely its well-definedness.
Its kernel is exactly
\[
 \frac{\{P\in\mathcal P_N:
              \mathcal I_{u,t}(v_hP)\in\mathscr R_N^{\rm per}\}}
       {h\mathcal P_{N-d_h}}.
 \tag{IPM.12d}
\]
Thus the remaining quotient by this kernel is isomorphic to
\(\mathcal I_{u,t}(v_h\mathcal P_N)/\mathscr R_N^{\rm per}\), by
the same representative map in both directions. This construction
does not presume that either the receiving subspace or the kernel is
zero. Cutoff inclusions enlarge \(\mathscr R_N^{\rm per}\), giving
the specified quotient arrows at successive cutoffs.

\subsection{A finite iterated reduction retains every derivative remainder}

On the same entire-amplitude domain define
\(\mathcal R_{u,t}a=t\mathcal Q_ha-u\partial_s\mathcal Q_ha\).
Repeated substitution of (IPM.10) proves, for every integer \(L\ge1\),
\[
 \begin{aligned}
 \mathcal I_\Gamma(a)
   &=\sum_{r=0}^{L-1}\mathcal I_\Gamma
                 \bigl(r_h(\mathcal R_{u,t}^{\,r}a)\bigr)
      +\mathcal I_\Gamma(\mathcal R_{u,t}^{\,L}a)\\
   &\quad+u\sum_{r=0}^{L-1}\partial_\Gamma
              \bigl(e^{f_t/u}\mathcal Q_h\mathcal R_{u,t}^{\,r}a\bigr).
 \end{aligned}
 \tag{IPM.13}
\]
The proof is induction: substitute (IPM.10) with amplitude
\(\mathcal R_{u,t}^{\,L}a\) for the final integral, retaining its
endpoint. Each operation preserves the proved entire-amplitude domain.
Thus global endpoint terms vanish separately; local ones remain. This
finite identity makes no convergence claim for its infinite expansion.
At \(t=0\) its \(r\)-th coefficient is the complete operator
\((-u)^r(\partial_s\mathcal Q_h)^r\), with its original order.

\subsection{Ordered tensor maps and every original support face}

For the \(i\)-th original factor let \(A_i\) denote the entire map
in (IPM.9) on polynomial coordinates, \(B_i\) its jet map, and
\(E_i=A_i-B_i\), explicitly (IPM.10) or (IPM.11). The full ordered
external tensor identity is
\[
 \bigotimes_{i=1}^kA_i-\bigotimes_{i=1}^kB_i
 =\sum_{i=1}^k
     \left(\bigotimes_{j<i}A_j\right)\otimes E_i\otimes
     \left(\bigotimes_{j>i}B_j\right).
 \tag{IPM.14}
\]
It follows by adding and subtracting consecutive products in which the
leftmost \(i\) factors have changed. It is an identity of multilinear
maps, hence of their tensor-linear extensions, without choosing a tensor
decomposition of the input. Every integral is absolutely convergent by
the proved one-factor bounds, so Fubini identifies the tensor target
with the original ordered product contours and orientation
\(ds_1\wedge\cdots\wedge ds_k\). The maps have degree zero and
introduce no extra tensor permutation sign.

Apply (IPM.14) to the literal polynomial \(P(\sum_i s_i)\).
Its arithmetic source remains
\(\mathcal V_NP=(\prod_i v_h(s_i))P(\sum_i s_i)\) of PGS.4.
On the finite observation the exact map is
\[
 \mathcal O^{\rm coeff}_N
    =\Pi_h^{\otimes k}M_{\upsilon_h^{-1}}^{\otimes k}
                     J^{(k)}\mathcal V_N
    =\Pi_h^{\otimes k}\iota J_N,
 \qquad \eta=M_{\upsilon_h}^{\otimes k}\iota.
 \tag{IPM.15}
\]
Here \(\iota[P]\) is the original successive remainder of
\(P(\sum_i s_i)\) in the entire ordered tensor coefficient space,
before cyclic restriction. The identity follows by multiplying the
specified inverse units into \(\eta\). More precisely put
\(\mathcal O_N^{\rm jet}=\Pi_h^{\otimes k}\eta J_N\).
Then (IPM.14) evaluates
\(\mathcal O_N^{\rm entire}-\mathcal O_N^{\rm jet}\), and
\[
 \mathcal O_N^{\rm entire}-\mathcal O_N^{\rm coeff}
 =(\mathcal O_N^{\rm entire}-\mathcal O_N^{\rm jet})
   +\Pi_h^{\otimes k}
       (M_{\upsilon_h}^{\otimes k}-I)\iota J_N .
 \tag{IPM.15f}
\]
This follows by subtracting the two finite maps explicitly.
The entire observation relative to the coefficient observation
therefore retains this additional full-unit term as well as the
derivative remainder. No independent tensor direction is identified
with the cyclic image.

The entire observation on the original cyclic relation is computed
without asserting that it vanishes. Let
\(P_0(s_1,\ldots,s_k)=\chi(\Sigma)Q(\Sigma)\), with
\(\Sigma=\sum_i s_i\), and successively put
\[
 P_i=\operatorname{rem}_{h(s_i)}P_{i-1},\qquad
 A_i=(P_{i-1}-P_i)/h(s_i)\quad(1\le i\le k).
 \tag{IPM.15a}
\]
These are ordinary monic divisions in the stated order. Since \(\chi\)
annihilates the original sum action on the full tensor coefficient
space, the final simultaneous remainder \(P_k\) is zero. Consequently
\(P_0=\sum_i h(s_i)A_i\), with every coefficient of the successive
quotients retained. Set \(a_i=(\prod_jv_h(s_j))A_i\), and let
\(\mathcal I^{(k)}\) be the ordered product version of (IPM.7), with
phase \(\sum_i(\Phi_h(s_i)-t_i s_i)/u\). Integration by parts in
the \(i\)-th variable gives the exact vector identity
\[
 \mathcal I^{(k)}\bigl((\prod_jv_h(s_j))\chi(\Sigma)Q(\Sigma)\bigr)
 =\sum_{i=1}^k
       \left(t_i\mathcal I^{(k)}(a_i)
                 -u\mathcal I^{(k)}(\partial_{s_i}a_i)\right).
 \tag{IPM.15b}
\]
Each term has the global convergence established above; the polynomial
dependence on the other variables permits Fubini. On truncated chains
the right side also contains \(u\) times the \(i\)-th endpoint
evaluation integrated over all other variables. In the original
exterior complex its primitive is
\((-1)^{i-1}a_i\,ds_1\wedge\cdots\widehat{ds_i}\cdots\wedge ds_k\).
Wedge insertion contributes a second \((-1)^{i-1}\), so the product
sign is positive, exactly as in the original theta relation.

Let \(O_N:\mathcal P_N\to\mathbb C^{d_h^k}\) be this full entire
observation. Its finite receiving subspace
\(\mathscr R_{N,k}^{\rm per}=O_NB_N\mathcal P_{N-q}\)
is computed by (IPM.15a)--15b. Replacing \(h,\mathcal T_h,E_h\)
in (IPM.12a)--12d by \(\chi,\mathcal V_N,E\) gives the exact
cochain square and universal quotient
\[
 E\longrightarrow
           \mathbb C^{d_h^k}/\mathscr R_{N,k}^{\rm per},
 \quad[P]\longmapsto[O_NP].
 \tag{IPM.15c}
\]
All source, target, kernel and cutoff maps in that assertion are
proved by the identical monic relation presentation; (IPM.15b)
supplies their actual receiving vectors.

At a present coefficient face \(A\), take the direct sum of each
displayed map over its unchanged labelled slots. For \(A\subset B\),
zero insertion commutes with every map because all are linear and
send zero to zero in each existing slot. This proves the square on
coefficients and, by \(f^\tau(0^\bullet)=0^\bullet\), on supported
zeros over the original \(\mathfrak b_\tau\). External absence remains
the separate \(\tau\). Original leg masks, Fourier transports and
proper relation-source labels are unchanged. On their induced top
cohomology the residual kernel is still
\((dC_{\rm full}^{k-1}\cap C_\lambda^k)/dC_\lambda^{k-1}\);
its map sends the same representative to its full-source boundary,
and its inverse is the same representative modulo proper boundaries.
The additional entire observation has the receiving terms (IPM.12)--14
on those representatives. This proves the bridge without assuming it
factors through a quotient on which those terms need not vanish.

\subsection{The existing sectorial incidence maps act on this entire source}

On the marked slice \(t_1=\cdots=t_k=0\), retain the already constructed SPG endpoint-incidence matrices
\(E_\theta\in\operatorname{GL}_{d_h}(\mathbb Z)\) and the ordered
global BC cycles. The SPG sectorial cycles are the integral combinations
with row matrix \(E_\theta\). Integration of any entire one-form on a
finite compact contour respects these combinations and compact contour
deformations by Cauchy's theorem. The same remains true for the
present noncompact cycles: on each retained decaying sector the large
arcs tend to zero by (IPM.3)--6. On a lifted inverse-polynomial ray
the original phase is a linear decaying variable, while
\(|s|=O(r^{1/(d_h+1)})\); the actual \(v_h\) factor is bounded by
\(\exp(O(r^{1/(d_h+1)}\log(r+2)))\), and hence is dominated by the
linear exponential decay. At a finite critical point the local
coordinate is analytic and its inverse derivative has exponent
\(1/(m_\rho+1)-1>-1\), which is integrable. Thus the same incidence
identities apply to the entire-amplitude integral, with every repeated
critical block and every original path sign retained.

Writing \(O_N^\theta\) for the sectorial entire observation, the exact
maps on the unchanged source are consequently
\[
 \begin{gathered}
 O_N^\theta=E_\theta^{\otimes k}O_N,\qquad
 O_N^{\theta,+}=\mathfrak S_{\theta_0}^{\otimes k}O_N^{\theta,-},
 \quad \mathfrak S_{\theta_0}=E_{\theta_0+}E_{\theta_0-}^{-1},\\
 \mathscr R_{N,k}^{\mathrm{per},\theta}
      =E_\theta^{\otimes k}\mathscr R_{N,k}^{\rm per},\qquad
 [z]\longmapsto[E_\theta^{\otimes k}z]:
 \frac{\mathbb C^{d_h^k}}{\mathscr R_{N,k}^{\rm per}}
       \xrightarrow{\sim}
 \frac{\mathbb C^{d_h^k}}{\mathscr R_{N,k}^{\mathrm{per},\theta}}.
 \end{gathered}
 \tag{IPM.15d}
\]
For distinct leg directions use their ordered external matrices in
place of the repeated tensor. The quotient inverse is induced by the
ordered inverse matrices. These identities prove preservation of the
full observation kernel and transport of its precise relation image;
they do not assert injectivity of \(O_N\). Its sectorial source Gram is
exactly
\[
 (O_N^\theta)^*O_N^\theta
      =O_N^*(E_\theta^*E_\theta)^{\otimes k}O_N.
 \tag{IPM.15e}
\]
The original SPG30 matrix \(E_\theta\Pi_hU_h^{-1}\), with
\(U_h=M_{\upsilon_h}\), remains the valid finite coefficient
observation \(\mathcal O^{\rm coeff}\) of (IPM.9), after applying
\(J_h\). Equations (IPM.11), (IPM.14), and (IPM.15d) are the
explicit additional map relating it to \(O_N^\theta\). Hence no
existing Stokes matrix needs to be invented or recomputed to obtain
this relation; the new work is the actual entire-source receiving
map and its original-norm Gram.

\section{Reconstruction of the analytic primary period calculation}
\label{sec:ipm-primary-period}

This section reconstructs the analytic statements missing from the
delivered MB20--44 span. Both delivered copies of INTEGRATION.patch
also jump from MB19 to MB45; these proofs are new reconstruction from
the retained definitions, not a claim to have recovered missing bytes.
For the actual complete primary \(h_\rho(s)=(s-\rho)^m\), put
\(d=m+1\),
\(\kappa=-(-\rho)^d/d\), \(u=v^d\), \(t=av^m\),
\(s=\rho+vx\). The exact phase is
\[
 \frac{\Phi_{h_\rho}(s)-ts}{u}
    =F(v,a)+x^d/d-ax,\qquad
 F(v,a)=\kappa/v^d-\rho a/v.
 \tag{IPM.16}
\]
No constant, Jacobian or parameter has been removed.

Let \(\theta_j=(\pi+2\pi j)/d\),
\(\ell_j(r)=re^{i\theta_j}\),
\(\Gamma_j=-\ell_0+\ell_j\), \(1\le j\le m\), and define
\[
 (H_m(a))_{jb}=\int_{\Gamma_j}e^{x^d/d-ax}x^b\,dx,
 \quad0\le b<m,\qquad
 D_v=\operatorname{diag}(v^{b+1})_{b=0}^{m-1}.
 \tag{IPM.17}
\]
All these integrals are entire in \(a\) by their fixed leading decay.
The translation matrix from the original ascending \(s\) frame to
the ascending \(z=s-\rho\) frame is
\( (C_\rho)_{bj}=\binom jb\rho^{j-b}\) for \(b\le j\), zero
otherwise. Binomial expansion gives its inverse \(C_{-\rho}\) and
determinant one. Substitution in each oriented integral proves
\[
 \Pi_\rho(v,a)=e^{F(v,a)}H_m(a)D_vC_\rho.
 \tag{IPM.18}
\]
The powers \(b+1\) include \(ds=v\,dx\). This is the period matrix
of polynomial one-forms; entire theta amplitudes are related to it by
(IPM.9)--14, retaining all additional terms.

Let \(Ne_b=e_{b+1}\) for \(b<m-1\), \(Ne_{m-1}=0\), and
\(Re_{m-1}=e_0\), \(Re_b=0\) otherwise. Differentiation in \(a\)
inserts \(-x\). The final column reduces by
\(\int_\Gamma(x^m-a)e^{x^d/d-ax}dx
=[e^{x^d/d-ax}]_{\partial\Gamma}=0\). Therefore
\[
 H_m'(a)=-H_m(a)(N+aR).
 \tag{IPM.19}
\]
For \(m\ge2\) the trace on the right is zero; for \(m=1\),
\(N=0\) and \(R=I\), so its trace is \(a\).

Write \(\zeta=e^{2\pi i/d}\), \(\beta_b=(b+1)/d\), and
\(W_{jb}=\zeta^{j(b+1)}-1\). Direct radial substitution gives
\[
 H_m(0)_{jb}
   =d^{\beta_b-1}e^{\pi i\beta_b}\Gamma(\beta_b)W_{jb}.
 \tag{IPM.20}
\]
The complete determinant, with its phase in the original order, is
\[
 \det H_m(a)=
 (2\pi)^{m/2}e^{3\pi i m(m+1)/4}
 \begin{cases}1,&m\ge2,\\ e^{-a^2/2},&m=1.\end{cases}
 \tag{IPM.21}
\]
To prove its constant, the full \(d\)-by-\(d\) Fourier matrix
\(V_{jr}=\zeta^{jr}\), \(0\le j,r<d\), becomes a block matrix
with lower block \(W\) after subtracting its row zero from every
other row. Thus \(\det W=\det V\). Its Vandermonde factors are
\[
 \zeta^q-\zeta^p
  =2i\sin\frac{\pi(q-p)}d\,
                         e^{\pi i(p+q)/d},\quad0\le p<q<d.
\]
Every sine here is positive. The sum of \(p+q\) over pairs is
\(d(d-1)^2/2\), and there are \(d(d-1)/2\) pairs. The determinant
phase is therefore \(\pi(d-1)(3d-2)/4\). Its absolute value is
\(d^{d/2}\), since \(V^*V=dI\), proved by summing geometric series.
The positive gamma product is
\(\prod_{r=1}^{d-1}\Gamma(r/d)=(2\pi)^{(d-1)/2}/\sqrt d\),
by the gamma multiplication formula. The column powers in (IPM.20)
multiply to \(d^{-m/2}e^{\pi i m/2}\). Combining these quantities
gives (IPM.21) at zero; (IPM.19) proves its entire dependence on \(a\).
In particular the factor at \(m=1\) is
\(-i\sqrt{2\pi}e^{-a^2/2}\), with the prescribed orientation.

For \(m=2\), let \(\zeta=e^{2\pi i/3}\). The Airy contour
definition on the oriented rays from angle \(-\pi/3\) to \(\pi/3\)
is \(2\pi i\operatorname{Ai}(a)\). The negative of this contour
is \(\Gamma_2\); hence its integral is
\(-2\pi i\operatorname{Ai}(a)\). Substitution \(x=\zeta y\)
carries that Airy contour to \(\Gamma_1\), with factor \(\zeta\)
from the differential and argument \(\zeta a\). Derivatives with
respect to \(a\) supply the negative second columns. Consequently
\[
 \begin{gathered}
 H_2(a)=2\pi i
 \begin{pmatrix}
  \zeta\operatorname{Ai}(\zeta a)&-\zeta^2\operatorname{Ai}'(\zeta a)\\
  -\operatorname{Ai}(a)&\operatorname{Ai}'(a)
 \end{pmatrix},\qquad \det H_2(a)=2\pi i.
 \end{gathered}
 \tag{IPM.22}
\]
This verifies the pasted Airy formula and its determinant directly on
the stated contours, rather than choosing their signs from a determinant.

\subsection{Full constants in the rectangular cyclic estimate}

For \(|a|\le A\) retain
\[
 \begin{gathered}
 C_m(A)=\frac m{m+1}2^{1/m}A^{(m+1)/m},\quad
 B_b(A)=\frac{2e^{C_m(A)}}{m+1}
       \bigl(2(m+1)\bigr)^{(b+1)/(m+1)}
                         \Gamma\bigl((b+1)/(m+1)\bigr),\\
 U_m(A)^2=m\sum_{b=0}^{m-1}B_b(A)^2,\quad
 D_m(A)=\begin{cases}(2\pi)^{m/2},&m\ge2,\\
                 \sqrt{2\pi}e^{-A^2/2},&m=1,\end{cases}\quad
 L_m(A)=D_m(A)/U_m(A)^{m-1}.
 \end{gathered}
 \tag{IPM.23}
\]
The maximum of \(Ar-r^{m+1}/(2(m+1))\) occurs at
\(r=(2A)^{1/m}\) and is \(C_m(A)\), including \(A=0\).
Thus the absolute value of either radial integrand is bounded by
\(e^{C_m(A)}r^be^{-r^{m+1}/(2(m+1))}\). Integration of both rays
gives \(B_b(A)\). The Hilbert--Schmidt norm is at most \(U_m(A)\).
The determinant absolute value is at least \(D_m(A)\), by (IPM.21)
and \(\Re(a^2)\le|a|^2\). The product of the largest \(m-1\)
singular values is at most \(U_m(A)^{m-1}\); hence
\[
 L_m(A)^2I\preceq H_m(a)^*H_m(a)\preceq U_m(A)^2I.
 \tag{IPM.24}
\]
Congruence by \(D_vC_\rho\) and multiplication by \(|e^F|^2\)
give the exact full period bounds from the intake, with both factors
and the entire translation matrix still present.

For \(k\) identical complete primary factors put
\(q_\rho=k(m-1)+1\), \(X=S-k\rho\), and
\[
 \begin{gathered}
 \iota_X X^l=
  \sum_{\substack{0\le b_i<m\\|\mathbf b|=l}}
            \frac{l!}{\prod_i b_i!}\,z^{\mathbf b},
       \quad0\le l<q_\rho,\qquad
 d_l=(l!)^2\sum_{\substack{0\le b_i<m\\|\mathbf b|=l}}
                           \prod_i(b_i!)^{-2},\\
 D_{k,v}=\operatorname{diag}(v^{k+l})_{l=0}^{q_\rho-1},\quad
 (D_v^{\otimes k})\iota_X=\iota_XD_{k,v},\quad
 \iota_X^*\iota_X=\operatorname{diag}(d_l).
 \end{gathered}
 \tag{IPM.25}
\]
The multinomial theorem gives every coefficient. Each degree
\(0\le l\le k(m-1)\) has at least one admissible tuple, so \(d_l>0\).
Different total degrees have disjoint monomial supports; their vectors
are independent and orthogonal in this specified coordinate metric.
These facts prove injectivity and both matrix identities. They neither
discard nor quotient the tensor directions complementary to the image.
Tensoring (IPM.24) and restricting by this actual inclusion gives
\[
 L_m(A)^{2k}\operatorname{diag}(d_l)
 \preceq\iota_X^*\left(\bigotimes_iH_m(a_i)^*H_m(a_i)\right)\iota_X
 \preceq U_m(A)^{2k}\operatorname{diag}(d_l).
 \tag{IPM.26}
\]
Translation from \(S\) to \(X\), congruence by \(D_{k,v}\), and
the scalar \(|\exp(k\kappa/v^d-\rho\sum_i a_i/v)|^2\) restore the
original rectangular period Gram. The dimension \(q_\rho\) belongs
to this retained primary image; it is not substituted for the full
quartet dimension \(q=[1+k(m-1)](k+1)^2\) of PGS.

\subsection{Actual local coefficients, incomplete gamma and remainder}

For a root of the full packet, use the original germ
\(s=\rho+\zeta_\rho(y)\) with
\(\Phi_h(s)=\kappa_\rho+y^d/d\), \(d=m_\rho+1\).
The entire theta source gives the literal local analytic amplitude
\[
 a_P(y)=v_h(\rho+\zeta_\rho(y))
        P(\rho+\zeta_\rho(y))\zeta_\rho'(y).
 \tag{IPM.27}
\]
At \(t\ne0\), multiply this amplitude by
\(\exp[-t(\rho+\zeta_\rho(y))/u]\); its supremum in the following
bounds must then be that of this actual parameter-dependent amplitude.
The unmodified MB68 bound is the marked \(t=0\) bound.

Let \(0<R_1<R_0\) lie strictly inside the coordinate disc, and let
\(M_{P,R_0}=\max_{|y|=R_0}|a_P(y)|\). Cauchy's coefficient bound is
\(|[y^n]a_P|\le M_{P,R_0}R_0^{-n}\). The formal reduction coefficient
of \(u^\ell[y^bdy]\), \(0\le b<m_\rho\), is
\[
 [y^{b+\ell d}]a_P\,(-1)^\ell
                      \prod_{j=0}^{\ell-1}(b+1+jd),
\]
whose absolute value is at most
\[
 M_{P,R_0}R_0^{-b}
       \left(\frac d{R_0^d}\right)^\ell\ell!.
 \tag{IPM.28}
\]
Indeed \(0<(b+1)/d<1\), so each factor \(b+1+jd\le d(j+1)\).
This is the complete Gevrey bound on the specified amplitude, not a
convergence statement for its formal reduction.

On a ray \(y=re^{i\theta}\) with
\(\Re(e^{id\theta}/u)\le-\delta/|u|\), put
\(\lambda_\theta=-e^{id\theta}/(du)\), so \(\Re\lambda_\theta>0\).
Using its branch with \(|\arg\lambda_\theta|<\pi/2\), each retained
truncated monomial has the exact integral
\[
 \int_0^{R_1e^{i\theta}}e^{y^d/(du)}y^n\,dy
 =\frac{e^{i(n+1)\theta}}d
       \lambda_\theta^{-(n+1)/d}
       \gamma\left(\frac{n+1}d,\lambda_\theta R_1^d\right).
 \tag{IPM.29}
\]
This follows by \(y=re^{i\theta}\), then \(z=\lambda_\theta r^d\);
the lower incomplete gamma follows that radial segment on the declared
branch. Multiplication by \(e^{\kappa_\rho/u}\) and the signs of the
two chains gives their actual truncated contribution. The Taylor tail
beginning at degree \(n\) is bounded on \(|y|\le R_1\) by
\(M_{P,R_0}|y|^n/((1-R_1/R_0)R_0^n)\). Integration of two such
segments, followed only for this upper bound by extension to infinity,
proves
\[
 |\mathrm{error}|
 \le\frac{2M_{P,R_0}|e^{\kappa_\rho/u}|}
       {(1-R_1/R_0)R_0^n d}
       \left(\frac{d|u|}\delta\right)^{(n+1)/d}
                       \Gamma\left(\frac{n+1}d\right).
 \tag{IPM.30}
\]
Thus MB68 is verified on its marked local domain. Finite endpoints are
the actual relative-chain boundary; their transport to the global cycles
uses the original SPG sectorial maps and keeps their signed incidence
matrices. Nothing in (IPM.29)--30 replaces a truncated monomial by its
integral over an infinite ray.

\section{The boundary-coordinate transport on the actual graph source}
\label{sec:ipm-graph-transport}

Use the full original cyclic \(E=\mathbb C[S]/\chi\), \(q=\deg\chi\),
\(\mathcal P_N\), \(N\ge q-1\), original remainder \(J_N\),
section \(s_N\), relation map \(B_N:Q\mapsto\chi Q\), and every
form of PGS, PGD, PGM, AEX and CMI. In particular
\[
 \begin{gathered}
 U_N=(\mathcal V_N,\mathsf T_k\mathcal V_N),\quad Y=(0,L),\quad
 W_N=U_N^*U_N,\quad Z_N=U_N^*Y,\quad K=Y^*Y,\\
 \widehat M_N=(U_N-YJ_N)^*(U_N-YJ_N)
   =W_N-Z_NJ_N-J_N^*Z_N^*+J_N^*KJ_N.
 \end{gathered}
 \tag{IPM.31}
\]
These are the actual full-unit source maps of PGS.2--9;
\(Z_N\) is their original cross pairing and \(K\) their phase-source
Gram. They have not been replaced by a coordinate period metric.

For any specified invertible \(U:E\to E\), including the retained
boundary map \(D_{k,v}^{-1}\) conjugated by the exact translation,
put
\[
 F_N=I-s_NJ_N+s_NUJ_N,\qquad
 F_N^{-1}=I-s_NJ_N+s_NU^{-1}J_N.
 \tag{IPM.32}
\]
Monic division gives the direct decomposition
\(\mathcal P_N=s_NE\oplus B_N\mathcal P_{N-q}\). On it these
maps are \(U\oplus I\) and \(U^{-1}\oplus I\). This proves both
inverse composites, \(J_NF_N=UJ_N\), and \(F_NB_N=B_N\).
The remainder section is the same polynomial at all cutoffs, so the
maps commute with every admitted degree inclusion. In particular one
common \(F_{2q}\) transports all four original windows.

The transported graph itself, not only its determinant, is
\[
 \begin{gathered}
 U_N'=U_NF_N,\quad Y'=YU,\quad
 W_N'=F_N^*W_NF_N,\quad Z_N'=F_N^*Z_NU,\quad K'=U^*KU,\\
 U_N'-Y'J_N=(U_N-YJ_N)F_N,\quad
 \widehat M_N'=F_N^*\widehat M_NF_N,
 \quad (B_N)^*W_N'B_N=(B_N)^*W_NB_N.
 \end{gathered}
 \tag{IPM.33}
\]
The middle identity uses exactly \(J_NF_N=UJ_N\); expanding its
Gram proves every cross term, including the two negative signs.
It also proves the unchanged relation Gram rather than identifying a
relation with zero before taking its source norm.

For any of the positive forms \(M_N\), let \(R_N\) and \(G_N\)
be its original minimum section and quotient Gram. At fixed quotient
coordinate \(x\), the map \(p\mapsto F_Np\) carries \(J_Np=x\)
bijectively to \(J_N(F_Np)=Ux\), and is an isometry from
\(F_N^*M_NF_N\) to \(M_N\). Uniqueness of the positive minimum gives
\[
 R_N'=F_N^{-1}R_NU,\qquad G_N'=U^*G_NU.
 \tag{IPM.34}
\]
This proof applies also to the full graph form in (IPM.31). On a common
affine form path \(M_x\), with \(U\) fixed in \(x\), all its
source, relation and quotient projectors conjugate by \(F_{2q}^{-1}\).
Its relative derivative and endpoint operator satisfy
\[
 (M_x')^{-1}\dot M_x'=F_{2q}^{-1}M_x^{-1}\dot M_xF_{2q},\qquad
 (M_0')^{-1}M_1'=F_{2q}^{-1}M_0^{-1}M_1F_{2q}.
 \tag{IPM.35}
\]
Hence cyclic trace preserves the signed integrand of CMI.8 and RMT7,
similarity preserves its full relative spectrum, and the isometry
preserves all angle, overlap, trace-norm and Hilbert--Schmidt bounds.
The Neumann polynomials, their centred residuals and the signed radii
of CMI.14--16 conjugate by the same map term by term. This proves their
transport on their exact original source, including the possible zero
of the graph four-window scalar treated in AEX.

Writing \(V_N=\det G_N\), the common factor is
\(V_N'=|\det U|^2V_N\). For the identical-primary boundary map,
\[
 \log|\det U|^2
     =-2\left(kq_\rho+\frac{q_\rho(q_\rho-1)}2\right)\log|v|.
 \tag{IPM.36}
\]
Its coefficient is the sum of the original exponents \(k+l\).
Both numerator and denominator of the four-window expression have
two factors, so this term cancels exactly. A common scalar exponential
line contributes \(q\log|e^F|^2\) to each determinant and cancels by
the same four signs. These are exact evaluations of the transported
cost. No sign or size estimate from (IPM.23)--26 makes that zero
contribution into an arithmetic gain.

\subsection{Exact Gram effect of the additional entire-amplitude map}

On any fixed original \(\mathcal P_N\), assemble the global entire
observation in (IPM.14) into a matrix \(P_N\). Let \(P_N^{[L]}\)
be a specified finite reduction or local approximation whose remainder
matrix \(E_N=P_N-P_N^{[L]}\) includes the global tails and every
local endpoint. Matrix multiplication, with no discarded cross term,
gives
\[
 P_N^*P_N-(P_N^{[L]})^*P_N^{[L]}
   =(P_N^{[L]})^*E_N+E_N^*P_N^{[L]}+E_N^*E_N.
 \tag{IPM.37}
\]
This is the period Gram discrepancy on the same source. It is not an
assertion of equality with \(\widehat M_N\). Its exact positive
source comparison is
\(\widehat M_N^{-1}P_N^*P_N\), self-adjoint in the original
\(\widehat M_N\) metric because pairing with that form gives
\(\langle P_Np,P_Nq\rangle\).
Here positive means positive semidefinite: its kernel is exactly
\(\ker P_N\), since its quadratic value is \(\|P_Np\|^2\).
No injectivity of this entire-amplitude map, and hence no positive
definiteness of this comparison, is assumed. Its exact adjoint has
type
\[
 P_N^\dagger=\widehat M_N^{-1}P_N^*:
       \mathbb C^{d_h^k}\longrightarrow\mathcal P_N.
 \tag{IPM.37a}
\]
Pairing on \(\mathcal P_N\) uses the unchanged graph form and pairing
on the target uses its stated coordinate form. For the receiving
quotient (IPM.15c), let \(p_{\perp,N}\) be Euclidean orthogonal
projection onto \((\mathscr R_{N,k}^{\rm per})^\perp\). The quotient
is identified isometrically with this orthogonal complement by
taking the unique perpendicular representative. Its actual matrix
is \(p_{\perp,N}O_Ns_N\), with kernel exactly (IPM.12d)'s tensor
version, and its Gram is
\(s_N^*O_N^*p_{\perp,N}O_Ns_N\). It is positive definite only on
the resulting quotient by that exact kernel, where its quadratic
value is the squared norm of a nonzero receiving image.

For the explicit column errors \(\|(E_N)_j\|\le e_j\) constructed in ECB1--5 below,
the actual source error satisfies
\[
 \begin{gathered}
 \|E_N\widehat M_N^{-1/2}\|^2
 \le\|E_N\widehat M_N^{-1/2}\|_{\rm HS}^2
 \le\sum_{i,j}|(\widehat M_N^{-1})_{ij}|e_ie_j=:\varepsilon_N^2,\\
 \left\|\widehat M_N^{-1/2}
   \bigl[P_N^*P_N-(P_N^{[L]})^*P_N^{[L]}\bigr]
                        \widehat M_N^{-1/2}\right\|
 \le2\|P_N^{[L]}\widehat M_N^{-1/2}\|\varepsilon_N+\varepsilon_N^2.
 \end{gathered}
 \tag{IPM.38}
\]
The first inequality is the operator/HS inequality; expansion of its
squared trace and \(|\langle E_i,E_j\rangle|\le e_ie_j\) proves
the second. The final bound follows by applying submultiplicativity to
all three exact terms in (IPM.37). Thus local incomplete-gamma estimates
have a precise receiving operator on the original graph norm. The
inverse source Gram, every mass and every full-unit factor remain.

The original PGM mixed residual is transported as
\(\mathscr D_k'=\mathscr D_kU\), with exact Gram
\[
 (\mathscr D_k')^*\mathscr D_k'
      =U^*(\Omega_k^{\rm PGS}-K_{\rm low})U.
 \tag{IPM.39}
\]
This identity is a direct congruence of the original operator and its
receiving weighted Hilbert space; its strict positivity for \(k\ge2\)
survives because \(U\) is invertible. It supplies no uniform lower
bound as \(k\) varies. In particular \(\Omega_k^{\rm PGS}\) here is
the PGS graph residual, not the differently typed scalar correction
denoted \(\Omega_k\) in AAM14.

The remaining calculation on the program's actual source is now the
signed contribution of (IPM.11)--14 and (IPM.37), transported by the
already constructed SPG sectorial maps on their marked \(t=0\)
domain, to the CMI three-path scalar
at the four original cutoffs. All its receiving maps and integrals are
defined above. Its comparison must use the exact AAM14 threshold and
the full quartet value of \(q\); a bound for one primary rectangular
coordinate form does not evaluate those original cross pairings.
The present proof computes the additional map and its domain, verifies
the delivered local estimates, and proves the complete cancellation
claimed in MB70--71. It does not claim an RH contradiction or an
actual offcritical zero.

\paragraph{Primary formula records.}
The gamma multiplication identity used in (IPM.21) is
\href{https://dlmf.nist.gov/5.5.E6}{DLMF 5.5.6}; the Airy contour
definition used in (IPM.22) is
\href{https://dlmf.nist.gov/9.5.E4}{DLMF 9.5.4}. The identification of
the original \(g=2\xi\) with the theta expression (IPM.2) is also
recorded by \href{https://dlmf.nist.gov/25.4.E4}{DLMF 25.4.4} and
\href{https://dlmf.nist.gov/25.5.E13}{DLMF 25.5.13--14}; its needed
growth and convergence estimates were proved explicitly above.

% END COMPLETE CORRECTED SOURCE period_metric.tex


% BEGIN COMPLETE CORRECTED SOURCE explicit_column_bounds.tex
\subsection{A constructed global approximation and all its column errors}

Here is a concrete choice of the finite approximation used in IPM.37--38.
Put \(n=d_h\), choose \(R_0>R_1\ge\max(1,2R_h)\), and let
\(a_b(s)=v_h(s)s^b\), for \(0\le b\le N\).
Retain the actual supremum \(M_b=\max_{|s|=R_0}|a_b(s)|\);
it is at most \(2^n R_0^{b-n}\mathcal X(R_0)\) by IPM.4.
For each parameter \(t_i\) define
\[
 E_i(r)=\exp\left[
 -\frac{r^{n+1}}{(n+1)|u|}
 +\frac{\sum_{\nu=0}^n|c_{i,\nu}|r^\nu}{|u|}\right],
 \qquad
 f_{t_i}(s)=\frac{s^{n+1}}{n+1}
                    +\sum_{\nu=0}^nc_{i,\nu}s^\nu .
 \tag{ECB1}
\]
This majorizes the absolute exponential on every original ray,
including on \(0\le r\le R_1\); no coefficient of the phase is
changed. Approximate the one-factor integral of \(a_b\) by integrating
its Taylor polynomial through degree \(L-1\) on the two original ray
segments of radius \(R_1\), with their original orientations.
Call this vector \(Q_{i,b}^{[L]}\), and its exact whole-ray vector
\(Q_{i,b}\). Define the following explicit finite numbers:
\[
\begin{split}
 e_{i,b}^{\rm entry}
 &=\frac{2M_b}{(1-R_1/R_0)R_0^L}
                      \int_0^{R_1} E_i(r)r^L\,dr\\
 &\quad+2^{n+1}\int_{R_1}^\infty
                      E_i(r)r^{b-n}\mathcal X(r)\,dr,\\
 C_{i,b}
 &=2M_b\int_0^{R_1}\frac{E_i(r)}{1-r/R_0}\,dr
       +2^{n+1}\int_{R_1}^\infty
                      E_i(r)r^{b-n}\mathcal X(r)\,dr .
\end{split}\tag{ECB2}
\]
Cauchy's coefficient estimate bounds the omitted Taylor sum by
\(M_b(r/R_0)^L/(1-R_1/R_0)\).
Integration gives the first term in the error; IPM.4 gives the
two-ray exterior term. The same geometric series bounds the whole
amplitude on the compact segments and proves the bound \(C_{i,b}\).
Convergence of the exterior terms was proved in IPM.3--6.
Thus every entry of \(Q_{i,b}-Q_{i,b}^{[L]}\) is bounded by
\(e_{i,b}^{\rm entry}\), every entry of \(Q_{i,b}\) by \(C_{i,b}\),
and every entry of \(Q_{i,b}^{[L]}\) by
\(C_{i,b}+e_{i,b}^{\rm entry}\).

For the original source column \(S^j=(\sum_i s_i)^j\), define its
approximation by the complete multinomial expansion
\[
 (P_N^{[L]})_j
 =\sum_{\substack{b_1+\cdots+b_k=j\\b_i\ge0}}
       \frac{j!}{\prod_i b_i!}
                   \bigotimes_i Q_{i,b_i}^{[L]} .
 \tag{ECB3}
\]
This uses unrestricted polynomial degrees, before any finite-jet
quotient. The exact column has the same formula with \(Q_{i,b_i}\).
The telescoping product identity therefore proves the actual
Euclidean column bound
\[
\begin{split}
 \|(E_N)_j\|&\le e_j,\\
 e_j&=n^{k/2}
 \sum_{\substack{b_1+\cdots+b_k=j\\b_i\ge0}}
 \frac{j!}{\prod_i b_i!}
 \sum_{i=1}^k e_{i,b_i}^{\rm entry}
       \prod_{r<i} C_{r,b_r}
       \prod_{r>i}(C_{r,b_r}+e_{r,b_r}^{\rm entry}) .
\end{split}\tag{ECB4}
\]
Indeed every product-contour entry has that scalar bound, and the
target has exactly \(n^k\) entries. This constructs the finite
approximation and all its error constants without a presumed
arithmetic interval certificate. Its coefficients are the original
Taylor coefficients of \(v_h(s)s^b\); the integrals still have the
original polynomial phase and declared contours.
Local-coordinate approximations can instead retain IPM.29--30
and their explicit endpoint terms; they are not substituted for the
just constructed global approximation.

Using these particular \(e_j\) in IPM.38 now gives its concrete
original-source estimate
\[
 \varepsilon_N^2
 =\sum_{i,j}|(\widehat M_N^{-1})_{ij}|e_i e_j,\qquad
 \|E_N\widehat M_N^{-1/2}\|\le\varepsilon_N .
 \tag{ECB5}
\]
For clarity, its proof is expansion of
\(\operatorname{Tr}(E_N^*E_N\widehat M_N^{-1})\):
each coefficient has absolute value at most \(e_i e_j\), by
Cauchy--Schwarz. The trace bounds the square of the operator norm,
giving ECB5. Substitution into the three terms of IPM.37 proves
the bound there on the full Gram discrepancy.
Every inverse source Gram entry is present. No uniform estimate
as \(k\), \(q_k\), or the cutoff grows follows by omitting these
entries or by treating the fixed finite integrals as constants
independent of their displayed parameters.

% END COMPLETE CORRECTED SOURCE explicit_column_bounds.tex


% BEGIN COMPLETE CORRECTED SOURCE mixed_attachment.tex
\section{Attachment to the original leg complex and every coefficient face}
\label{sec:intake-mixed-attachment}

The following construction specifies the coefficient and source arrows
used to attach the reconstructed boundary calculation to the defining
total object.  It retains the original outer labels and is a diagram of
complexes, rather than a colimit identifying those labels.
All tensor products in this paragraph are over \(\mathbb C\).
Let \(W\subset V\) be one of the original admitted source subspaces,
\(\widehat W=\mathcal F W\), where \(\mathcal F\) is the original
Fourier transform.  On even Schwartz functions
\(\mathcal F^2=I\), with the original \(2\pi\) convention.
Use the same target \(\mathscr B_\Omega\) as the original source
complex.  In degrees zero and one put
\[
 C_W=[W\xrightarrow{\Theta}\mathscr B_\Omega],\qquad
 T_W=[W\oplus\widehat W
              \xrightarrow{d_T}\mathscr B_\Omega],\qquad
 d_T(\phi,\psi)=\Theta(\phi-\mathcal F\psi).
 \tag{IMS1}
\]
These are the admitted positive and joint source presentations,
with the target restriction inherited from the original diagram.
Define actual cochain maps \(p:T_W\to C_W\) and \(i:C_W\to T_W\) by
\[
 \begin{gathered}
 p^0(\phi,\psi)=\phi-\mathcal F\psi,\quad p^1=I,\\
 i^0(\phi)=(\phi/2,-\mathcal F\phi/2),\quad i^1=I,\qquad
 d_Ti^0=\Theta,\quad \Theta p^0=d_T,\quad pi=I .
 \end{gathered}\tag{IMS2}
\]
All equalities follow by substitution and \(\mathcal F^2=I\).
In particular \(p^0\) lands in \(W\), because
\(\mathcal F\widehat W=W\).  The complementary projection is
\[
 (I-ip)^0(\phi,\psi)
   =\left(\frac{\phi+\mathcal F\psi}{2},
           \mathcal F\frac{\phi+\mathcal F\psi}{2}\right),
 \qquad (I-ip)^1=0 .
 \tag{IMS3}
\]
Its image is the actual complex
\([\,\{(w,\mathcal Fw):w\in W\}\to0\,]\).
The maps \(i^0(\phi)\) and \((w,\mathcal Fw)\) have intersection zero:
applying \(p^0\) gives \(\phi=0\), and then \(w=0\).
Every vector is their sum by (IMS3), and the differential on the
second summand is zero.  This proves an isomorphism of complexes
\(T_W\simeq C_W\oplus[W\to0]\), with the displayed maps and no
removal of its degree-zero diagonal cycle.

The two one-leg lifts to the same joint target are
\[
 i_+^0(\phi)=(\phi,0),\qquad
 i_-^0(\phi)=(0,-\mathcal F\phi),\qquad i_\pm^1=I .
 \tag{IMS4}
\]
Here \(i_-\) includes the sign that identifies the negative-leg
complex with \(C_W\); the original negative-leg coefficient is
\(-\mathcal F\phi\).  Direct calculation gives
\[
 i_+^0-i^0=(\phi/2,\mathcal F\phi/2),\qquad
 i_-^0-i^0=(-\phi/2,-\mathcal F\phi/2).
 \tag{IMS5}
\]
Thus the two lifts differ by specified diagonal cycles.  No
degree-minus-one homotopy is claimed to kill those cycles.
An inclusion \(W\subset W'\) commutes with all these maps, because
its Fourier image is the corresponding inclusion
\(\widehat W\subset\widehat W'\).

Let \(L^\bullet\to M^\bullet\) be any of the actual coefficient
cochain maps in the reconstructed calculation: the polynomial packet
map, the two free terms of a ramified lattice map, or their specified
ordered tensor map.  Use its displayed free complex when calculating
a derived specialization.  On the total tensor complex the
differential is, for homogeneous \(c\in C^r\),
\[
 d(c\otimes l)=d_Cc\otimes l+(-1)^r c\otimes d_Ll .
 \tag{IMS6}
\]
The maps \(I_C\otimes f\), \(p\otimes I\), \(i\otimes I\) and the two
leg inclusions commute with this differential.  For example,
\(d(I_C\otimes f)(c\otimes l)
=d_Cc\otimes f(l)+(-1)^r c\otimes d_Mf(l)\);
using \(d_Mf=fd_L\) gives \((I_C\otimes f)d(c\otimes l)\).
The other identities follow from (IMS2) and exactly the same
calculation.  They also commute with one another on each simple
tensor and hence on every finite sum.  This proves the tensor
attachment and its signs on the complexes themselves.  It does not
assert that a non-\(\mathbb C[s]\)-linear analytic section is balanced
over that ring; the original strict free resolutions supply any
additional balancing in the defining object.

For each independent coefficient face \(A\) at an original outer
label, form the direct sum of its labelled copies of these complexes.
For \(A\subset B\) let \(z_{AB}\) insert zero in the newly admitted
slots.  For every displayed cochain map \(f\),
\[
 f_B z_{AB}=z_{AB}f_A,\qquad
 \ker f_A=\bigoplus_{\alpha\in A}\ker f,\qquad
 \operatorname{im}f_A=\bigoplus_{\alpha\in A}\operatorname{im}f .
 \tag{IMS7}
\]
The first identity evaluates to \(f(c_\alpha)\) in an old slot and
zero in a new slot on both sides.  A vector is in the kernel exactly
when each of its coordinates is; choosing each coordinate preimage
proves the image equality.  The same proof applies to the explicit
derived kernels and ordinary quotients computed from the free
ramified resolution. The identical direct-sum calculation applies in
the separate \(\ell\)-adic coefficient category to the actual
varying-parameter restriction map of the finite-field chapter;
no \(\mathbb C\)-linear map into that category is asserted.
Where that restriction
map has zero source and a nonzero target, its cokernel is retained in
each labelled slot; it is not changed into a fibrewise isomorphism.

Lifting a coefficient map to the original supported carrier sends
\(x^\bullet\) to \(f(x)^\bullet\) and \(\tau\) to \(\tau\).
Consequently a represented zero remains the supported zero of its
receiving fibre, whereas an absent input remains absent.  The empty
coefficient face is the zero complex at its unchanged present outer
label.  An original synchronization map changing the outer or leg
label still acts by its original specified zero insertions; (IMS7)
proves its commutation with each coefficient map on every surviving
slot.  No assertion that an outer support map reflects bottom is
needed.

Finally, if \(C\subset C'\) is an original degreewise injective
two-term source transition with identical degree-one target \(B\),
the actual kernel in degree-one cohomology is
\[
 \ker\bigl(B/dC^0\longrightarrow B/d'C'^0\bigr)
            =d'C'^0/dC^0 .
 \tag{IMS8}
\]
Indeed a class represented by \(b\in B\) maps to zero precisely
when \(b\in d'C'^0\), and changing its representative adds
\(dC^0\).  This is a cohomology kernel, while the degreewise kernel
of an injective cochain transition is zero.  If the degree-one
target is itself restricted, the same proof replaces the numerator
by its intersection with that target.  The comparison therefore
retains exactly the original proper-source kernel, including on
every independent face.

The global finite-jet period observation, the whole-amplitude source
map and its explicit discrepancy in the analytic chapter are
additional observations on these retained presentations.  Their
agreement or factorization is used only where that chapter proves
the corresponding map.  In particular the relation image is kept
before passing to any of its specified receiving quotients.

% END COMPLETE CORRECTED SOURCE mixed_attachment.tex


% BEGIN COMPLETE CORRECTED SOURCE coordinating_next_wave/source_snapshot/accepted/holonomy_hcd_complete.tex
\begingroup
\def\C{\mathbb C}
\def\R{\mathbb R}
\def\av{\operatorname{av}}
\def\im{\operatorname{im}}
\section{Continuous holonomy cochains and the original relation metric}



This continuation proves the continuous-phase version of H29--H31
and the exact maps to the completed targets in H64 of the delivered
\emph{All-holonomy descent of the original theta quotient}.
The delivered NOTE.tex was read completely, all 1152 lines;
its SHA-256 is
\begin{center}\small
\texttt{444f82a8fa420e6f471ef448442085f3a93acb137069e1505d155faab051b881}.
\end{center}
The archive SHA-256 is
\begin{center}\small
\texttt{f0506a27cd4c9012f43fa70c3d1f8f996bbb6db01b0eea2944b34dcf290ca8f0}.
\end{center}
All source coordinates, masses, phases, relation columns, and full
primary multiplicities below are retained. No finite test or Lean
execution is used as a proof.
The sampled-density argument uses the full analytic proof
\emph{Periodized source density review}, PSD8--PSD14,
whose TeX SHA-256 is
\begin{center}\small
\texttt{51102189bcc7a45610464143a5c0bee21c08059c317024e6e856bc7de993fe74}.
\end{center}
The completed-fibre calculation develops H64 and is also proved
in the companion completed-quotient and holonomy-intake derivations.
Its repetition here supplies the maps needed for this continuous
common-image complex, including the mean-zero contraction and
the precise weighted relation gap.

\subsection{Original coefficient spaces and actual phase norms}

Fix the original monic polynomial $\chi$, its full degree $q\ge1$,
and an integer $N\ge q-1$. Put
\[
 \begin{gathered}
 E=\C[S]/(\chi),\quad P_N=\C[S]_{\le N},\quad
 J=J_N:P_N\longrightarrow E,\quad\\
 B_N=\ker J=\chi P_{N-q},\quad A=M_S:E\longrightarrow E,
 \end{gathered}
 \tag{HCD1}
\]
where $P_{-1}=0$. Use the original increasing-power coefficient
bases. Let $B:\C^r\to P_N$, $r=N-q+1$, be multiplication by
the literal $\chi$, so that $\im B=B_N$. At $r=0$ it is the
unique map from the zero coefficient space.

Let $\Theta=[0,2\pi)$ carry its actual Fourier phase measure
$d\mu(\theta)=d\theta/(2\pi)$. Retain the measurable Hermitian
phase source matrices $M_\theta$ and their mean
\[
 M=\int_\Theta M_\theta\,d\mu,\qquad
 0<a I_{P_N}\preceq M_\theta\preceq b I_{P_N}<\infty
 \quad\text{almost everywhere},\qquad a,b>0.
 \tag{HCD2}
\]
The original source isometry H13--H17 identifies $M$ with the
original source Gram, including its full constant mass.
For the actual source, the bounds in (HCD2) are obtained from
H39--H41, also proved in HG3, as follows. Let
$\Psi_N:P_N\to L^2(\R;\mathcal K)$ be the original finite
source columns after the half-density map, with
$\mathcal K=L^2(\R^{k-1},dz)$. Fix $\beta>0$ for the retained
weighted source Grams, whose finite entries are supplied by
the original source regularity:
\[
 M_{\beta,\pm}=(e^{\pm\beta r}\Psi_N)^*
                    (e^{\pm\beta r}\Psi_N),\qquad
 \kappa_{\beta,N}=
 \max_{v\ne0}\frac{v^*(M_{\beta,+}+M_{\beta,-})v}{v^*Mv}.
 \tag{HCD2a}
\]
These are the original weights $x_k^{\pm2\beta}$; all original
source measures remain present. The maximum is finite and positive
because $M>0$ and the two weighted Grams are finite positive forms.
For $\mathcal C(s)=\int\Psi_N(r)^*\Psi_N(r+s)\,dr$ and $s\ge0$,
write the integrand at $v$ as the inner product of
$e^{-\beta r}\Psi_N(r)v$ and
$e^{\beta(r+s)}\Psi_N(r+s)v$, multiplied by $e^{-\beta s}$.
Cauchy--Schwarz, followed by translation in the second integral,
gives
\[
 |v^*\mathcal C(s)v|
 \le e^{-\beta s}
 \sqrt{(v^*M_{\beta,-}v)(v^*M_{\beta,+}v)} .
\]
For $-s$ the same bound follows from
$\mathcal C(-s)=\mathcal C(s)^*$.
Unfolding the literal Zak sum gives
$M_\theta=\sum_{j\in\mathbb Z}e^{ij\theta}\mathcal C(jL)$.
The correlation bound and polarization give absolute convergence
of each matrix entry, uniformly in phase. Pairing $j$ and $-j$,
using $2\sqrt{uv}\le u+v$, and summing the geometric series prove
\[
 -\delta_LM\preceq M_\theta-M\preceq\delta_LM,\qquad
 \delta_L=\frac{\kappa_{\beta,N}}{e^{\beta L}-1}.
 \tag{HCD2b}
\]
For any chosen $0<\eta<1$, the explicit finite circumference
\[
 L\ge\frac1\beta\log(1+\kappa_{\beta,N}/\eta)
 \quad\Longrightarrow\quad
 \delta_L\le\eta<1
 \tag{HCD2c}
\]
therefore supplies (HCD2), with
$a=(1-\delta_L)\lambda_{\min}(M)$ and
$b=(1+\delta_L)\lambda_{\max}(M)$ in the original coefficient
frame. The phase mean is still exactly $M$: every nonzero
Fourier mode integrates to zero against $d\theta/(2\pi)$.
The concrete choice $\eta=1/2$ gives
$L=\beta^{-1}\log(1+2\kappa_{\beta,N})$ and $\delta_L=1/2$.
These are actual finite-source bounds at the stated cutoff;
no bound uniform in $N$ or the arithmetic packet is asserted.
The argument below also applies to any fixed $L$ for which
(HCD2) is already established.

Define the Hilbert space
\[
 \mathscr P=L^2(\Theta;P_N,M_\theta\,d\mu),\qquad
 \|P\|_{\mathscr P}^2=\int_\Theta P(\theta)^*M_\theta
 P(\theta)\,d\mu.
\]
By (HCD2), its elements are exactly coefficientwise $L^2$
functions, with equivalent norms. Equality means equality
almost everywhere. The coefficientwise Bochner integral
$\av P=\int P\,d\mu$ exists, since coefficientwise $L^2$
implies $L^1$ on this finite phase measure. Every fixed
coefficient map commutes with this integral, entry by entry.

The closed relation space and common-image space are
\[
 \mathscr B=L^2(\Theta;B_N),\qquad
 \mathscr P^{\rm eq}
 =\{P\in\mathscr P:JP(\theta)=x\text{ a.e. for a fixed }x\in E\}.
 \tag{HCD3}
\]
They carry the restricted source norms. They are closed:
$J$ is continuous on coefficientwise $L^2$, $B_N=\ker J$
is closed, and constant $E$-valued functions form a closed
subspace. For the last assertion, averaging followed by
repetition is a bounded idempotent with precisely that image.

\subsection{The continuous cochain splitting and its full contraction}

Use complexes in degrees $0,1$ with actual inclusion differentials:
\[
 \mathcal C_N=[B_N\hookrightarrow P_N],\qquad
 \mathscr C_N=[\mathscr B\hookrightarrow\mathscr P^{\rm eq}].
 \tag{HCD4}
\]
Repetition $\Delta v(\theta)=v$ and averaging $\av$ are cochain
maps. In degree one, $J\av P=\int JP\,d\mu=x$; in degree
zero the mean belongs to $B_N$. The equality $\av\Delta=I$
uses the literal phase measure. Repetition is an isometry:
\[
 \|\Delta v\|_{\mathscr P}^2
 =v^*\left(\int M_\theta\,d\mu\right)v=v^*Mv.
\]
Averaging is bounded with the explicit estimate
\[
 \|\av P\|_M^2
 \le b\|\av P\|_{\rm coeff}^2
 \le b\int\|P(\theta)\|_{\rm coeff}^2\,d\mu
 \le (b/a)\|P\|_{\mathscr P}^2.
 \tag{HCD5}
\]
This uses scalar Cauchy--Schwarz for the coefficient vector
and does not replace the arithmetic mass by one.

Put $\mathscr B^0=\{b\in\mathscr B:\av b=0\}$. The inverse
continuous coefficient maps of complexes are
\[
 \begin{gathered}
 \mathscr C_N\simeq
 \mathcal C_N\oplus[\mathscr B^0\xrightarrow{I}\mathscr B^0],\\
 P\longmapsto(\av P,P-\Delta\av P),\qquad
 (v,b^0)\longmapsto\Delta v+b^0.
 \end{gathered}
 \tag{HCD6}
\]
In degree one, $J(P-\Delta\av P)=x-x=0$ and its mean is
zero. In degree zero every value already belongs to $B_N$.
Both composites restore the original vector. Inclusion
carries these components to inclusion and identity,
respectively, proving the isomorphism of complexes.

The degree-minus-one map
\[
 h:\mathscr P^{\rm eq}\longrightarrow\mathscr B,\qquad
 hP=P-\Delta\av P
 \tag{HCD7}
\]
satisfies $dh+hd=I-\Delta\av$. In degree one this is
$dhP=P-\Delta\av P$; in degree zero it is
$hd b=b-\Delta\av b$. On the complementary complex it
is the identity from degree one to degree zero. Boundedness
follows from (HCD5) and the isometry of repetition.
It is not asserted orthogonal or isometric.

The common quotient map is
\[
 \mathfrak j:\mathscr P^{\rm eq}\longrightarrow E,\qquad
 P\longmapsto JP(\theta)\text{ a.e.}
 \tag{HCD8}
\]
It is onto by repeating a right inverse of $J$, and its
kernel is exactly $\mathscr B$. Inclusion is injective.
Therefore $H^0(\mathscr C_N)=0$ and $H^1(\mathscr C_N)=E$,
with the isomorphism induced by the original full $J$.
Repetition and averaging induce the identity on this $E$.

Multiplication by $S$ maps $P_N$ to $P_{N+1}$ and $B_N$
to $B_{N+1}$. Applied pointwise, it sends common image $x$
to common image $Ax$ and commutes with repetition, mean
and $h$. These are maps between the stated cutoffs.
On degree-one cohomology the action is the full companion
$A$, including its nilpotent primary parts. No trace of
the infinite-dimensional contractible complement is assumed.

\subsection{The minimum norm of the common image}

Retain the original and phasewise minimum sections
\[
 \begin{aligned}
 G&=(JM^{-1}J^*)^{-1},& C&=M^{-1}J^*G,\\
 G_\theta&=(JM_\theta^{-1}J^*)^{-1},&
 C_\theta&=M_\theta^{-1}J^*G_\theta,\qquad\\
 &&\overline G&=\int G_\theta\,d\mu.
 \end{aligned}
 \tag{HCD9}
\]
Surjectivity and positivity make every inverse defined.
Continuity of matrix inversion proves measurability.
For any fixed coefficient right inverse $C_{\rm ref}$ of $J$,
\[
 \frac{a}{\|J\|_{\rm coeff}^2}I_E
 \preceq G_\theta
 \preceq b\|C_{\rm ref}\|_{\rm coeff}^2I_E.
 \tag{HCD10}
\]
Indeed any lift $p$ of $x$ has
$\|x\|\le\|J\|\|p\|$, proving the lower bound, and the
particular lift $C_{\rm ref}x$ proves the upper bound.
The formula for $C_\theta$ then also gives boundedness.
All following integrals and sections are consequently defined.

Direct multiplication gives $JC_\theta=I$, and for $b_0\in B_N$,
\[
 b_0^*M_\theta C_\theta=b_0^*J^*G_\theta=0.
 \tag{HCD11}
\]
The section $\mathscr Sx(\theta)=C_\theta x$ has common image
$x$. Any other $P$ with that image has
$P-\mathscr Sx\in\mathscr B$. Phasewise orthogonality gives
\[
 \|P\|_{\mathscr P}^2
 =x^*\overline Gx+\|P-\mathscr Sx\|_{\mathscr P}^2.
 \tag{HCD12}
\]
Thus the quotient metric of the continuous complex is exactly
$\overline G>0$ on its full $E$. For the constant-source
complex the identical argument gives $G$. Repetition is
isometric on sources, and its induced identity
$(E,G)\to(E,\overline G)$ has the additional continuous
relation directions available in the target.

\subsection{The positive relation gap and its mean-zero decomposition}

For $r>0$, put
$X_\theta=(B^*M_\theta B)^{-1}B^*M_\theta C$.
Both sections have right-inverse image $I_E$, so their difference
lies in $B_N$ and equals $BX_\theta$; multiplication by
$B^*M_\theta$ proves the formula and uniqueness.
At $r=0$, both sections equal the unique inverse of $J$ and
$X_\theta$ is the zero map into the zero space.
Expanding $C=C_\theta+BX_\theta$ and using (HCD11) proves
\[
 \boxed{
 \begin{aligned}
 G-\overline G=\mathscr D
 &=\int(C-C_\theta)^*M_\theta(C-C_\theta)\,d\mu\\
 &=\int X_\theta^*(B^*M_\theta B)X_\theta\,d\mu\succeq0.
 \end{aligned}}
 \tag{HCD13}
\]
The integral of $C^*M_\theta C$ is $C^*MC=G$.
Every term is the norm of a field in the original relation
image. Multiplication by the literal $\chi$, followed by
the inherited finite tensor theta primitive, realizes
that field in the original cochain complex.

The relation to the particular mean-zero contraction is also exact.
Define
\[
 \overline C=\int C_\theta\,d\mu,\quad
 a_C=C-\overline C:E\to B_N,\quad
 h_\theta=C_\theta-\overline C:E\to B_N.
 \tag{HCD14}
\]
Then $\int h_\theta\,d\mu=0$, and $h\mathscr S$ is precisely
the field $h_\theta$. The original discrepancy is
$C-C_\theta=a_C-h_\theta$ and need not have mean zero.
Since $a_C$ has relation-valued columns, phase and
common-source orthogonality give
\[
 \begin{split}
 \int a_C^*M_\theta h_\theta\,d\mu
 &=-a_C^*M\overline C\\
 &=-a_C^*M(C-a_C)=a_C^*Ma_C.
 \end{split}
 \tag{HCD15}
\]
The adjoint cross term is the same Hermitian matrix.
Expanding all terms of (HCD13) therefore proves
\[
 \boxed{\mathscr D
 =\int h_\theta^*M_\theta h_\theta\,d\mu-a_C^*Ma_C.}
 \tag{HCD16}
\]
This retains the positive relation norm and the complete
weighted cross terms. A mean-zero coefficient field is
not presumed orthogonal to constant relations for the
varying metric. At $N=q-1$, $B_N=0$,
$C_\theta=C=J^{-1}$, $a_C=h_\theta=0$ and $\overline G=G$.

\subsection{Actual sampled Hilbert spaces, including zero weights}

Fix the original circumference $L>0$, tensor degree $k\ge1$,
and $c=k/2$. Retain
\[
 \begin{gathered}
 u_n(\theta)=(2\pi n+\theta)/L,\qquad
 s_n(\theta)=c+iu_n(\theta),\\
 w_n(\theta)=\frac{2\pi}{L}m_{h,k}(u_n(\theta))\ge0,
 \qquad n\in\mathbb Z.
 \end{gathered}
 \tag{HCD17}
\]
The Fourier coefficient is literally $2\pi/L$, with
$m_{h,k}=w_h^{*k}$, $w_h=|(g/h)(1/2+iu)|^2/(2\pi)$
and $g=2\xi$. Sampling gives
\[
 \mathsf T_\theta P=(P(s_n(\theta)))_n,\qquad
 \|\mathsf T_\theta P\|_\theta^2
 =\sum_n w_n(\theta)|P(s_n(\theta))|^2=P^*M_\theta P.
 \tag{HCD18}
\]
This is the actual phase source Gram H17.

Define $\mathcal H_\theta$ to be weighted square-summable
sequences modulo sequences supported where $w_n(\theta)=0$.
Thus zero weights, possible when $k=1$, have an explicit
null subspace. For $k\ge2$, positivity of the original
convolution makes every sample weight positive.
At $k=1$, retain the exact ideal equality
$(\chi_{h,1})=(h)$. If $a_h\ne0$ is the original leading
coefficient of $h$, then $\chi_{h,1}=h/a_h$ is its monic
generator, and $\chi_{h,1}=h$ in the original monic packet
convention. Indeed a polynomial annihilates multiplication by
$S$ on $\C[S]/(h)$ if and only if its action on $[1]$ is
zero, which is equivalent to belonging to $(h)$.
The analytic source remains the original $g/h$, with
the factor $a_h$ retained in that quotient and its norm.
Now take a complete selected root
$\rho=1/2+iu_n(\theta)$ whose common order in $g$ and $h$
is $m$. The two Taylor expansions have leading coefficients
$g^{(m)}(\rho)/m!$ and $h^{(m)}(\rho)/m!$, so the canceled
entire quotient has the exact nonzero value
\[
 (g/h)(\rho)=\frac{g^{(m)}(\rho)}{h^{(m)}(\rho)},\qquad
 w_n(\theta)=\frac1L
 \left|\frac{g^{(m)}(\rho)}{h^{(m)}(\rho)}\right|^2>0.
 \tag{HCD18a}
\]
This retains both cancellation factorials and the full
Fourier mass $(2\pi/L)(1/(2\pi))$. It is the selected-root
calculation PSD5 and PSD41. Other $k=1$ coordinates where
$g/h$ vanishes still give the explicit null subspace above.
For a concrete measurable realization use the isometry
$v\mapsto(\sqrt{w_n(\theta)}v_n)_n$ into ordinary
$\ell^2(\mathbb Z)$, with zero coordinates at zero weights.
Its inverse on its range is
$v_n=y_n/\sqrt{w_n(\theta)}$ at positive weights; other
coordinates give the null class. This specifies the norm
map without changing a mass. The range projections are
measurable diagonal matrices; their images of the countable
standard basis give a measurable field.

The original gamma-contour and alternating-zeta calculation
PSD8--PSD13 proves, for every $0<\gamma<\pi/2$, the finite
constants
\[
 A_\gamma=\|e^{\gamma|\cdot|}w_h\|_\infty,\qquad
 B_\gamma=\|e^{\gamma|\cdot|}w_h\|_1.
\]
That calculation retains the leading coefficient of $h$,
its full degree, and the continuous canceled values on the
compact interval containing its selected roots. With
$f_\gamma(v)=e^{\gamma|v|}w_h(v)$, the ordinary triangle
inequality on the convolution variables and integration
of the remaining factors give, for every real $u$,
\[
 e^{\gamma|u|}m_{h,k}(u)
 \le f_\gamma^{*k}(u)
 \le A_\gamma B_\gamma^{k-1}.
 \tag{HCD19a}
\]
For the second inequality bound one factor by $A_\gamma$
and integrate each of the other $k-1$ factors with its
literal Lebesgue measure. For $k=1$ the expression is
$A_\gamma$, with $B_\gamma^0=1$.
For any $0<2\epsilon<\gamma$,
\[
 \sup_{0\le\theta<2\pi}
 \sum_nw_n(\theta)e^{2\epsilon|u_n(\theta)|}<\infty.
 \tag{HCD19}
\]
Indeed $|u_n(\theta)|\ge(2\pi/L)(|n|-1)$, so with
$d=\gamma-2\epsilon>0$ the sum is at most
\[
 \frac{2\pi}{L}A_\gamma B_\gamma^{k-1}
 e^{2\pi d/L}
 \left(1+\frac{2e^{-2\pi d/L}}{1-e^{-2\pi d/L}}\right).
\]
For every fixed polynomial $Q$, the constant
$\sup_{u\in\R}|Q(c+iu)|^2e^{-d|u|/2}$ is finite.
Inserting it into the same geometric estimate, with
$d/2$ in place of $d$, proves uniform exponential
summability of $w_n(\theta)|Q(s_n(\theta))|^2$.
This includes $Q=\chi$ and every fixed polynomial
relation used below. All original constants in the
weights are unchanged. The full source proof of
PSD8--PSD14 is a dependency of this continuation;
the remaining discrete density argument is given
explicitly here with its shifted Fourier phase.

\subsection{Density and the complete fibrewise Hilbert quotient}

Here is a direct density proof with the zero weights retained.
If $f\in\mathcal H_\theta$ is orthogonal to all polynomials,
put $a_n=w_n(\theta)\overline{f_n}$. Cauchy--Schwarz and
(HCD19) give $\sum|a_n|e^{\epsilon|u_n|}<\infty$.
Hence $F(z)=\sum_n a_ne^{zu_n}$ is holomorphic on
$|\Re z|<\epsilon$, with termwise derivatives on smaller
closed strips. Every derivative at zero is zero, since
polynomials in $S$ and in $(S-c)/i$ span the same space.
The Taylor series and the identity theorem give $F=0$
on the connected strip. On the imaginary axis,
\[
 F(it)=e^{it\theta/L}\sum_n a_ne^{2\pi int/L}.
\]
The series is absolutely uniformly convergent.
Multiplication by $e^{-it\theta/L}e^{-2\pi imt/L}$ and
integration over $0\le t\le L$ give $a_m=0$.
At positive weights this implies $f_m=0$; at zero weights
the class is already null. Thus polynomials are dense.
The identical argument proves polynomial density for
the weights $w_n|\chi(s_n)|^2$.

Multiplication by the original polynomial is the isometry
\[
 U_\chi:\ell^2(w_n|\chi(s_n)|^2)\longrightarrow\mathcal H_\theta,
 \qquad f\longmapsto(\chi(s_n)f_n)_n.
 \tag{HCD20}
\]
Its image is precisely the sequences vanishing at the
positive-weight sampled roots of $\chi$. To prove the
converse inclusion, take such a sequence $v$ and set
$f_n=v_n/\chi(s_n)$ when $\chi(s_n)\ne0$, and zero
elsewhere. Its squared domain norm equals $\|v\|_\theta^2$.
Thus the image is closed. Polynomial density in the domain
now proves
\[
 \mathcal R_\theta:=
 \overline{\chi\C[S]}^{\,\mathcal H_\theta}
 =\{v:v_n=0\text{ if }w_n(\theta)>0,\ \chi(s_n(\theta))=0\}.
 \tag{HCD21}
\]
The receiving Hilbert quotient is exactly
\[
 \begin{gathered}
 \mathcal Q_\theta=\mathcal H_\theta/\mathcal R_\theta
 \simeq
 \bigoplus_{\substack{n:\chi(s_n(\theta))=0\\w_n(\theta)>0}}\C,\\
 \|v\|_{\mathcal Q_\theta}^2
 =\sum_{\rm displayed\ n}w_n(\theta)|v_n|^2.
 \end{gathered}
 \tag{HCD22}
\]
The quotient map is orthogonal coordinate restriction.
No weight equal to zero is inverted.

\subsection{The complete primary kernel before phase integration}

Let $\Lambda_\theta^+$ be the distinct sampled roots of
$\chi$ with positive weight, and set
$p_\theta(S)=\prod_{\lambda\in\Lambda_\theta^+}(S-\lambda)$,
using $p_\theta=1$ for an empty product.
The finite map is
\[
 \rho_\theta:E\longrightarrow\mathcal Q_\theta,\qquad
 [P]_\chi\longmapsto(P(\lambda))_{\lambda\in\Lambda_\theta^+}.
 \tag{HCD23}
\]
It is well defined, and is onto by Lagrange interpolation
at these distinct points, using degree less than
$|\Lambda_\theta^+|\le q$. Vanishing at all those points
is equivalent to divisibility by $p_\theta$. Since
$p_\theta\mid\chi$,
\[
 \ker\rho_\theta=(p_\theta)/(\chi)\subset E.
 \tag{HCD24}
\]
This specifies the zero-weight-corrected kernel in H64:
a root at a zero-mass coordinate contributes no receiving
Hilbert line.

Write the full original factorization as
$\chi=\prod_{\lambda\in\Lambda}(S-\lambda)^{m_\lambda}$.
The complete Taylor map is
\[
 \begin{split}
 E&\longrightarrow\prod_{\lambda\in\Lambda}
 \C[\varepsilon_\lambda]/(\varepsilon_\lambda^{m_\lambda}),\\
 [P]&\longmapsto\left(
 \sum_{j=0}^{m_\lambda-1}\frac{P^{(j)}(\lambda)}{j!}
 \varepsilon_\lambda^j\right)_\lambda.
 \end{split}
 \tag{HCD25}
\]
For an explicit inverse, let
$q_\lambda=\chi/(S-\lambda)^{m_\lambda}$ and let
$u_\lambda$ be the Taylor polynomial of $1/q_\lambda$
at $\lambda$ through degree $m_\lambda-1$.
Then $e_\lambda=[u_\lambda q_\lambda]_\chi$ is one
in the $\lambda$ factor and zero in every other factor,
by divisibility by each original prime power. It follows
that $\sum e_\lambda=1$ and the factors are orthogonal
idempotents. The inverse is
$\sum_\lambda e_\lambda\sum_j a_{\lambda,j}(S-\lambda)^j$.

In these full primary coordinates, $\rho_\theta$ takes
the constant coefficient of each active sampled factor.
Its kernel in that factor is $(\varepsilon_\lambda)$,
retaining every order $1,\ldots,m_\lambda-1$.
For an unsampled factor or a zero-weight sampled factor,
the kernel is the entire primary algebra, including
its constant term. The original action is multiplication
by $\lambda+\varepsilon_\lambda$; the receiving action
is the simple scalar $\lambda$. Evaluation intertwines
these actions with the displayed nilpotent kernel.

\subsection{The measurable completed field and its vanishing map}

Every sampled root has real part $c$. For a fixed root $\lambda$
with $\Re\lambda=c$, its only possible phase is
\[
 \theta_\lambda=L\Im\lambda\pmod{2\pi},\quad
 0\le\theta_\lambda<2\pi,\qquad
 n_\lambda=(L\Im\lambda-\theta_\lambda)/(2\pi)\in\mathbb Z.
\]
Thus the set $\Theta_\chi$ of possible phases is finite.
For $\theta\notin\Theta_\chi$, $\mathcal Q_\theta=0$.
At an exceptional phase the positive-weight condition
in (HCD22) is still retained.

In the measurable realization after (HCD18), the projection
onto $\mathcal Q_\theta$ has diagonal entries
\[
 {\bf1}_{\{w_n(\theta)>0\}}
 {\bf1}_{\{\chi(s_n(\theta))=0\}},\qquad n\in\mathbb Z.
\]
These are measurable and give an explicit field of projections.
The direct integral satisfies
\[
 \mathscr Q=\int_\Theta^\oplus\mathcal Q_\theta\,d\mu=0,
 \tag{HCD26}
\]
because every measurable section has squared norm supported
on the finite measure-zero set $\Theta_\chi$.
This is an almost-everywhere assertion; each exceptional
nonzero receiving line is still exactly the weighted line
in (HCD22).

For $x\in E$, the section $\theta\mapsto\rho_\theta x$
is measurable: it is given by polynomial evaluations and
the displayed coordinate masks. It defines
\[
 \rho:E\longrightarrow\mathscr Q,\qquad
 x\longmapsto[\theta\mapsto\rho_\theta x]_{\rm a.e.},
 \qquad \ker\rho=E.
 \tag{HCD27}
\]
Every pointwise exceptional value has zero direct-integral
class. This does not assert that those pointwise maps vanish,
or that the original nonreduced coefficient algebra was zero.

\subsection{The source-to-completion cochain map}

The direct-integral source is the weighted space with squared
norm $\int_\Theta\sum_n w_n(\theta)|v_n(\theta)|^2\,d\mu$.
Let $\mathscr H$ denote this space, and let $\mathscr R$
be its closed subspace of sections belonging to
$\mathcal R_\theta$ almost everywhere. The coordinate
projection above has kernel $\mathscr R$.

The measurable evaluation map is
\[
 \mathsf T:\mathscr P\longrightarrow\mathscr H,\qquad
 P\longmapsto(P(\theta,s_n(\theta)))_{\theta,n}.
\]
The integrated identity (HCD18) proves it is an isometry.
Evaluation of a relation-valued section is in $\mathscr R$,
because the original section is a multiple of the literal
$\chi$. Hence there is a cochain map
\[
 [\mathscr B\hookrightarrow\mathscr P^{\rm eq}]
 \longrightarrow[\mathscr R\hookrightarrow\mathscr H],
 \tag{HCD28}
\]
using evaluation in both degrees. On a common image $x$,
the receiving quotient value is $\rho_\theta x$, since
representatives differ by an original relation.
The induced degree-one map is exactly (HCD27).

The global relation subspace is also the closure of measurable
finite polynomial relation sections. To prove this without
a measurable choice of approximating polynomials, suppose
$v\in\mathscr H$ is orthogonal to every section
$\mathbf1_U(\theta)\chi(S)S^j$, for measurable $U$ and $j\ge0$.
Its fibre inner product with $\chi S^j$ is integrable by
Cauchy--Schwarz, and its integral over every $U$ is zero,
so that inner product is zero almost everywhere.
The original exponential sample bounds make every fixed
polynomial section square-integrable.
Intersecting the countably many resulting full-measure sets,
(HCD21) gives $v(\theta)\in\mathcal Q_\theta$ there.
Conversely every such section is orthogonal to the relations.
The orthogonal-complement characterization proves that their
closed span is exactly $\mathscr R$.

Since $\mathscr Q=0$, $\mathscr R=\mathscr H$ as Hilbert
spaces, and the completed complex is the identity complex.
Its degree-one cohomology is zero. The finite continuous
common-image complex has positive Gram $\overline G$
on its full $E$ by (HCD8)--(HCD12).
The intervening map and its whole kernel are precisely
(HCD23)--(HCD28).

\subsection{The receiving supported zero and all marked fibres}

Integrate in the stated coefficient fibre, then apply the
original marked support construction. At a fixed original
label $\ell$, a linear map $f$ has the lift
\[
 (\ell,x)\longmapsto(\ell,f(x)),\qquad
 \tau\longmapsto\tau.
 \tag{HCD29}
\]
For (HCD27), every represented input maps to the receiving
supported zero $(\ell,0)$, and external absence maps to
external absence $\tau$. Their inverse images are respectively
the full represented fibre and the singleton external absence.
At a fixed phase, the represented kernel is the entire
ideal (HCD24), with its full primary decomposition (HCD25).
No integral on arbitrary mixed-support totals is presumed.

If the empty-packet case $q=0$ is separately admitted,
$E=0$, $B_N=P_N$, and the finite and continuous coefficient
complexes are identity complexes. The quotient Gram
determinant is the empty determinant one. The analytic
source and its literal norm remain present before the
quotient. Its marked receiving object still has $\tau$
and the represented zero, with the same maps (HCD29).

\subsection{Exact outcome and scope}

The continuous common-image complex has the explicit bounded
cochain contraction (HCD7) on its complement and retains
the whole nonreduced quotient in degree one.
Its minimum metric is $\overline G$, and its difference
from the constant-source metric is the positive original
relation norm (HCD13). The mean-zero decomposition retains
the weighted cross term (HCD15), giving (HCD16).

The subsequent completion has receiving weighted value
spaces (HCD22), primary kernels (HCD24)--(HCD25), measurable
field (HCD26), and actual cochain map (HCD28).
Its receiving coefficient direct integral is zero because
only finitely many phases can sample a fixed polynomial root.
Every exceptional phase, nilpotent component and zero-weight
coordinate remains in the stated maps and kernels.
No arithmetic volume asymptotic or RH conclusion is inferred.


\endgroup

% END COMPLETE CORRECTED SOURCE coordinating_next_wave/source_snapshot/accepted/holonomy_hcd_complete.tex


% BEGIN COMPLETE CORRECTED SOURCE coordinating_next_wave/phase_graph_averaged_bridge_prepared.tex
\section{The exact phase-graph map to the original averaged quotient and its completion}
\label{sec:pgb-averaged-bridge}

Fix the original full quartet and all the data of PGS.1--PGS.41.
Thus $m,k\geq1$, $d=4m$, $\rho=1/2+\delta+i\gamma$,
$0<\delta<1/2$, $\gamma>2$,
\[
 h(s)=\prod_{r\in\{\rho,\bar\rho,1-\rho,1-\bar\rho\}}(s-r)^m,
 \qquad v_h(s)=\frac{2\xi(s)}{h(s)},\qquad
 \Phi'=h,\quad\Phi(0)=0.
\]
Retain $s_i=1/2+it_i$, $\Sigma=\sum_i s_i=k/2+iu$,
$u=\sum_i t_i$, $\Phi(1/2+it)=A+iq_\Phi(t)$, $A>0$,
and $\tau_k=\sum_i\Phi(s_i)=kA+i\sum_iq_\Phi(t_i)$.
The polynomial $q_\Phi$ has degree $4m+1$ and leading
coefficient $1/(4m+1)$. These are the original objects, and the
full-divisor hypothesis of the original counterfactual packet is
unchanged. Let
\[
 \begin{gathered}
 \mathscr H_k=L^2\left(\mathbb R^k,\prod_i\frac{dt_i}{2\pi}\right),
 \qquad \mathsf T_k f=\tau_k f,
 \quad\operatorname{dom}\mathsf T_k
 =\{f\in\mathscr H_k:\tau_kf\in\mathscr H_k\},\\
 \chi=\chi_{h,k},\quad q=[1+k(m-1)](k+1)^2,
 \quad E=\mathbb C[S]/(\chi),\\
 \mathcal P_N=\mathbb C[S]_{\leq N},\quad N\geq q-1,
 \qquad B_N:\mathcal P_{N-q}\longrightarrow\mathcal P_N,
 \quad Q\longmapsto\chi Q,\\
 J_N:\mathcal P_N\longrightarrow E,\quad P\longmapsto[P],
 \qquad s:E\longrightarrow\mathcal P_{q-1},
 \quad [P]\longmapsto\operatorname{rem}_{\chi}P.
 \end{gathered}
 \tag{PGB.1}\label{eq:pgb-data}
\]
Here $\mathcal P_{-1}=0$. Every finite matrix below uses the
original increasing monomial frame, including the increasing
remainder frame of $E$. All Hilbert inner products are conjugate
linear in their first argument. Monic division proves
$\ker J_N=\operatorname{im}B_N$, injectivity of $B_N$, and
$J_Ns=I_E$.

Write the actual polynomial observation and defect as
\[
 \begin{gathered}
 (\mathcal V_NP)(t)=\Bigl(\prod_i v_h(s_i)\Bigr)P(\Sigma),
 \qquad Lx=\Bigl(\prod_i v_h(s_i)\Bigr)F_x(s_1,\ldots,s_k),\\
 F_x=\sum_{b=0}^{q-1}x_b\,
 \operatorname{rem}_{h(s_1),\ldots,h(s_k)}(\tau_k\Sigma^b),
 \qquad D_N=\mathsf T_k\mathcal V_N-LJ_N,\\
 \mathcal A_N=(\mathcal V_N,D_N):
 \mathcal P_N\longrightarrow\mathscr H_k\oplus\mathscr H_k,
 \qquad M_N=\mathcal V_N^*\mathcal V_N,
 \quad\widehat M_N=\mathcal A_N^*\mathcal A_N=M_N+D_N^*D_N.
 \end{gathered}
 \tag{PGB.2}\label{eq:pgb-original-maps}
\]
The definition of $F_x$ retains every original ordered tensor
remainder and every factor $v_h$. Its equality with
$\mathsf S_kC_k\eta x$ and the injectivity of $L$ are the
fully proved PGS.4 construction. In particular $F_x=0$ implies
$x=0$. The source decay of PGS and PGD places all the displayed
columns in the multiplication domain. The injectivity of
$\mathcal V_N$ makes both finite Grams positive definite.

\subsection{The positive cost of adding the literal defect}

Suppress $N$ only inside the next calculation. Define
\[
 \begin{gathered}
 H=B^*MB,\qquad
 C=s-BH^{-1}B^*Ms,\qquad G=C^*MC=(JM^{-1}J^*)^{-1},\\
 T=DB:\mathcal P_{N-q}\longrightarrow\mathscr H_k,
 \qquad H_+=H+T^*T=B^*\widehat M B,\\
 Q_D=I_{\mathscr H_k}-TH_+^{-1}T^*.
 \end{gathered}
 \tag{PGB.3}\label{eq:pgb-augmentation-data}
\]
If $N=q-1$, every map through the zero relation space is zero,
both relation determinants are one, and $Q_D=I$.
For positive relation dimension, $H,H_+$ are positive definite.
The formula for $C$ is proved by writing a lift as $sx+By$ and
completing its $M$ square. It gives $JC=I$ and $B^*MC=0$.
To verify the inverse formula for $G$, decompose every polynomial
as $CJp+By$; pairing against $Cx$ proves $MC=J^*G$.
Multiplication first by $M^{-1}$ and then by $J$ gives the formula.

The bounded operator $Q_D$ is positive and has the exact inverse
\[
 Q_D=(I_{\mathscr H_k}+TH^{-1}T^*)^{-1},\qquad
 0<Q_D\preceq I_{\mathscr H_k}.
 \tag{PGB.4}\label{eq:pgb-positive-inverse}
\]
Indeed $I+TH^{-1}T^*$ is self-adjoint and at least $I$.
The product of this operator with $I-TH_+^{-1}T^*$ is $I$:
the coefficient between $T$ and $T^*$ in the remaining term is
\[
 H^{-1}-H_+^{-1}-H^{-1}T^*TH_+^{-1}
 =H^{-1}-(I+H^{-1}T^*T)H_+^{-1}=0.
\]
Taking adjoints proves the product in the opposite order.
On the finite-dimensional range of $T$ positivity follows by
diagonalizing the positive operator $TH^{-1}T^*$; on its
orthogonal complement the inverse is $I$. This also proves
$Q_D\succeq(1+\|TH^{-1}T^*\|)^{-1}I$ and its bounded
positive square root without requiring an infinite-dimensional
spectral assertion.

Every lift is $Cx+By$. Its unchanged graph norm is
\[
 \begin{aligned}
 \|\mathcal A(Cx+By)\|^2
 &=x^*Gx+y^*Hy+\|DCx+Ty\|^2\\
 &=x^*\bigl(G+C^*D^*Q_DDC\bigr)x
 +(y+H_+^{-1}T^*DCx)^*H_+
                  (y+H_+^{-1}T^*DCx).
 \end{aligned}
 \tag{PGB.5}\label{eq:pgb-augmentation-square}
\]
Expansion proves this identity with its two positive cross terms
in the first square and the negative minimized contribution
$-C^*D^*TH_+^{-1}T^*DC$ retained. Therefore
\[
 \begin{gathered}
 \widehat R_N=C_N-B_N(H_{+,N})^{-1}T_N^*D_NC_N,\\
 \widehat G_N=G_N+\mathcal E_N^{\rm aug},\qquad
 \mathcal E_N^{\rm aug}=C_N^*D_N^*Q_{D,N}D_NC_N>0,\\
 \ker\mathcal E_N^{\rm aug}=\ker(D_NC_N)=\{0\}.
 \end{gathered}
 \tag{PGB.6}\label{eq:pgb-augmentation-quotient}
\]
The section has $J_N\widehat R_N=I$ and is orthogonal to the
relation columns in $\widehat M_N$. Thus uniqueness of a minimum
identifies it with the original section in PGS.13 and PGS.20;
the displayed $\widehat G_N$ is that same graph quotient.
The kernel assertion follows from the strict bounded positivity
of $Q_{D,N}$ proved above. The actual map whose Gram is the
new positive term is
$x\mapsto Q_{D,N}^{1/2}D_NC_Nx$ into $\mathscr H_k$.

Here the kernel can be completely evaluated on the actual source.
If $P\ne0$ has degree $n$ and leading coefficient $p_n$, the
polynomial $\tau_kP(\Sigma)$, as a polynomial in $s_1$, has
degree $n+d+1$ and leading coefficient $p_n/(d+1)$.
Indeed the original monic $h$ gives $\Phi$ leading coefficient
$1/(d+1)$, while the terms $\Phi(s_j)P(\Sigma)$ with $j>1$
have $s_1$ degree only $n$. On the other hand every variable
degree of $F_{JP}$ is at most $d-1$ by its original successive
remainders. Consequently $\tau_kP(\Sigma)-F_{JP}$ is a nonzero
polynomial. Its restriction under $s_i=1/2+it_i$ is nonzero
almost everywhere outside its polynomial zero set; the product
$\prod_i v_h(s_i)$ is nonzero almost everywhere because its
one-variable factors are nonzero entire functions. Thus
$D_NP\ne0$ in $\mathscr H_k$. This proves injectivity of
$D_N$. Since $J_NC_N=I$, it also proves injectivity of
$D_NC_N$, all asserted zero kernels, and the strict sign in
PGB.6. In particular $G_N<\widehat G_N$ even at the empty
relation window.

There is a second realization retaining both original graph
components. Put $\mathfrak a_N=\widehat R_N-C_N:E\to\operatorname{im}B_N$.
Orthogonality of $\mathcal V_NC_N$ and $\mathcal V_N\mathfrak a_N$ gives
\[
 \begin{gathered}
 \mathcal A_N\widehat R_N
 =(\mathcal V_NC_N,0)
       +(\mathcal V_N\mathfrak a_N,D_N\widehat R_N),\\
 \mathcal E_N^{\rm aug}=\mathfrak a_N^*M_N\mathfrak a_N
        +(D_N\widehat R_N)^*(D_N\widehat R_N).
 \end{gathered}
 \tag{PGB.7}\label{eq:pgb-augmentation-map}
\]
The two displayed summands in the first line are orthogonal.
Taking their Grams proves the second line and also proves its
equality with PGB.6, rather than merely comparing two presentations.

\subsection{Composition with the accepted all-holonomy quotient}

At each fixed $N$, retain the original HCD circumference, denoted
here by $L_{\rm circ}$ to distinguish that scalar from the
observation map $L:E\to\mathscr H_k$. Thus $L_{\rm circ}$
is exactly the quantity called $L$ in HCD, with the same value.
Retain the original holonomy matrices of HCD.2--HCD.2c.
Concretely choose the already
computed original $\beta>0$, weighted source constant
$\kappa_{\beta,N}$ and any $0<\eta<1$, and take
$L_{\rm circ}\geq\beta^{-1}\log(1+\kappa_{\beta,N}/\eta)$.
With $d\mu=d\theta/(2\pi)$ the HCD proof then gives
\[
 \begin{gathered}
 \int_0^{2\pi}M_{N,\theta}\,d\mu=M_N,
 \qquad(1-\eta)M_N\preceq M_{N,\theta}
                              \preceq(1+\eta)M_N,\\
 C_{N,\theta}=M_{N,\theta}^{-1}J_N^*G_{N,\theta},\qquad
 G_{N,\theta}=(J_NM_{N,\theta}^{-1}J_N^*)^{-1},\qquad
 \overline G_N=\int_0^{2\pi}G_{N,\theta}\,d\mu.
 \end{gathered}
 \tag{PGB.8}\label{eq:pgb-accepted-holonomy}
\]
No phase graph matrix is substituted into this accepted definition.
The full relative Hilbert fibre and the original half-density
source used to construct $M_{N,\theta}$ remain those of HCD.
Uniform finite-dimensional equivalence makes every following
coefficient field measurable and square integrable.

For each phase set
\[
 X_{N,\theta}=(B_N^*M_{N,\theta}B_N)^{-1}
                          B_N^*M_{N,\theta}C_N,
 \qquad C_N-C_{N,\theta}=B_NX_{N,\theta}.
 \tag{PGB.9}\label{eq:pgb-original-relation-field}
\]
At zero relation dimension $X_{N,\theta}$ is the zero map.
Both sections lift the same vector, so their difference lies in
$\operatorname{im}B_N$. Pairing the difference with $B_N$ in
the phase metric, and using $B_N^*M_{N,\theta}C_{N,\theta}=0$,
proves exactly the displayed formula. Squaring this orthogonal
decomposition of $C_N$ and integrating proves
\[
 \mathscr D_N^{\rm hol}:=G_N-\overline G_N
 =\int_0^{2\pi}(C_N-C_{N,\theta})^*M_{N,\theta}
                      (C_N-C_{N,\theta})\,d\mu\succeq0.
 \tag{PGB.10}\label{eq:pgb-original-gap}
\]
This is the exact accepted HCD.13 relation cost. In particular
the graph augmentation does not remove it.

Define
\[
 \mathscr B_N=L^2(\Theta;\operatorname{im}B_N,
                                  M_{N,\theta}\,d\mu),
 \quad
 \mathfrak F_Nx=
 \left(Q_{D,N}^{1/2}D_NC_Nx,
       [\theta\mapsto(C_N-C_{N,\theta})x]\right)
 \in\mathscr H_k\oplus\mathscr B_N.
 \tag{PGB.11}\label{eq:pgb-full-defect-map}
\]
Its complete matrix identity and kernel are
\[
 \begin{gathered}
 \boxed{\quad
 \widehat G_N-\overline G_N
 =\mathcal E_N^{\rm aug}+\mathscr D_N^{\rm hol}
 =\mathfrak F_N^*\mathfrak F_N>0,\quad}\\
 \ker\mathfrak F_N
 =\{x:D_NC_Nx=0,\ C_Nx=C_{N,\theta}x\text{ a.e.}\}=\{0\}.
 \end{gathered}
 \tag{PGB.12}\label{eq:pgb-main-bridge}
\]
Indeed the two orthogonal direct-sum target norms are precisely
PGB.6 and PGB.10. A sum of two nonnegative real norms vanishes
exactly when both vanish; the positivity in PGB.4 and that of
the phase source metric give the stated full kernel.
The actual injectivity proved after PGB.6 evaluates that kernel
as zero and makes the displayed matrix difference strictly positive.

The identity on coefficients consequently defines the actual
contraction $(E,\widehat G_N)\to(E,\overline G_N)$. It is
induced by an explicit map of the original two-term complexes:
send $p$ in $(\mathcal P_N,\widehat M_N)$ to the constant
field $\theta\mapsto p$ in the HCD common-image space, and
send a relation $B_Ny$ to that same constant relation field.
The inclusion differentials commute, $J_Np$ is unchanged, and
\[
 \int p^*M_{N,\theta}p\,d\mu=p^*M_Np
                  \leq p^*\widehat M_Np.
\]
Thus the source map is a contraction and its induced quotient
map is the displayed identity. More precisely, the complete
isometric dilation of this quotient contraction is
\[
 x\longmapsto
 \bigl(x,Q_{D,N}^{1/2}D_NC_Nx,
             [\theta\mapsto(C_N-C_{N,\theta})x]\bigr)
 \quad\hbox{from }(E,\widehat G_N)
 \hbox{ to }(E,\overline G_N)\oplus\mathscr H_k\oplus\mathscr B_N.
 \tag{PGB.13}\label{eq:pgb-isometric-dilation}
\]
Its squared norm identity is PGB.12; polarization proves every
cross pairing. This gives the precise relation between the graph
quotient, the accepted average, and the common finite-source
relation obstruction on the unchanged coefficient object $E$.

\subsection{Exponential moments and the complete graph quotient}

For clarity the density argument needed for the degree limit is
proved here. If a finite positive Borel measure $\nu$ on
$\mathbb R$ has $\int e^{\epsilon|u|}\,d\nu<\infty$ for some
$\epsilon>0$, then polynomials are dense in $L^2(\nu)$.
To prove this, let $f$ be orthogonal to them and put
$d\lambda=\bar f\,d\nu$. Cauchy--Schwarz gives
\[
 \int e^{\epsilon|u|/2}\,|d\lambda|
 \leq\|f\|_{L^2(\nu)}
       \left(\int e^{\epsilon|u|}\,d\nu\right)^{1/2}<\infty.
\]
The Fourier integral $F(z)=\int e^{izu}\,d\lambda(u)$ is
holomorphic on $|\operatorname{Im}z|<\epsilon/2$:
on a smaller closed strip every differentiated integrand is
dominated by a constant times $e^{\epsilon|u|/2}|d\lambda|$.
Every derivative at zero vanishes by polynomial orthogonality.
The Taylor expansion and then overlapping analytic discs give
$F=0$ throughout this connected strip, in particular on its
real axis.

Here is the needed uniqueness on that axis. Let
$g_t(x)=(4\pi t)^{-1/2}\exp(-x^2/(4t))$, $t>0$.
The elementary Gaussian integral gives
\[
 g_t(v)=\frac1{2\pi}\int_{\mathbb R}e^{-t\zeta^2}e^{i\zeta v}
                                                  \,d\zeta.
\]
For completeness, differentiating the latter integral in $v$
and integrating by parts in $\zeta$ gives the differential
equation $I'(v)=-vI(v)/(2t)$; its value at zero is
$I(0)=\sqrt{\pi/t}$, obtained by squaring the Gaussian integral
and using polar coordinates. These statements prove the formula
with its constants. Absolute integrability and Fubini now give
$\lambda*g_t=0$ because $F(-\zeta)=0$.
For $\varphi\in C_c(\mathbb R)$ it follows that
$\int(\varphi*g_t)\,d\lambda=0$. The functions
$\varphi*g_t$ converge uniformly to $\varphi$: split the
unit-mass Gaussian integral into $|y|<a$, where uniform
continuity controls the difference, and $|y|\geq a$, where
its mass tends to zero. Thus $\int\varphi\,d\lambda=0$ for
every such $\varphi$. Regularity of a finite Borel measure on
the real line, followed by continuous approximations to indicators
of compact sets and open sets, gives $\lambda=0$.
It follows that $f=0$ $\nu$-almost everywhere. The orthogonal
complement criterion for closed subspaces proves the density claim.

Apply this to the original two densities
\[
 \begin{gathered}
 w(t)=|v_h(1/2+it)|^2/(2\pi),\quad m_k=w^{*k},\\
 r_k(u)=\int_{\sum t_i=u}(1+|\tau_k|^2)\prod_iw(t_i),
 \qquad\mathscr L_k=L^2(\mathbb R,r_k(u)\,du).
 \end{gathered}
 \tag{PGB.14}\label{eq:pgb-original-density}
\]
All fibre integrals use elimination of $t_k$ and measure
$dt_1\cdots dt_{k-1}$, with determinant-one change of variables.
For $k=1$ the fibre integral is evaluation.
The proved PGD.11 bound is
$w(t)\leq C_he^{-\alpha|t|}(1+|t|)^{-p_0}$,
$\alpha=\pi/2$, $p_0=8m-9/2>1$.
For fixed $k,h$ and $0<\epsilon<\alpha$, both
\[
 \int e^{\epsilon|u|}r_k(u)\,du<\infty,
 \qquad
 \int e^{\epsilon|u|}|\chi(k/2+iu)|^2r_k(u)\,du<\infty
 \tag{PGB.15}\label{eq:pgb-actual-exponential-moments}
\]
follow directly on the product source. Indeed
$|u|\leq\sum_i|t_i|$, while $1+|\tau_k|^2$ and
$|\chi(k/2+iu)|^2$ are bounded by fixed products of
powers of $(1+|t_i|)$. Their product with
$\prod_i e^{-(\alpha-\epsilon)|t_i|}$ is integrable.
The same proof holds with $r_k$ replaced by $m_k$.
Every constant here depends on the fixed original packet and
$k$; no uniform statement in $k$ has been inserted.

The density is positive almost everywhere by PGS.35, and the
polynomial $\chi(k/2+iu)$ has only finitely many real zeros.
Consequently multiplication by that unchanged polynomial is
the onto isometry
\[
 L^2(\mathbb R,|\chi(k/2+iu)|^2r_k(u)\,du)
 \longrightarrow\mathscr L_k,
 \qquad f(u)\longmapsto\chi(k/2+iu)f(u).
 \tag{PGB.16}\label{eq:pgb-relation-density-map}
\]
Its inverse is division by $\chi(k/2+iu)$ away from its finite
zero set, with arbitrary values on that null set. Equality of
the two squared norms proves both the isometry and the inverse
domain assertion. PGB.15 and the density proof therefore imply
\[
 \overline{\{\chi(k/2+iu)Q(k/2+iu):Q\in\mathbb C[S]\}}
              ^{\,\mathscr L_k}=\mathscr L_k.
 \tag{PGB.17}\label{eq:pgb-dense-actual-relations}
\]
Polynomials in $u$ and in $S=k/2+iu$ have exactly the same span:
the forward substitution has inverse $u=(S-k/2)/i$.
This proves the density assertion on the original coordinate,
not on an independently substituted source.

Recall the actual isometry and observation from PGS.36--PGS.40:
\[
 \begin{gathered}
 \mathscr J_k:\mathscr L_k\longrightarrow
                      \mathscr H_k\oplus\mathscr H_k,
 \qquad
 \mathscr J_k f=
 \left(\prod_iv_h(s_i)f(u),\tau_k\prod_iv_h(s_i)f(u)\right),\\
 Y=(0,L),\qquad g=\mathscr J_k^*Y,
 \quad g_x(u)=z(u)x/r_k(u),\quad
 z_b(u)=\int_{\sum t_i=u}\overline{\tau_k}F_b(s)\prod_iw(t_i),\\
 Y_\perp=(I-\mathscr J_k\mathscr J_k^*)Y,
 \qquad\Omega_k=Y_\perp^*Y_\perp>0.
 \end{gathered}
 \tag{PGB.18}\label{eq:pgb-full-observation}
\]
These maps are the original isometry and mixed observation;
the adjoint is exactly $z/r_k$ with its conjugate phase.
To recall why the strict sign holds, fibre Cauchy--Schwarz gives
$|z(u)x|^2\leq(r_k(u)-m_k(u))x^*\Lambda(u)x$, where
$\Lambda_{bc}(u)=\int_{\sum t_i=u}\overline{F_b}F_c\prod_iw$.
It follows that
$\Omega_k\succeq\int(m_k/r_k)\Lambda\,du>0$;
for nonzero $x$, the integral of $x^*\Lambda x$ is
$\|Lx\|^2>0$, and $m_k/r_k>0$ almost everywhere.
This argument also proves $g_x\in\mathscr L_k$ and the stated
adjoint by direct pairing. Thus the preceding uses of these
objects involve actual bounded maps on the specified domains.

Put $\iota P(u)=P(k/2+iu)$. For every original polynomial $P$
write $\mathcal A P$ and $JP$ for the compatible union of the
finite-window maps, whose definitions agree on every inclusion.
Then
the exact source identity is
\[
 \mathcal A P=
 \mathscr J_k(\iota P-gJ P)-Y_\perp J P,
 \qquad
 \|\mathcal A P\|^2=
 \|\iota P-gJ P\|_{\mathscr L_k}^2+(JP)^*\Omega_k(JP).
 \tag{PGB.19}\label{eq:pgb-completion-source-identity}
\]
It follows by substituting
$Y=\mathscr J_kg+Y_\perp$ into $\mathcal A=U-YJ$.
The two summands lie in orthogonal closed subspaces. Define
\[
 \begin{gathered}
 \mathscr K_\infty=
 \mathscr J_k\mathscr L_k\ \mathbin{\oplus^\perp}\ Y_\perp E,
 \qquad
 \Xi:\mathscr L_k\oplus(E,\Omega_k)\longrightarrow\mathscr K_\infty,
 \quad (f,x)\longmapsto\mathscr J_k f-Y_\perp x,\\
 \Xi^{-1}(v)=
 \left(\mathscr J_k^*v,-\Omega_k^{-1}Y_\perp^*v\right).
 \end{gathered}
 \tag{PGB.20}\label{eq:pgb-completion-isometry}
\]
The space is closed because the first summand is the range of
an isometry and the second is finite dimensional and orthogonal.
The formulas $\mathscr J_k^*Y_\perp=0$ and
$Y_\perp^*Y_\perp=\Omega_k$ prove both inverse identities,
and prove the exact squared norm
$\|\Xi(f,x)\|^2=\|f\|^2+x^*\Omega_kx$.
This is a complete isometric parametrization with its minus
sign determined by the original subtraction in $D$.

In fact $\mathscr K_\infty$ is exactly the closure of the
original polynomial graph image. Inclusion follows from PGB.19.
Conversely take $(f,x)$ in the domain of $\Xi$.
The affine set
\[
 \{\iota(sx+\chi Q)-g_x:Q\in\mathbb C[S]\}
\]
is dense in $\mathscr L_k$ by PGB.17. Choose $Q_j$ for which
these functions converge to $f$ and put $P_j=sx+\chi Q_j$.
Then $JP_j=x$ for every $j$, and PGB.19 proves
$\mathcal AP_j\to\Xi(f,x)$. This proves surjectivity onto
the completed graph from limits of original polynomials,
retaining the prescribed quotient coordinate throughout.
In the same manner
\[
 \overline{\mathcal A(\chi\mathbb C[S])}
      =\mathscr J_k\mathscr L_k,
 \qquad
 \mathscr K_\infty/\overline{\mathcal A(\chi\mathbb C[S])}
       \simeq(E,\Omega_k).
 \tag{PGB.21}\label{eq:pgb-completed-quotient}
\]
The quotient map on the whole completed source is the bounded
map $v\mapsto-\Omega_k^{-1}Y_\perp^*v$; it sends
$\mathcal AP$ to $JP$ and has precisely the displayed closed
relation space as kernel. Its section is $x\mapsto-Y_\perp x$.
This explicitly computes the cochain completion and every map.

This completed source is a closed linear relation in
$\mathscr H_k\oplus\mathscr H_k$, with its vertical part fully
determined. Indeed $Y=\mathscr J_kg+Y_\perp$ gives
\[
 \Xi(f,x)=\mathscr J_k(f+g_x)-Yx,
 \qquad
 \mathscr K_\infty\cap(\{0\}\oplus\mathscr H_k)
          =\{(0,-Lx):x\in E\}.
 \tag{PGB.21a}\label{eq:pgb-vertical-kernel}
\]
If the first component vanishes, its squared norm is
$\int|f+g_x|^2m_k\,du=0$. Since $m_k>0$ almost everywhere,
$f=-g_x$ almost everywhere, also as elements of $\mathscr L_k$;
the source vector is therefore $-Yx$. Conversely
$\Xi(-g_x,x)=-Yx$ proves the reverse inclusion.
In particular the original densely defined defect operator
$\mathcal VP\mapsto DP$ on the closure of the polynomial
observation source is not closable. For any fixed $x\ne0$,
PGB.17 supplies $P_j=sx+\chi Q_j$ with
$\iota P_j\to0$ in $\mathscr L_k$. The original first
components tend to zero, while
$DP_j=\mathsf T_k\mathcal VP_j-Lx\to-Lx\ne0$.
The two convergence assertions follow from the two component
norms of $\mathscr J_k\iota P_j$. This witnesses the entire
vertical subspace without asserting that the completed relation
is the graph of a single-valued operator.

There is also an explicit map to the accepted unweighted
completion. Put $\mathscr L_k^0=L^2(m_kdu)$ and
\[
 \mathscr J_k^0:\mathscr L_k^0\to\mathscr H_k,
 \quad f\mapsto\prod_i v_h(s_i)f(u),\qquad
 I_{r,m}:\mathscr L_k\to\mathscr L_k^0,\quad f\mapsto f.
\]
The first is an isometry by the original convolution integral;
the second is a contraction because $r_k\geq m_k$.
The first-coordinate projection $\pi_1$ obeys
$\pi_1\Xi(f,x)=\mathscr J_k^0I_{r,m}(f+g_x)$ and
$\pi_1\mathscr J_k f=\mathscr J_k^0I_{r,m}f$.
Since the same exponential-moment argument proves that the
unweighted relation closure is all of $\mathscr J_k^0\mathscr L_k^0$,
these formulas define a bounded cochain map
\[
 [\mathscr J_k\mathscr L_k\hookrightarrow\mathscr K_\infty]
 \xrightarrow{\ \pi_1\ }
 [\mathscr J_k^0\mathscr L_k^0\xrightarrow{I}
                         \mathscr J_k^0\mathscr L_k^0].
 \tag{PGB.21b}\label{eq:pgb-to-unweighted-completion}
\]
It agrees on the original polynomial complex with
$\mathcal AP\mapsto\mathcal VP$ and
$\mathcal A(\chi Q)\mapsto\mathcal V(\chi Q)$.
Its degree-one map is $(E,\Omega_k)\to0$ with kernel all
of $E$. Thus the positive completed graph quotient maps exactly
to the zero unweighted completion in the accepted R66 source;
its nonzero metric and the receiving zero object are connected
by this specific cochain map and its full kernel.

\subsection{The degree limit and the exact obstruction to a uniform comparison}

For any $x\in E$, minimizing PGB.19 over $P\in\mathcal P_N$
with $JP=x$ gives
\[
 x^*\widehat G_Nx
 =x^*\Omega_kx+
 \inf_{Q\in\mathcal P_{N-q}}
       \|\iota(sx+\chi Q)-g_x\|_{\mathscr L_k}^2.
 \tag{PGB.22}\label{eq:pgb-quotient-distance}
\]
At $N=q-1$ the infimum contains exactly $Q=0$.
The nested relation spaces and their proved density show that
the nonnegative second term decreases to zero for every $x$.
Using polarization on the fixed finite remainder frame gives
entrywise convergence, hence operator-norm convergence of the
finite matrices. Thus, at fixed $h,k$,
\[
 \widehat G_N\downarrow\Omega_k>0,
 \qquad G_{W,N}\downarrow0,
 \qquad\Delta_N\downarrow\Omega_k,
 \qquad G_N\downarrow0.
 \tag{PGB.23}\label{eq:pgb-three-limits}
\]
For $G_{W,N}$ repeat the distance calculation with $g=0$ and
without the $\Omega_k$ term, using the same $r_k$ and PGB.17.
For $\Delta_N$, its PGS.40 expression is
$\Omega_k+((I-\mathsf P_N)g)^*((I-\mathsf P_N)g)$;
polynomial density in $\mathscr L_k$ makes the latter term
decrease to zero. For $G_N$ repeat the affine approximation
argument in $L^2(m_kdu)$, whose exponential moments and dense
relation image were proved by the same PGB.15--PGB.17 argument.
These are matrix limits in the unchanged frame, not estimates
in the separate tensor variable $k$.

The original averaged minimum form is defined at every
$L_{\rm circ}>0$, including a phase with a singular finite
matrix, by its actual source map $T_\theta$ into the sampled
Hilbert space of HCD.18. Let $P_\theta^{\rm rel}$ be orthogonal
projection onto the finite-dimensional closed subspace
$T_\theta(\operatorname{im}B_N)$. Then
\[
 \begin{aligned}
 x^*G_{N,\theta}x
 &=\|(I-P_\theta^{\rm rel})T_\theta C_Nx\|^2,\\
 x^*(G_N-\overline G_N)x
 &=\int\|P_\theta^{\rm rel}T_\theta C_Nx\|^2\,d\mu.
 \end{aligned}
\]
The first equality follows by minimizing over the actual affine
source image; the second follows from its orthogonal
decomposition and $\int\|T_\theta C_Nx\|^2d\mu=x^*G_Nx$.
The projections are measurable: for the map $T_\theta B_N$,
their finite-rank formula is the pointwise limit as $t\downarrow0$
of $(T_\theta B_N)(B_N^*M_{N,\theta}B_N+tI)^{-1}
(T_\theta B_N)^*$. Diagonalization of the finite Gram proves
this limit, including its zero eigenvalues. The displayed
integrands are bounded above by the integrable full source
norm, so both minimum forms and integrals are defined.
Consequently $0\preceq\overline G_N\preceq G_N$ at every
original circumference, with no inverse imposed on a singular
phase. Positivity of the finite $\overline G_N$ follows directly:
for $k\geq2$ every original sampled convolution weight is
positive; for $k=1$ all weights are positive outside a countable
set of phases determined by the isolated real zeros of $v_h$.
At such a phase a nonzero polynomial cannot vanish on infinitely
many distinct sample coordinates, so its finite source Gram
and its quotient Gram are positive definite. Integrating any
nonzero quotient vector gives a strictly positive norm.

The identity $G_N\to0$ is therefore precisely consistent with
the accepted R66/CR conclusion. For these actual averaged
minimum forms the augmentation identity gives
\[
 \mathcal E_N^{\rm aug}=\widehat G_N-G_N\longrightarrow\Omega_k,
 \qquad
 \widehat G_N-\overline G_N\longrightarrow\Omega_k
 \tag{PGB.24}\label{eq:pgb-comparison-limit}
\]
at every fixed $L_{\rm circ}>0$.
This conclusion also holds for any choices $L_{\rm circ}=L_N$ satisfying
PGB.8 at each finite cutoff; no interchange of the limit with
a changing phase source was used, since the order bound alone
gives $\|\overline G_N\|\leq\|G_N\|\to0$.

Where $\overline G_N>0$, which the accepted source proves at
every finite $N$, the comparison has the explicit divergence
\[
 \inf_{x\ne0}
   \frac{x^*\widehat G_Nx}{x^*\overline G_Nx}
 \ \geq\ \frac{\lambda_{\min}(\Omega_k)}{\|G_N\|}
 \ \longrightarrow\ +\infty.
 \tag{PGB.25}\label{eq:pgb-uniform-obstruction}
\]
Indeed the numerator is at least
$\lambda_{\min}(\Omega_k)\|x\|^2$, while the denominator
is at most $\|G_N\|\|x\|^2$. Both constants refer to the same
original remainder frame. Hence no finite constant independent
of $N$ can bound $\widehat G_N$ above by that constant times
the accepted $\overline G_N$. The direction
$\overline G_N\preceq\widehat G_N$ remains the exact
cochain-induced contraction PGB.12--PGB.13. This proves an
actual relation and its precise obstruction, rather than
inferring lack of relation from the obstruction.

Finally fix the same original circumference at all four endpoints.
For the inverse-form realization PGB.8, choose it using the
original common maximal source $N_*=2q+2$ as in R63 and restrict
its bound to every endpoint source. Alternatively use the
proved minimum-form realization at any fixed circumference.
The literal endpoint convention remains
$(q-1,q,2q-1,2q)$ with signs $(+,+,-,-)$. Set
$K_N^{\rm av}=\overline G_N^{-1/2}\widehat G_N
\overline G_N^{-1/2}>I$ at these finite endpoints.
Direct determinant multiplication gives the exact signed bridge
\[
 \begin{aligned}
 &\bigl[\log\det\widehat G_{q-1}+\log\det\widehat G_q
        -\log\det\widehat G_{2q-1}-\log\det\widehat G_{2q}\bigr]\\
 &\quad-
 \bigl[\log\det\overline G_{q-1}+\log\det\overline G_q
        -\log\det\overline G_{2q-1}-\log\det\overline G_{2q}\bigr]\\
 &=\log\det K_{q-1}^{\rm av}+\log\det K_q^{\rm av}
       -\log\det K_{2q-1}^{\rm av}-\log\det K_{2q}^{\rm av}.
 \end{aligned}
 \tag{PGB.26}\label{eq:pgb-four-endpoint-bridge}
\]
Each separate logarithm on the right is nonnegative; their
signed sum has no sign consequence from that fact alone.
The identity retains all four original endpoints, including
their possible repeated appearance when $q=1$.

All the maps on $E$ and polynomial coefficient spaces above
extend to a fixed original active mask by applying the same
linear map in each coordinate. For PGB.12 the represented-zero
kernel is exactly the stated kernel in each coordinate; for
the completed quotient PGB.21 it is the stated closed relation
space. The empty present tuple stays present. External absence
maps to external absence. At any of the already specified
compatible realizations $R\to\mathbb C$, these complex-linear
maps are $R$-linear with the given action, and their scalar
extension sends $b\otimes x$ to $b\otimes f(x)$. No new
residue-field realization, elimination of arithmetic nilpotents,
or identification of supported zero with absence is made.

% END COMPLETE CORRECTED SOURCE coordinating_next_wave/phase_graph_averaged_bridge_prepared.tex


% BEGIN COMPLETE CORRECTED SOURCE coordinating_next_wave/mixed_residual_completion_addendum.tex
\section{The sealed strict mixed residual inside the completed graph quotient}
\label{sec:pgmb}

Retain the entire original packet, measures, coordinates and maps
of PGS.1--41, PGM.1--12 and PGB.1--26, including PGB.21a--21b.
The PGM source used here has SHA-256
\texttt{105a4a7fc7a4f1853aea5884f58fc45a6dae74aab52a86784401e3662f310b40}.
Write $\mathscr D_k^{\rm mix}$ for exactly the map called
$\mathscr D_k$ in that source. Its original domain is
$E=\mathbb C[S]/(\chi_{h,k})$, and its original target is
\[
 \mathscr K_k^{\rm mix}
 =L^2\left(\mathbb R^k,
  \left[1+k^2A^2+\left(\sum_iq_\Phi(t_i)\right)^2\right]
                         \prod_iw(t_i)\,dt_i\right),
 \qquad w(t)=|v_h(1/2+it)|^2/(2\pi).
 \tag{PGMB.1}\label{eq:pgmb-domain}
\]
This target is the PGM weighted scalar space, distinct from the
closed two-component relation $\mathscr K_\infty$ of PGB.20.
With $u=\sum_it_i$, $\tau_k=kA+i\sum_iq_\Phi(t_i)$,
$F_x$ the full original successive remainder observation, and
$g_x=z(u)x/r_k(u)$ the actual PGS observation adjoint, the map is
\[
 (\mathscr D_k^{\rm mix}x)(t)
 =\frac{\overline{\tau_k}F_x(s)}{1+|\tau_k|^2}-g_x(u),
 \qquad
 (\mathscr D_k^{\rm mix})^*\mathscr D_k^{\rm mix}
                       =\Omega_k-K_{\rm low}.
 \tag{PGMB.2}\label{eq:pgmb-mixed-map}
\]
All adjoints use the displayed original weighted target and the
original increasing remainder frame of $E$. The second equality
is the sealed PGM.5 exact Gram, with its complete proof from the
orthogonal sum-coordinate projection. PGM.6--9 proves
\[
 \ker\mathscr D_k^{\rm mix}=\{0\},\quad
 \Omega_k-K_{\rm low}>0\quad(k\geq2),
 \qquad
 \mathscr D_1^{\rm mix}=0,\quad\Omega_1=K_{\rm low}.
\]
In particular this is not the finite polynomial defect
$D_N=\mathsf T_k\mathcal V_N-LJ_N$ of PGB.2: its domain and
target differ, and the latter is injective even when $k=1$ by
the original leading-coefficient proof after PGB.6.

The completed quotient of PGB.21 has the exact isometric map
\[
 \mathfrak I_k:(E,\Omega_k)\longrightarrow
       (E,K_{\rm low})\oplus\mathscr K_k^{\rm mix},
 \qquad x\longmapsto(x,\mathscr D_k^{\rm mix}x).
 \tag{PGMB.3}\label{eq:pgmb-completed-isometry}
\]
Indeed the squared target norm is
$x^*K_{\rm low}x+\|\mathscr D_k^{\rm mix}x\|^2
=x^*\Omega_kx$ by PGMB.2. Polarization proves equality of
all inner products; the first coordinate proves injectivity.
Projection to that coordinate is the actual identity contraction
$(E,\Omega_k)\to(E,K_{\rm low})$. Its loss is exactly the
mixed map Gram, with zero loss kernel for $k\geq2$ and full
loss kernel $E$ for $k=1$. The quotient map itself has kernel
zero in both cases. Its composition with the completed source
quotient is explicitly
$v\mapsto-\Omega_k^{-1}Y_\perp^*v$, now regarded as a vector
of $(E,K_{\rm low})$; its source kernel remains precisely
$\mathscr J_kL^2(r_kdu)$. Thus the strict PGM contribution is
embedded in the existing completed morphism without altering
its sign, quotient coordinate or relation kernel.

There is a compatible finite-window identity. Let
$\widehat R_N:E\to\mathcal P_N$ be the same actual minimum
section of PGS.13 and PGB.6, and define the bounded map
\[
 f_N:E\longrightarrow L^2(r_kdu),\qquad
 f_Nx=\iota\widehat R_Nx-g_x,\qquad
 (\iota P)(u)=P(k/2+iu).
 \tag{PGMB.4}\label{eq:pgmb-finite-tail}
\]
Since $J_N\widehat R_N=I_E$, PGB.19 evaluated on that section
gives $\widehat G_N=\Omega_k+f_N^*f_N$. Inserting PGMB.2
gives the full original Gram and its isometric realization
\[
 \begin{gathered}
 \widehat G_N=K_{\rm low}
       +(\mathscr D_k^{\rm mix})^*\mathscr D_k^{\rm mix}+f_N^*f_N,\\
 (E,\widehat G_N)\longrightarrow
 (E,K_{\rm low})\oplus\mathscr K_k^{\rm mix}\oplus L^2(r_kdu),
 \qquad x\longmapsto(x,\mathscr D_k^{\rm mix}x,f_Nx).
 \end{gathered}
 \tag{PGMB.5}\label{eq:pgmb-three-term-isometry}
\]
Each term is the norm of its displayed original map, so expanding
the direct-sum norm proves this identity, and polarization
proves every cross pairing. PGB.22--23 proves
$f_N^*f_N=\widehat G_N-\Omega_k\to0$ at fixed original $h,k$.
It follows that $\|f_N\|\to0$, since its squared operator norm
in the fixed coefficient frame equals $\|f_N^*f_N\|$.
The limiting isometry is exactly PGMB.3 with zero third
component. Therefore the strict mixed contribution survives
the polynomial degree limit for every $k\geq2$, whereas the
finite affine-approximation contribution tends to zero.
The original relation $0<K_{\rm low}<\Omega_k\leq\widehat G_N$
is now strict in its first comparison for every $k\geq2$.
No uniform eigenvalue lower bound in the distinct parameter
$k$ is inferred from this fixed-packet argument.

The first summand also has its exact original Hilbert-source
realization. Define
\[
 \begin{gathered}
 L_{\rm low}:E\longrightarrow\mathscr H_k,\qquad
 L_{\rm low}x=
 \frac{\bigl(\prod_i v_h(s_i)\bigr)F_x(s)}
                   {\sqrt{1+|\tau_k|^2}},\qquad
 L_{\rm low}^*L_{\rm low}=K_{\rm low},\\
 \mathfrak U_k:E\longrightarrow
                   \mathscr H_k\oplus\mathscr K_k^{\rm mix},
 \qquad\mathfrak U_kx=(L_{\rm low}x,\mathscr D_k^{\rm mix}x),
 \qquad\mathfrak U_k^*\mathfrak U_k=\Omega_k.
 \end{gathered}
 \tag{PGMB.6}\label{eq:pgmb-original-source-isometry}
\]
The square root is the positive real square root of the
unchanged scalar $1+k^2A^2+(\sum_iq_\Phi(t_i))^2$.
The map lies in $\mathscr H_k$ since this reciprocal multiplier
has absolute value at most one and $Lx\in\mathscr H_k$.
Taking its original product integral gives exactly the entries
of PGS.17, proving the first Gram identity. PGMB.2 proves the
second. Thus $\mathfrak U_k$ is an isometry from
$(E,\Omega_k)$ onto its actual closed finite-dimensional image.
Its bounded left inverse on the whole displayed target is
\[
 (a,b)\longmapsto
 \Omega_k^{-1}\bigl(L_{\rm low}^*a+
                         (\mathscr D_k^{\rm mix})^*b\bigr).
\]
Composition with $\mathfrak U_k$ gives $I_E$ by the computed
Gram. The reverse composite is self-adjoint, idempotent, and
has exactly $\mathfrak U_k(E)$ as its range, so it is the
orthogonal projection in the original direct-sum Hilbert metric.

Let $j_\infty:\mathscr K_\infty\to E$ be the exact completed
quotient map $j_\infty v=-\Omega_k^{-1}Y_\perp^*v$ of PGB.21.
The complete original-source realization is
\[
 \begin{gathered}
 \mathfrak C_k=\mathfrak U_kj_\infty:
 \mathscr K_\infty\longrightarrow
                  \mathscr H_k\oplus\mathscr K_k^{\rm mix},\\
 \mathfrak C_k\Xi(f,x)
      =(L_{\rm low}x,\mathscr D_k^{\rm mix}x),\qquad
 \ker\mathfrak C_k=\mathscr J_kL^2(r_kdu),\qquad
 \|\mathfrak C_k\Xi(f,x)\|^2=x^*\Omega_kx.
 \end{gathered}
 \tag{PGMB.7}\label{eq:pgmb-completed-source-composite}
\]
The identities follow from $j_\infty\Xi(f,x)=x$ and
PGMB.6. Injectivity of $\mathfrak U_k$ and the computed
kernel of $j_\infty$ prove the complete kernel assertion.
This is a map on the actual completed polynomial graph source;
it factors through its specifically proved cyclic quotient.
It asserts no factorization through an unidentified larger
or proper-source quotient. Its norm equals the distance to
the completed relation subspace by the orthogonal coordinates
in PGB.20. In particular it is a contraction on that whole
completed source and an isometry on its quotient.

For $k\geq2$ the second component alone detects the entire
quotient. Its inverse on the mixed image is exactly the
already computed PGM.12 map
\[
 \mathscr R_k^{\rm mix}
 = (\Omega_k-K_{\rm low})^{-1}(\mathscr D_k^{\rm mix})^*,
 \qquad
 \mathscr R_k^{\rm mix}\mathscr D_k^{\rm mix}=I_E,
 \qquad
 \mathscr R_k^{\rm mix}\,\operatorname{pr}_2\mathfrak C_k
                                  =j_\infty.
 \tag{PGMB.8}\label{eq:pgmb-mixed-inverse-composite}
\]
The first composition follows by substituting PGMB.2; the
second follows by composing that identity with $j_\infty$.
For $k=1$ the second component is identically zero by PGM.9;
the first component still has positive Gram $K_{\rm low}$
and therefore retains the complete quotient.

The exact relative Gram and determinant cost can be written
without changing the original coefficient frame. Set
\[
 \begin{gathered}
 H_k^{\rm mix}=\Omega_k-K_{\rm low},\qquad
 S_k^{\rm mix}=K_{\rm low}^{-1/2}H_k^{\rm mix}K_{\rm low}^{-1/2},\qquad
 q=q_k=[1+k(m-1)](k+1)^2,\\
 0\leq\lambda_{k,1}\leq\cdots\leq\lambda_{k,q_k}
       \quad\hbox{the eigenvalues of }S_k^{\rm mix}.
 \end{gathered}
\]
These are generalized eigenvalues of the actual two original
forms $H_k^{\rm mix},K_{\rm low}$; the positive congruence
specifies their computation and does not replace either form.
Every $\lambda_{k,j}>0$ for $k\geq2$ and every
$\lambda_{1,j}=0$. Finite-dimensional Hermitian diagonalization
and determinant multiplication give
\[
 \begin{gathered}
 K_{\rm low}^{-1/2}\Omega_kK_{\rm low}^{-1/2}
                                     =I+S_k^{\rm mix},\\
 (1+\lambda_{k,1})K_{\rm low}\preceq\Omega_k
             \preceq(1+\lambda_{k,q_k})K_{\rm low},\\
 c_k^{\rm mix}:=\log\frac{\det\Omega_k}{\det K_{\rm low}}
       =\sum_{j=1}^{q_k}\log(1+\lambda_{k,j}),\qquad
 \begin{cases}c_1^{\rm mix}=0,\\c_k^{\rm mix}>0& (k\geq2).
 \end{cases}
 \end{gathered}
 \tag{PGMB.9}\label{eq:pgmb-relative-cost}
\]
The extremal constants are exact: test the original vector
$x=K_{\rm low}^{-1/2}v$ for a corresponding eigenvector $v$.
All $k$ dependence is displayed through the unchanged degree
$q_k$, the full matrices (PGS.17), (PGS.40), and (PGMB.2),
and their stated eigenvalues. No bound for these eigenvalues
uniform in $k$ follows merely from their strict positivity.

At finite $N$ put
$S_{k,N}^{\rm tail}=\Omega_k^{-1/2}f_N^*f_N\Omega_k^{-1/2}$
and let its nonnegative eigenvalues be
$\nu_{k,N,1},\ldots,\nu_{k,N,q_k}$. Then
\[
 \log\frac{\det\widehat G_N}{\det K_{\rm low}}
  =\sum_{j=1}^{q_k}\log(1+\lambda_{k,j})
       +\sum_{j=1}^{q_k}\log(1+\nu_{k,N,j}),\qquad
 \nu_{k,N,j}\longrightarrow0\quad(N\to\infty; h,k\text{ fixed}).
 \tag{PGMB.10}\label{eq:pgmb-finite-determinant-cost}
\]
Indeed PGMB.5 gives
$\widehat G_N=\Omega_k+f_N^*f_N$, so multiplying determinants
after the positive congruence yields
$\det\widehat G_N/\det\Omega_k=\det(I+S_{k,N}^{\rm tail})$.
The stated limit follows from the proved norm convergence of
$f_N$ and the fixed positive matrix $\Omega_k$.
Taking the four original signed endpoints cancels the same
$\log\det\Omega_k$ or $c_k^{\rm mix}$ with coefficients
$1+1-1-1=0$; their residual expression is exactly the signed
sum of these four tail determinant costs. This cancellation
does not remove $\Omega_k$ from the tail congruences and gives
no estimate at the distinct growing endpoints $N\asymp q_k$
as $k$ varies.

On each retained coordinate mask all these maps act coordinatewise.
The represented-zero kernel of $\mathfrak C_k$ is the product
of its displayed completed relation kernels, and the represented
kernel of its second component is that same product for
$k\geq2$ and the whole completed fibre for $k=1$.
The present empty tuple remains present, and external absence
maps to external absence with its singleton inverse image.
Compatible coefficient realizations are exactly those specified
in PGM.10, with the full kernel of their realized tensor map;
composition with the injective quotient realizations above
does not remove that earlier supported-zero fibre.

The finite graph-to-average morphism PGB.12--13 and the
completed first-coordinate cochain map PGB.21b retain their
original domains and kernels. In particular the former still
has its strictly positive two-term loss relative to
$\overline G_N$, and the latter still induces
$(E,\Omega_k)\to0$ with full coefficient kernel. The added
mixed isometry computes the interior structure of that same
$\Omega_k$ and supplies no identification of either target
with the other. The signed four endpoints in PGB.26 and PGS.34
retain their complete differences; none receives a new sign
from PGMB.2 or PGMB.5.

% END COMPLETE CORRECTED SOURCE coordinating_next_wave/mixed_residual_completion_addendum.tex


\end{document}
